% SPDX-FileCopyrightText: 2024-2026 Fabian Januszewski
% SPDX-License-Identifier: CC-BY-4.0

%======================================================================
% cli_manual.tex --- [CLI], "hacs: the command line and the
% OC console. A user manual"
%
%   Created in hecke-extensions slot 18, package WP-M (owner ruling
%   2026-10-03: a full LaTeX/PDF user manual of the command-line interface
%   and the console). 1.5.15: the campaign key hz_theta_order (S27 TS). 1.5.14: LEVERS-26 retired env arms (wording). 1.5.13: the program is hacs throughout; the census-runbook sections 8.7-8.8 replaced by a neutral section (release audit). 1.5.11: machine names and private paths removed (public release). 1.5.10: the env row ENGINE_VNNI2_FLOOR (S25 PERF-VF). 1.5.9: slot-23 env arms retired (S25 LEVERS-25). 1.5.8: Ctrl-C / interrupt at cancellation points (S25 NBCONFIRM S5). 1.5.7: LaTeX polish stage 1 (uniform author block, dedication; no table or count changed). 1.5.6: the verb hecke_structure, campaign keys structure / algebra / hz_cert (S25 HZ-REC). 1.5.5: retired GPU levers (S24 close, CL-R2 audit M6). 1.5.4: descent charpoly `exact' (S24 FMTMP). 1.5.3: no argument prints the usage, exit 1 (S24 CL-7). 1.5.2: the verb display (NB-R34). 1.5.1: referee fixes (DOCS-FIX). Restructured by subject at 1.5.0 (slot 24, DOCS-CLI;
%   owner direction 2026-10-29: a logically consistent manual that makes the
%   mathematical, algorithmic and implementation structure transparent, not
%   the history of the implementation).
%   Build: docs/latex/build_cli_manual.sh <outdir>   (outdir in scratch,
%          never the tree: pdflatex, makeindex, pdflatex x2).
%
%   STATUS: NOT BINDING. This manual defines nothing mathematical. The
%   mathematics is the binding contract [C] (docs/latex/gl2_contract.tex),
%   cited as {\ttfamily[C,~<label>]} (scanned by
%   docs/matrix-engine/tools/label_census.py), and the extension
%   specification [E] (docs/latex/hecke_extensions.tex), cited as
%   \Ecite{<label>}. The console's binding design is [P]
%   (docs/console/00-console-interface-plan.md).
%
%   GENERATED PARTS: every file under docs/latex/cli_manual/ is generated
%   from the binary by docs/latex/tools/gen_cli_manual_tables.sh; never edit
%   one by hand. Every file docs/latex/cli_manual_runs/*.txt (a real
%   command-line output, \runoutput{<name>}) is generated by
%   docs/latex/tools/gen_cli_manual_runs.sh. CliManualCoverageValidation
%   (fast) fails when a table is stale, when a parser gains a token without
%   a manual row, when an \input of a generated file is missing, when an
%   excerpt line below does not occur in its script, when a quoted
%   transcript or run output does not exist, when a verb has no worked
%   example (allow-list docs/latex/cli_manual_gate.tsv), or when an engine,
%   descent or L23 lever, or a listed rollback lever, is missing from the
%   environment chapter.
%
%   EXCERPTS: an \excerptfrom{<script>} line followed by an ocexcerpt block
%   quotes lines of src/hacs/Console/scripts/<script>
%   verbatim ("..." marks an omission); \octranscript{<name>} inputs the
%   frozen <name>.expected transcript itself, so no output is ever retyped.
%======================================================================
\documentclass[11pt,openany]{report}
\usepackage[T1]{fontenc}
\usepackage{lmodern}
\usepackage{microtype}
\DisableLigatures[-]{encoding=T1,family=tt*}
\usepackage[a4paper,margin=2.3cm]{geometry}
\usepackage{amsmath,amssymb}
\usepackage{booktabs}
\usepackage{array}
\usepackage{longtable}
\usepackage{enumitem}
\usepackage{xcolor}
\usepackage{fvextra}
\usepackage{multicol}
\usepackage[most]{tcolorbox}
\usepackage{makeidx}
\usepackage[colorlinks=true,linkcolor=blue,citecolor=blue,urlcolor=blue]{hyperref}
\usepackage{xurl}
\makeindex

\setlength{\parindent}{0pt}
\setlength{\parskip}{0.55ex plus 0.2ex}
\setcounter{tocdepth}{1}
\setlength{\emergencystretch}{3em}

\newcommand{\F}{\mathbb{F}}
\newcommand{\Z}{\mathbb{Z}}
\newcommand{\Q}{\mathbb{Q}}
\newcommand{\Fl}{\mathbb{F}_{\ell}}
\newcommand{\Sym}{\mathrm{Sym}}
\newcommand{\SL}{\mathrm{SL}}
\newcommand{\tors}{\mathrm{tors}}
\newcommand{\Ecite}[1]{{\ttfamily[E,~#1]}}
\newcommand{\Pcite}[1]{[P]~#1}
\newcommand{\app}{\texttt{hacs}}
\newcommand{\cli}[1]{\texttt{#1}}
\DeclareRobustCommand{\envname}[1]{\texttt{\detokenize{#1}}}
\newcommand{\env}[1]{\envname{#1}\index{#1@\envname{#1} (environment)}}
\newcommand{\envx}[1]{\texttt{#1}}
\newcommand{\file}[1]{\nolinkurl{#1}}
\newcommand{\vb}[1]{\texttt{#1}}
\newcommand{\clitok}[1]{\texttt{#1}}

% --- generated-table environments (docs/latex/cli_manual/*.tex) -------------
\newenvironment{tokentable}[2]{%
  \small
  \begin{longtable}{@{}>{\raggedright\arraybackslash}p{0.27\linewidth}>{\raggedright\arraybackslash}p{0.22\linewidth}>{\raggedright\arraybackslash}p{0.45\linewidth}@{}}
  \toprule \textbf{#1} & \textbf{#2} & \textbf{Meaning}\\ \midrule\endhead
  \bottomrule\endlastfoot}{%
  \end{longtable}}
\newenvironment{membertable}{%
  \small
  \begin{longtable}{@{}>{\raggedright\arraybackslash}p{0.28\linewidth}>{\raggedright\arraybackslash}p{0.12\linewidth}>{\raggedright\arraybackslash}p{0.54\linewidth}@{}}
  \toprule \textbf{Member} & \textbf{Kind} & \textbf{Description}\\ \midrule\endhead
  \bottomrule\endlastfoot}{%
  \end{longtable}}
\newenvironment{settingtable}{%
  \footnotesize\setlength{\tabcolsep}{3pt}
  \begin{longtable}{@{}>{\raggedright\arraybackslash}p{0.21\linewidth}>{\raggedright\arraybackslash}p{0.08\linewidth}>{\raggedright\arraybackslash}p{0.06\linewidth}>{\raggedright\arraybackslash}p{0.26\linewidth}>{\raggedright\arraybackslash}p{0.34\linewidth}@{}}
  \toprule \textbf{Key} & \textbf{Family} & \textbf{Kind} & \textbf{Default / source} & \textbf{Description}\\ \midrule\endhead
  \bottomrule\endlastfoot}{%
  \end{longtable}}
\newenvironment{verbgroup}[1]{%
  \section{Group \texttt{#1}}\label{verbgroup:#1}}{}
% #5: the overload's read_only tag ([P] v3.10 Sec. 18.8): ro / rw / untagged.
\newcommand{\ocromark}[1]{\def\ocrotmp{#1}\def\ocroro{ro}\def\ocrorw{rw}%
  \ifx\ocrotmp\ocroro\textsf{\footnotesize[served]}\else
  \ifx\ocrotmp\ocrorw\textsf{\footnotesize[not served]}\else\textsf{\footnotesize[untagged]}\fi\fi}
\newcommand{\ocverb}[5]{%
  \par\smallskip\noindent{\raggedright\hangindent=1.5em #2#4\hfill\ocromark{#5}\par}\nopagebreak
  {\leftskip=1.5em\noindent\small #3\par}}
\newcommand{\octype}[3]{%
  \section*{\texttt{#1}}#3\nopagebreak
  \addcontentsline{toc}{section}{\texttt{#1}}%
  \noindent #2\par\nopagebreak}
\newcommand{\ocnomembers}{\par\noindent(no members)\par}
\newcommand{\ocsetting}[1]{\texttt{#1}}
\newenvironment{validationlist}{%
  \begin{multicols}{2}\scriptsize\raggedright}{\end{multicols}}
\newcommand{\ocvalidation}[2]{\noindent\texttt{#1}\hfill\textsf{#2}\par}

% --- verbatim material --------------------------------------------------------
\DefineVerbatimEnvironment{clihelp}{Verbatim}{fontsize=\footnotesize,frame=single,framesep=2mm,rulecolor=\color{black!40},breaklines=true}
\DefineVerbatimEnvironment{ocexcerpt}{Verbatim}{fontsize=\small,frame=leftline,framesep=3mm,rulecolor=\color{blue!50},breaklines=true}
\DefineVerbatimEnvironment{shell}{Verbatim}{fontsize=\small,frame=single,framesep=2mm,rulecolor=\color{black!30},breaklines=true}
\DefineVerbatimEnvironment{spec}{Verbatim}{fontsize=\small,frame=leftline,framesep=3mm,rulecolor=\color{black!40}}
\newcommand{\scriptsdir}{../../src/hacs/Console/scripts}
\newcommand{\excerptfrom}[1]{\par\smallskip\noindent{\small\textit{From} \nolinkurl{Console/scripts/#1}:}\par\nopagebreak}
\newcommand{\octranscript}[1]{\par\noindent{\small\textit{Frozen transcript} \nolinkurl{Console/scripts/#1.expected}:}\par\nopagebreak
  \VerbatimInput[fontsize=\small,frame=single,framesep=2mm,rulecolor=\color{green!40!black},breaklines=true,breakanywhere=true]{\scriptsdir/#1.expected}}
% A real command-line output, generated by docs/latex/tools/gen_cli_manual_runs.sh (never typed).
\newcommand{\runoutput}[1]{\par\noindent{\small\textit{Real output} \nolinkurl{docs/latex/cli_manual_runs/#1.txt}:}\par\nopagebreak
  \VerbatimInput[fontsize=\footnotesize,frame=single,framesep=2mm,rulecolor=\color{orange!60!black},breaklines=true,breakanywhere=true,obeytabs=true,tabsize=4]{cli_manual_runs/#1.txt}}

\newtcolorbox{danger}{colback=red!4,colframe=red!60!black,fonttitle=\bfseries,title=Owner-gated: writes files}
\newtcolorbox{note}{colback=blue!3,colframe=blue!40!black,fonttitle=\bfseries,title=Note}

% GENERATED by docs/latex/tools/gen_cli_manual_tables.sh, i.e.
% `CLI_MANUAL_EMIT=<dir> hacs CliManualCoverageValidation`, from the
% registries of the build it runs in.  Do not edit by hand: regenerate when a
% verb, type, setting, validation or command-line token changes.
\newcommand{\CMNumVerbs}{165}
\newcommand{\CMNumOverloads}{178}
\newcommand{\CMNumVerbGroups}{7}
\newcommand{\CMNumTypes}{31}
\newcommand{\CMNumSettings}{43}
\newcommand{\CMNumMathSettings}{11}
\newcommand{\CMNumEngineSettings}{32}
\newcommand{\CMNumValidations}{294}
\newcommand{\CMNumFast}{75}
\newcommand{\CMNumStandard}{171}
\newcommand{\CMNumSlow}{48}
\newcommand{\CMNumAll}{246}
\newcommand{\CMNumReadOnly}{169}
\newcommand{\CMNumReadOnlyVerbs}{159}
\newcommand{\CMNumNotReadOnly}{9}


\title{\app: the command line and the OC console\\[0.4ex]\large A user manual}
\author{Fabian Januszewski\\ Universit\"at Paderborn\\ \texttt{fabian.januszewski@math.uni-paderborn.de}}
\date{Version 1.5.15, 2026-11-13 (restructured by subject; not binding; history in Appendix~\ref{app:versions})}

\begin{document}
\maketitle

\begin{abstract}
\begin{center}\emph{In memoriam Prof.\ Dr.\ G\"unter Harder (1938--2025).}\end{center}
\app{} computes the cohomology of $\SL(2,\Z)$ and of its congruence subgroups $\Gamma_0(N)$ with coefficients in the
symmetric powers $\Sym^n$ of the standard representation, over a finite field $\Fl$, over $\Z/\ell^k$ and, by a
descent, over $\Z$; the Hecke operators on it; their primary decompositions; and an extension layer of Hecke-module
invariants (maximal ideals, local Hecke algebras, Loewy structure, the Dickson $\theta$ map, Hida theory). This manual
explains what these objects are, how the program computes them, and how to drive it: the command line, the
interactive and scripted console with its language OC, campaigns, the residue atlas, benchmarks and the engine
switches. Every example is a real output of the program --- a frozen console transcript or a command run by a
generator script --- and the reference tables of the \CMNumVerbs{} console verbs, \CMNumTypes{} object types,
\CMNumSettings{} settings and \CMNumValidations{} validations are generated from the binary. It is written for a
mathematician who wants to compute, not for someone changing the code.
\end{abstract}

\tableofcontents

%======================================================================
\chapter{Getting started}\label{ch:start}
%======================================================================

\section{What the program computes}\label{sec:overview}

\paragraph{What the program computes.}
Let $\Gamma=\SL(2,\Z)$ and let $M_n=\Sym^n(\Fl^2)$ be the $n$-th symmetric power of the standard
representation, with the action fixed by {\ttfamily[C,~def:rhon]}. \app{} computes the equivariant
cohomology $H^q(\Gamma; M_n)$, $q\in\{0,1,2\}$, on a seven-vertex model of the orbifold
$\Gamma\backslash\mathfrak{H}$; for $6\in\Fl^\times$ this is group cohomology
({\ttfamily[C,~rmk:orbifold]}). On top it computes the Hecke operators $T_p$, $U_\ell$ (the same code
path at $p=\ell$), $T_{p^r}$ and $T(p,p)$ on cohomology ({\ttfamily[C,~def:hecke]},
{\ttfamily[C,~thm:relations]}, {\ttfamily[C,~prop:diamond]}), the certified primary decomposition of
$H^q$ over a commuting Hecke family, and the extension layer of [E]: the maximal ideals of the Hecke
algebra, local Hecke algebras $\mathbb{T}_\mathfrak{m}$, Loewy and socle series, the cone and boundary
cohomology, the involution $\iota^*$ ({\ttfamily[C,~def:iota]}), the Dickson long exact sequence
(\Ecite{sec:thetaobj}) and the $\Z/\ell^k$ coefficient route (\Ecite{sec:ring}).
All arithmetic is exact: residues modulo a prime $\ell$ (up to $2^{64}$ in places), or modulo $\ell^k$.

\paragraph{One binary, several front ends.}
Everything is one executable, \app, whose first argument selects a front end
(Section~\ref{sec:cli}):
\begin{itemize}[nosep]
\item the \textbf{validation harness} (\cli{--list}, \cli{--all}, \cli{--all-full}, validation names) --
  the project's test suite, which is also the most reliable way to reproduce a published number;
\item the \textbf{benchmark driver} \cli{--bench}: one pipeline point with timings;
\item the \textbf{campaign driver} \cli{--campaign}: a parameter sweep that writes the mathematics to
  JSON and TSV;
\item the \textbf{atlas builder} \cli{--atlas}: the residue atlas $\Sigma(\ell)$;
\item the \textbf{console} (\cli{--console}, \cli{--script}, \cli{-c}, \cli{-}): an interactive and
  scriptable expression language, OC, over the whole pipeline (Chapter~\ref{ch:console}).
\end{itemize}


\paragraph{How this manual is organised.}
By subject, from the mathematics to the machinery. Chapter~\ref{ch:objects} introduces the objects: the group, the
model, the local systems $\Sym^n$, the coefficient rings $\Fl$, $\Z/\ell^k$, $\Z$, and the levels $\Gamma_0(N)$.
Chapters~\ref{ch:field}--\ref{ch:ext} compute with them in increasing depth: cohomology over a field, Hecke operators,
integral cohomology by descent, and the extension layer with Hida theory. Chapter~\ref{ch:console} is the console and its
language; Chapter~\ref{ch:campaigns} the campaign driver for sweeps; Chapter~\ref{ch:perf} benchmarks and every
environment variable. Chapter~\ref{ch:reference} has a glossary and troubleshooting, and the appendices are the generated
reference tables. Command-line tokens, console verbs, types, settings and environment variables are indexed by name.

\paragraph{Sources of truth.}
This manual restates; it does not legislate. The mathematics is the binding contract [C]
(\file{docs/latex/gl2_contract.tex}) and the extension specification [E]
(\file{docs/latex/hecke_extensions.tex}), cited by label. The console design is [P]
(\file{docs/console/00-console-interface-plan.md}); \file{docs/console/console.md} is the console's own
generated Markdown reference; the campaign format is \file{docs/matrix-engine/campaign.md}; the
per-validation encyclopaedia is \file{docs/validation-catalogue.md}; benchmark reference figures live
in \file{docs/performance-reference.md}. Where this manual and the code disagree, the code wins and
\vb{CliManualCoverageValidation} should have failed (Appendix~\ref{app:regen}).

\begin{note}
\textbf{Generated parts.} The usage texts in boxes, the token tables, and all of Appendices~\ref{app:verbs}--\ref{app:validations} are generated from the binary
by \file{docs/latex/tools/gen_cli_manual_tables.sh}; the counts in the prose
(\CMNumVerbs{} verbs, \CMNumValidations{} validations, \dots) are generated too. They describe the build
the manual was generated from.
\end{note}

\section{Quick start}\label{sec:quickstart}

Build once (Section~\ref{sec:build} has the details), then from the build directory:

\begin{shell}
./hacs --help                     # the top-level usage
./hacs --list | head              # registered validations and tiers
./hacs mod_arithmetic_validation  # run one validation; exit code 0 = pass
./hacs --bench hecke 48 134217757 2      # T_2 on H^*(SL(2,Z); Sym^48), one BENCH line
./hacs -c 'print(cohomology(symmetric_power(sl2z_model(), 10, 691)).dims)'
./hacs --console                  # interactive session, prompt "oc> "
\end{shell}
The build also creates \texttt{cohomology\_app}, a symbolic link to \texttt{hacs} (the program's former name), so older
scripts keep working; this manual uses \texttt{hacs} throughout.

The one-liner above computes $\dim H^q(\Gamma;\Sym^{10}\F_{691}^2)$ for $q=0,1,2$; the frozen
transcript of script \file{001_dimensions.oc} (Section~\ref{ex:dims}) shows the answer \texttt{[0, 3, 0]}.
A typical console session builds the flagship cell $(n,\ell)=(10,691)$ step by step:

\excerptfrom{006_flagship.oc}
\begin{ocexcerpt}
L = symmetric_power(MODEL, N_SYM, ELL);
H = cohomology(L);
print(H.dims);
T2 = hecke(H, P_HECKE);
A = T2.A(1);
print(charpoly(A));
D = decompose(T2, 1);
print(D.num_blocks);
\end{ocexcerpt}

(\vb{MODEL}, \vb{N\_SYM}, \vb{ELL}, \vb{P\_HECKE} are bound by the shared \file{prelude.oc} to
\vb{sl2z\_model()}, $10$, $691$, $2$.) Section~\ref{ex:flagship} quotes the full transcript.

\paragraph{Odd $n$.} On this model $H^\bullet(\Gamma;\Sym^n)$ vanishes for odd $n$, because $-I$ acts by
$(-1)^n$ ({\ttfamily[C,~rmk:negI]}). Use an even $n$ for anything you want to look at; the bench usage
text says so too.

\paragraph{Limits.} \cli{--bench} parses $\ell$ (and $n$, $p$, $r$) as 32-bit unsigned integers; a larger
value is refused (with the misleading message ``must be non-negative integers''). A campaign spec takes a
64-bit $\ell$ (but refuses \texttt{coefficient\_ring = Z/ell\^{}k} with $\ell^k\ge2^{64}$ together with
\texttt{ext = true}). \cli{--atlas} needs $5\le\ell\le100000$. The $U_\ell$ operator needs $\ell\le23170$
({\ttfamily[C,~lem:bounds]}). \texttt{coefficient\_ring = Z/ell\^{}<k>} needs $2\le k\le4096$.

\section{Building the program}\label{sec:build}

The build needs CMake $\ge 3.18$ ($\ge 3.21$ for the presets), a C++17 compiler and, unless it is configured
CPU-only (below), the CUDA toolkit (\texttt{nvcc}); the CUDA layer is optional at run time and never used by
default (Section~\ref{sec:env-engine}).

\begin{shell}
cmake -B build -S . -DCMAKE_BUILD_TYPE=RelWithDebInfo   # out-of-source (the only kind)
make -C build -j8
./build/hacs --list
\end{shell}

The sources live under \file{src/}. Builds are out-of-source only: an in-source configure (\verb|cmake .| in the
repository root) is refused with a one-line hint. A build directory's \file{CMakeCache.txt} pins the configuration it
was first made with, so start a fresh build directory (or remove its \file{CMakeCache.txt} and \file{CMakeFiles/})
before changing it.

\begin{center}\small
\begin{tabular}{@{}ll>{\raggedright\arraybackslash}p{0.33\linewidth}@{}}
\toprule
\textbf{\texttt{CMAKE\_BUILD\_TYPE}} & \textbf{Flags} & \textbf{Use}\\\midrule
\texttt{RelWithDebInfo} & \verb|-O2 -g -DNDEBUG| & the default; everyday runs and validations\\
\texttt{Debug} & \verb|-O0 -g| & debugging (much slower)\\
\texttt{Release} & \verb|-O3 -DNDEBUG| & --\\
\texttt{Bench} & \verb|-O3 -march=native -DNDEBUG| & timing runs\\
\texttt{Check} & \verb|-O1 -g -fsanitize=address,undefined| & sanitizer runs (debug checks live)\\
\texttt{CheckThread} & \verb|-O1 -g -fsanitize=thread| & ThreadSanitizer (run with \env{ENGINE_CUDA_DISABLE}\texttt{=1})\\
\bottomrule
\end{tabular}
\end{center}

\paragraph{CMake options.} None of them changes a published value:
\cli{-DHECKE\_PURE\_INTEGER=ON} disables the certified floating-point filter of the Hecke sign
predicates (bit-identical results, slower); \cli{-DENGINE\_SHADOW\_LEGACY=ON} runs engine and legacy
routes side by side and throws on the first differing entry; \cli{-DENGINE\_VERIFY\_DEVICE=ON}
re-executes every GPU node on the CPU and compares; \cli{-DENGINE\_COUNT\_OBJECTS=ON} adds object
counters to the \texttt{BENCH} line; \cli{-DOC\_CUDA\_ARCHITECTURES=\dots} overrides the GPU
architecture (default \texttt{120}; \texttt{87} on a Jetson Orin).

\paragraph{Build switches and presets.} Three options are decided before anything is compiled; a component
switched off is neither required nor probed. \cli{-DOC\_ENABLE\_CUDA=OFF} (default \texttt{ON}) builds CPU-only:
no CUDA toolkit is needed, the device layer is a stub that reports itself unavailable, and the device
validations print a \texttt{[SKIP]} line. \cli{-DOC\_BUILD\_DOCS=ON} (default \texttt{OFF}; needs pdflatex,
bibtex, makeindex) adds the targets \texttt{make docs} and \texttt{make docs\_<name>}, which build this manual
and the contracts into \texttt{<build>/docs-pdf}; they are never part of \texttt{all}.
\cli{-DOC\_BUILD\_NOTEBOOK=ON} (default \texttt{OFF}) compiles the served console \cli{--serve}
(Section~\ref{sec:serve}), and \texttt{cmake --install} puts the notebook gateway and page under
\file{<prefix>/share/oc-notebook/}. No switch changes a value, and the \cli{--list} count never depends on
one: a validation excluded by a switch keeps its name and prints \texttt{[SKIP] <name>: OC\_BUILD\_<X>=OFF}
(exit~0). \file{CMakePresets.json} holds \texttt{minimal} (a plain configure), \texttt{cpu-only},
\texttt{reference} (CUDA \texttt{sm\_120}, docs and notebook on), \texttt{jetson-orin} and
\texttt{jetson-orin-cpu} (NVIDIA Jetson Orin, aarch64), and \texttt{termux} / \texttt{aarch64-cpu} (CPU only,
\texttt{Release -O2}, stripped: Android under Termux and other small aarch64 devices):
\verb|cmake --preset cpu-only -B <dir>|. The release's \file{BUILDING.md} gives the exact commands for the four
targets (CUDA, CPU only, Jetson Orin, Termux) and what was tested on each.

\paragraph{The source path is baked in.} The binary records the source directory it was configured
from (\texttt{OC\_SOURCE\_DIR}); validations that read tracked files (the golden reference, the console
scripts, this manual) read them there. Do not delete or move that checkout while a binary built from
it is running.

\section{The command line}\label{sec:cli}

\app{} looks only at its \emph{first} argument to decide what to do; everything after it belongs to
the selected front end. The printed top-level usage (\cli{--help}, and the same text when no argument is
given) is frozen text; it names every front end, \cli{--atlas} and the console forms included:

% GENERATED by docs/latex/tools/gen_cli_manual_tables.sh, i.e.
% `CLI_MANUAL_EMIT=<dir> hacs CliManualCoverageValidation`, from the
% registries of the build it runs in.  Do not edit by hand: regenerate when a
% verb, type, setting, validation or command-line token changes.
\begin{clihelp}
Usage:
  ./hacs                 With no argument: print this usage text; exit 1.
  ./hacs --help          Print this usage text; exit 0.
  ./hacs --version       Print the version line (`hacs <version> (commit ...; contract ...)`); exit 0.
  ./hacs --list          List all registered validation names with their tier.
  ./hacs --all           Run the fast + standard tiers; exit code = number of failures.
  ./hacs --all-full      Run every tier including slow; exit code = number of failures.
  ./hacs <Name>          Run one registered validation by name; exit code = its own return code.
  ./hacs --bench <mode> <args...>
                                  Run the benchmark CLI (`--bench --help` for modes).
  ./hacs --campaign <spec-file> [options]
                                  Run a parameter sweep (`--campaign --help`; docs/matrix-engine/campaign.md).
  ./hacs --atlas <ell> [options]
                                  Build, close, verify or index a residue atlas (`--atlas --help`).
  ./hacs --console | -i | --script <file.oc> | -c '<code>' | - [modifiers]
                                  The OC script console (`--console --help`; docs/console/console.md).
  ./hacs --serve [options]
                                  The served read-only console, NDJSON on stdio (OC_BUILD_NOTEBOOK=ON builds only).
\end{clihelp}


The complete dispatch table of \file{main.cpp}, in the order it is tested:

% GENERATED by docs/latex/tools/gen_cli_manual_tables.sh, i.e.
% `CLI_MANUAL_EMIT=<dir> hacs CliManualCoverageValidation`, from the
% registries of the build it runs in.  Do not edit by hand: regenerate when a
% verb, type, setting, validation or command-line token changes.
\begin{tokentable}{Token}{Arguments}
\clitok{(none)} & -- & Prints the top-level usage text (as \texttt{-{}-{}help}) and exits 1, the usage-error code; builds no validation registry. (Until S24 CL-7 it ran the legacy restriction-morphism benchmark.) \\
\clitok{-{}-{}version}\index{version@\texttt{-{}-{}version} (command line)} & -- & Prints the one-line version \texttt{hacs \textless{}version\textgreater{} (commit \textless{}sha12\textgreater{}; contract [C] \textless{}x\textgreater{},\allowbreak{} [E] \textless{}y\textgreater{},\allowbreak{} [U] \textless{}z\textgreater{})} (the \texttt{VERSION} file, the configure-time commit, the contract versions) and exits 0; reads nothing else, builds no registry. Metadata only: no golden object, record hash or job id reads it. \\
\clitok{-{}-{}list}\index{list@\texttt{-{}-{}list} (command line)} & -- & Prints every registered validation as \texttt{name\textless{}TAB\textgreater{}[tier]}, in registration order; exit 0. \\
\clitok{-{}-{}all}\index{all@\texttt{-{}-{}all} (command line)} & \texttt{[-{}-{}shard \textless{}i\textgreater{}/\textless{}N\textgreater{}] [-{}-{}list]} & Runs the \texttt{fast} and \texttt{standard} tiers; exit code = number of failures; \texttt{slow} names are skipped and counted. \\
\clitok{-{}-{}all-{}full}\index{all-full@\texttt{-{}-{}all-{}full} (command line)} & \texttt{[-{}-{}shard \textless{}i\textgreater{}/\textless{}N\textgreater{}] [-{}-{}list]} & Runs every tier, \texttt{slow} included; exit code = number of failures. \\
\clitok{-{}-{}shard}\index{shard@\texttt{-{}-{}shard} (command line)} & \texttt{\textless{}i\textgreater{}/\textless{}N\textgreater{}} & (after \texttt{-{}-{}all} / \texttt{-{}-{}all-{}full}; S24 SHARD) runs shard i of N (1-based, N \textless{}= 64) of that tier's roster in this process: a deterministic partition balanced by expected wall time, every device name in one shard, registry order within a shard; the union over i is the roster, each name once. Prints \texttt{SHARD i/N names=\textless{}k\textgreater{} of \textless{}total\textgreater{} ...} first; exit code = its failures; with \texttt{-{}-{}list} prints the shard's names instead (exit 0); a malformed i/N is a usage error (exit 1). Driver: \texttt{docs/matrix-{}engine/tools/run\_all\_sharded.sh}. Not in the frozen help text. \\
\clitok{-{}-{}calibrate}\index{calibrate@\texttt{-{}-{}calibrate} (command line)}, \clitok{-{}-{}calibrate=thorough}\index{calibrate=thorough@\texttt{-{}-{}calibrate=thorough} (command line)} & \texttt{[-{}-{}repeat \textless{}R\textgreater{}] [-{}-{}threads \textless{}T\textgreater{}] [-{}-{}taskset \textless{}cpus\textgreater{}] [-{}-{}profile \textless{}path\textgreater{} \textbar{} -{}-{}no-{}write] [-{}-{}emit-{}reference]} & (NBCONFIRM S3) The per-device speed calibration of the console cost estimates: runs the fixed probe set of the embedded \texttt{calib} rows (six verb anchors inside today's silent envelopes, four fingerprints) as child processes at 1 thread and at min(16, CPUs) threads (\texttt{=thorough}: 3 repeats; default 1, about 20 s on the reference box), prints one \texttt{CALIB} row per probe, \texttt{NBCONFIRM\_CALIB ... verdict=OK} and the drift line \texttt{COST\_MODEL\_DRIFT ... verdict=OK\textbar{}DRIFT} (meaningful on the reference box only), and writes the device profile \texttt{\$\{HACS\_DEVICE\_PROFILE:-{}\$\{XDG\_CONFIG\_HOME:-{}\textasciitilde{}/.config\}/hacs/device\_profile.tsv\}} (\texttt{-{}-{}emit-{}reference}: prints \texttt{calib} rows instead). Advisory only: never a value, a record or a hash. Not in the frozen help text. \\
\clitok{\textless{}Name\textgreater{}} & \texttt{[\textless{}Name\textgreater{} ...]} & Runs the named validations in this process, in order; exit code = number of failures (an unknown name is a failure). \\
\clitok{-{}-{}bench}\index{bench@\texttt{-{}-{}bench} (command line)} & \texttt{\textless{}mode\textgreater{} \textless{}args...\textgreater{}} & The benchmark CLI: one pipeline point, phase timings, peak RSS and one \texttt{BENCH} line. \\
\clitok{-{}-{}campaign}\index{campaign@\texttt{-{}-{}campaign} (command line)} & \texttt{\textless{}spec-{}file\textgreater{} [options]} & The campaign driver: a parameter sweep of Hecke jobs with streamed JSON output. \\
\clitok{-{}-{}atlas}\index{atlas@\texttt{-{}-{}atlas} (command line)} & \texttt{\textless{}ell\textgreater{} [options] \textbar{} -{}-{}close \textless{}f\textgreater{} \textbar{} -{}-{}verify \textless{}f\textgreater{} \textbar{} -{}-{}index \textless{}f\textgreater{}} & Builds, closes, verifies or indexes a residue atlas Sigma(ell) outside the suite (\texttt{-{}-{}atlas -{}-{}help} for its options). \\
\clitok{-{}-{}help}\index{help@\texttt{-{}-{}help} (command line)}, \clitok{-{}h}\index{h@\texttt{-{}h} (command line)} & -- & Prints the top-level usage text; exit 0. \\
\clitok{-{}-{}console}\index{console@\texttt{-{}-{}console} (command line)}, \clitok{-{}i}\index{i@\texttt{-{}i} (command line)} & \texttt{[modifiers]} & Interactive OC session (prompt \texttt{oc\textgreater{} }). \\
\clitok{-{}-{}script}\index{script@\texttt{-{}-{}script} (command line)} & \texttt{\textless{}file.oc\textgreater{} [name=value ...] [modifiers]} & Runs an OC script file. \\
\clitok{-{}c}\index{c@\texttt{-{}c} (command line)} & \texttt{'\textless{}code\textgreater{}' [modifiers]} & Runs one line of OC script. \\
\clitok{-{}}\index{-@\texttt{-{}} (command line)} & \texttt{[modifiers]} & Runs the OC script read from standard input (stdin is the script, so no REPL is ever opened). \\
\clitok{-{}-{}serve}\index{serve@\texttt{-{}-{}serve} (command line)} & \texttt{[-{}-{}set key=value ...] [-{}-{}rss-{}cap-{}mib N] [-{}-{}max-{}reply-{}bytes N]} & The served console (protocol \texttt{oc-{}serve/1}): one read-only session driven by NDJSON on stdin / stdout; built only with \texttt{OC\_BUILD\_NOTEBOOK=ON} (otherwise an unknown token, exit 1). \\
\clitok{-{}-{}serve-{}atlas}\index{serve-atlas@\texttt{-{}-{}serve-{}atlas} (command line)} & \texttt{-{}-{}store \textless{}store.json\textgreater{} [-{}-{}index \textless{}f\textgreater{}] [-{}-{}expect-{}fnv 0x\textless{}hex\textgreater{}] [-{}-{}max-{}reply-{}bytes N] [-{}-{}rss-{}cap-{}mib N]} & Read-only queries on a \texttt{hecke\_atlas/v1} store (protocol \texttt{oc-{}serve-{}atlas/1}) by NDJSON on stdin / stdout: the header, cells, one system, filtered pages of systems; built only with \texttt{OC\_BUILD\_NOTEBOOK=ON} (otherwise an unknown token, exit 1). \\
\end{tokentable}


\paragraph{Rules that follow from the dispatch.}
\begin{itemize}[nosep]
\item With \emph{no} argument the program prints the usage text above and exits with code~1 (the usage
  error); it builds no validation registry. (Until slot~24 it ran a fixed legacy benchmark.)
\item Any first argument that is not one of the tokens above is taken as a \emph{validation name}, and
  so is every further argument: \cli{./hacs NameA NameB} runs both in one process.
  An unknown name counts as one failure.
\item The console modifiers (\cli{--set}, \cli{--log}, \cli{--quiet}, \cli{--verbose}, \texttt{name=value})
  are only recognised \emph{after} a console form: \cli{./hacs --script f.oc --set u\_ell=off}
  works, \cli{./hacs --set u\_ell=off --script f.oc} looks for a validation called
  \cli{--set}.
\item \cli{--bench micro \dots} and a GPU device policy in the environment (\envx{ENGINE\_DEVICE=gpu|hybrid},
  or \env{ENGINE_CUDA_PROBE}) install the CUDA hooks; every other run never creates a CUDA context.
\end{itemize}

\paragraph{Exit codes by front end.}
\begin{center}\small
\begin{tabular}{@{}>{\raggedright\arraybackslash}p{0.24\linewidth}>{\raggedright\arraybackslash}p{0.70\linewidth}@{}}
\toprule
\textbf{Front end} & \textbf{Exit code}\\\midrule
validations, \cli{--all}, \cli{--all-full} & the number of failing validations (the shell sees it modulo 256)\\
\cli{--list}, \cli{--help} & 0\\
\cli{--bench} & 0 success; 1 usage error; 2 an exception, or a \cli{micro} op whose backend is unavailable\\
\cli{--campaign} & 0 success; 1 usage or spec error; 2 at least one job failed\\
\cli{--atlas} & 0 success; 1 usage; 2 \cli{--verify} mismatch; 3 an exception; 5 a non-field coefficient ring in force\\
console forms & 0 ok; 1 usage; 2 parse error; 3 runtime error; 4 failed \vb{assert} / \vb{expect\_error}; 5 a legal setting this build cannot execute (Section~\ref{sec:errors})\\
\bottomrule
\end{tabular}
\end{center}

\section{The validation harness}\label{sec:validations}

There is no external test framework. A \emph{validation} is a registered function returning $0$ on
success; it prints its own \texttt{[PASS]}/\texttt{[FAIL]} lines. This build registers
\CMNumValidations{} of them, each with a \emph{tier}:

\begin{center}\small
\begin{tabular}{@{}lrl@{}}
\toprule
\textbf{Tier} & \textbf{Count} & \textbf{Run by}\\\midrule
\texttt{fast} & \CMNumFast & \cli{--all}, \cli{--all-full}, by name\\
\texttt{standard} & \CMNumStandard & \cli{--all}, \cli{--all-full}, by name\\
\texttt{slow} & \CMNumSlow & \cli{--all-full}, by name\\
\bottomrule
\end{tabular}
\end{center}

So \cli{--all} runs \CMNumAll{} validations and \cli{--all-full} all \CMNumValidations{}. The full
roster, sorted by name, is Appendix~\ref{app:validations}; what each one checks is in
\file{docs/validation-catalogue.md}.

\begin{itemize}[nosep]
\item \cli{--list} prints \texttt{name<TAB>[tier]} in registration order -- the same list the console
  verbs \vb{validations()} and \vb{validations("fast")} return.
\item An exception escaping a validation (an \texttt{ENGINE\_CHECK}, an out-of-range argument, \dots) is
  reported as that validation's failure with an \texttt{[EXCEPTION] <name>: \dots} line, and the run
  continues.
\item Validations armed by an environment variable print one \texttt{EFFECTIVE <name> arm=<X> cells=<k>}
  line when the arm actually runs. \textbf{Read that line} before trusting a sharded or cluster run: an
  environment variable that was silently ignored shows up as a missing or vacuous \texttt{EFFECTIVE}
  line.
\item Some validations are \emph{interlocked}: they print \texttt{[SKIP]} even under \cli{--all-full}
  unless their variable is set (Section~\ref{sec:env-gated}).
\item The Hecke acceptance gate (\vb{HeckeEndToEndValidation}, \vb{HeckeBadPrimeValidation},
  \vb{HeckeBenchmarkSmoke}) needs \cli{--all-full} or the three names, since
  \vb{HeckeBenchmarkSmoke} is \texttt{slow}.
\end{itemize}

From the console, \vb{run\_validation("Name")} runs one validation in process and returns its code; the
owner-gated emitters refuse there (Section~\ref{ex:harness}).

\section{How the examples in this manual are made}\label{sec:examples}

Two kinds of example appear, and neither is ever typed by hand. A \emph{console example} is part of one of the scripts
of \file{src/hacs/Console/scripts/}, which \vb{ConsoleScriptValidation} runs on every \cli{--all};
the output shown is the script's frozen transcript, input from the \file{.expected} file itself, and a quoted line of the
script is checked against the script by \vb{CliManualCoverageValidation}. A \emph{command-line example} (``Real output'')
is written by \file{docs/latex/tools/gen_cli_manual_runs.sh}, which runs the command shown on its first line (\texttt{\$
\dots}) against a built binary and appends its exit code; timings and memory figures in it are those of the reference
machine and are never part of a result. To run a console script yourself:
\begin{shell}
cd src/hacs/Console/scripts
../../../../build/hacs --script 006_flagship.oc
\end{shell}
Most scripts begin with \texttt{include("prelude.oc")}, which binds the flagship cell $(n,\ell)=(10,691)$ and checks that
the mathematical settings are at their defaults:

\excerptfrom{prelude.oc}
\begin{ocexcerpt}
N_SYM = 10;
ELL = 691;
P_HECKE = 2;
POINT_COH = "n=10 ell=691";
POINT_HECKE = "n=10 ell=691 p=2 r=1";
assert(fingerprint() == "0x0000000000000000",
"the smoke scripts pin the DEFAULT mathematical configuration");
MODEL = sl2z_model();
\end{ocexcerpt}

%======================================================================
\chapter{The mathematical objects}\label{ch:objects}
%======================================================================

This chapter fixes the vocabulary of the rest of the manual. The definitions are the binding contract's [C]; what
follows is a reader's guide to them, with the labels to look up.

\section{The group, the orbifold and the model}\label{sec:model}

Let $\Gamma=\SL(2,\Z)$. It acts on the upper half plane $\mathfrak H$ with finite stabilisers: $\pm I$ everywhere, a
cyclic group of order~4 at $i$ (generated by $S=\begin{pmatrix}0&-1\\1&0\end{pmatrix}$) and of order~6 at
$e^{2\pi i/3}$ (generated by $U=ST$). The quotient $\Gamma\backslash\mathfrak H$ is an \emph{orbifold}: a disc with two
cone points and a cusp. \app{} works on one fixed triangulation of it, the \emph{seven-vertex model}
(\vb{sl2z\_model()}, {\ttfamily[C,~rmk:orbifold]}), on which a local system is a coefficient module at each simplex
together with the monodromy relations between them. Two facts govern everything else:
\begin{itemize}[nosep]
\item \textbf{Away from $2$ and $3$ the model computes group cohomology.} The finite stabilisers have orders $2$, $4$
  and $6$; once $6$ is invertible in the coefficients, the equivariant cohomology of the model is $H^\bullet(\Gamma;M)$
  ({\ttfamily[C,~rmk:orbifold]}, {\ttfamily[C,~prop:bredon]}). Over $\F_2$ and $\F_3$ the same complex computes the
  \emph{Bredon} cohomology of the model, a different object; the group cohomology at $\ell=2,3$ is reached by the
  integral descent (Chapter~\ref{ch:integral}) and by the orbifold census (Section~\ref{sec:orbifold}).
\item \textbf{The model is finite.} Its cochain complex lives in degrees $0,1,2$, while the group cohomology over $\Z$
  continues forever in degrees $\ge2$ (periodically, with period~2, through the cup product with an Euler class). Every
  object therefore records its degrees, its top degree $q_{\max}$ and whether it is truncated.
\end{itemize}

\section{Local systems: the symmetric powers}\label{sec:localsys}

The coefficient module is $\Sym^n$ of the standard representation of $\Gamma$ on $R^2$: homogeneous polynomials of
degree $n$ in two variables, $\gamma\in\Gamma$ acting by substitution, with the convention fixed in
{\ttfamily[C,~def:rhon]} ($\rho_n$). Its rank is $d_V=n+1$. The weight of the corresponding modular forms is $k=n+2$:
by the Eichler--Shimura isomorphism, over a field $\Fl$ with $\ell>n+1$,
\[
  \dim H^1(\Gamma;\Sym^n\Fl^2) \;=\; 2\dim S_{n+2}(\SL_2(\Z)) + 1 \qquad (n>0 \text{ even}),
\]
the cusp forms counted twice and the Eisenstein series once. Because $-I$ acts on $\Sym^n$ by $(-1)^n$,
\textbf{everything vanishes for odd $n$} ({\ttfamily[C,~rmk:negI]}) over a field of characteristic $\ne2$; use an even
$n$. In the console, \vb{symmetric\_power(sl2z\_model(), n, ell)} builds the local system; the presets of
Section~\ref{ex:presets} put other local systems on other spaces.

\section{Coefficient rings}\label{sec:rings}

One mathematical setting, \vb{coefficient\_ring} (Section~\ref{sec:mathsettings}), chooses where the coefficients live.
It is a property of the whole session or campaign, recorded in every result:
\begin{center}\small
\begin{tabular}{@{}>{\raggedright\arraybackslash}p{0.15\linewidth}>{\raggedright\arraybackslash}p{0.38\linewidth}>{\raggedright\arraybackslash}p{0.39\linewidth}@{}}
\toprule
\textbf{Value} & \textbf{What is computed} & \textbf{Where}\\\midrule
\vb{F\_ell} (default) & vector spaces over the prime field $\Fl$ of the verb's $\ell$: dimensions, Hecke matrices, decompositions & every verb; Chapters~\ref{ch:field}, \ref{ch:hecke}, \ref{ch:ext}\\
\vb{Z/ell\^{}<k>} & finite $\Z/\ell^k$-modules: lengths and invariant factors, Hecke operators over $\Z/\ell^k$ (\Ecite{sec:ring}) & the pipeline verbs and most extension verbs; Section~\ref{sec:ringcoh}\\
\vb{Z} & finitely generated abelian groups $\Z^r\oplus\tors$ ({\ttfamily[C,~def:ringsemantics]}(3)) & the verb \vb{descent} and descent campaigns only; Chapter~\ref{ch:integral}\\
\bottomrule
\end{tabular}
\end{center}
A verb that does not honour the ring in force stops with exit code~5 and names itself; it never silently computes over
another ring. A non-default ring moves the settings fingerprint, so results over different rings never compare equal.

\section{Levels and characters: \texorpdfstring{$\Gamma_0(N)$}{Gamma0(N)} and nebentypus}\label{sec:levels}

For a level $N$, $\Gamma_0(N)\subset\Gamma$ is the subgroup of matrices with lower-left entry divisible by $N$, of index
$\psi(N)=N\prod_{p\mid N}(1+1/p)$. Its cohomology is computed on the \emph{same} seven-vertex model, by Shapiro's lemma:
$H^\bullet(\Gamma_0(N);\Sym^n)=H^\bullet(\Gamma;\mathrm{Ind}\,\Sym^n)$, the induced module being the permutation module on
the $\psi(N)$ cosets tensored with $\Sym^n$ ({\ttfamily[C,~sec:levelN]}). A \emph{nebentypus} twists the induced module by a
Dirichlet character $\chi$ modulo $N$, given by its Conrey label $N.m$ and reduced into $\Fl$. The verbs \vb{level\_n(N, n,
ell)} and \vb{nebentypus(N, m, n, ell)} return $H^0$, $H^1$ and the Hecke operators $T_p$ (and $U_p$ for $p\mid N$) of the
level; the descent of Chapter~\ref{ch:integral} takes a level too. At level~1 the dimension is the Eichler--Shimura count
above; at level~11 and weight~2 the two classes of the elliptic curve $X_0(11)$ appear beside one Eisenstein class (the
two cusps give $c-1=1$ Eisenstein class in weight~2, so $h^1=3=2g+(c-1)$):

\excerptfrom{052_s22_level_n_nebentypus.oc}
\begin{ocexcerpt}
A = level_n(1, 10, 7);
print("level 1, n = 10, ell = 7: h1 =", A.h1, " dim S_12 =", A.dim_cusp_forms, " charpoly T_2 =", A.ops[0].charpoly);
...
B = level_n(11, 0, 7);
print("level 11, n = 0, ell = 7: h0 =", B.h0, " h1 =", B.h1, " index =", B.index, " cusps =", B.cusps);
...
C = nebentypus(11, 1, 0, 7);
...
D = nebentypus(7, 2, 2, 13);
print("7.2, n = 2, ell = 13: order =", D.order, " h0 =", D.h0, " h1 =", D.h1, " charpoly T_2 =", D.ops[0].charpoly);
...
E = nebentypus(13, 8, 4, 17);
print("13.8, n = 4, ell = 17: parity_ok =", E.parity_ok, " h1 =", E.h1);
\end{ocexcerpt}
\octranscript{052_s22_level_n_nebentypus}
The characteristic polynomials of these verbs are coefficient lists \emph{from the constant term up}: \texttt{[4, 0, 1, 1]}
is $x^3+x^2+4$ over $\F_7$, i.e.\ $(x-4)^2(x-5)$ with $4\equiv\tau(2)=-24$ and $5\equiv1+2^{11}=2049$; the
trivial character reproduces $\Gamma_0(11)$ operator by operator; a character of the wrong parity ($\chi(-1)\ne(-1)^n$)
gives zero.

%======================================================================
\chapter{Cohomology over a field}\label{ch:field}
%======================================================================

The basic computation: build a local system, ask for its cohomology. A pipeline verb returns a lazy \emph{handle}; the work
runs on the first access to a computed member (Section~\ref{sec:types}).

\section{Dimensions and the golden reference}\label{ex:dims}

The scripts never compare with a copied number: every pin is read from the frozen golden reference with
\vb{golden(point, object)}, which refuses at a non-default fingerprint.

\excerptfrom{001_dimensions.oc}
\begin{ocexcerpt}
L = symmetric_power(MODEL, N_SYM, ELL);
H = cohomology(L);
assert(H.degrees == [0, 1, 2], "degrees = [0 .. q_max]");
for q in H.degrees {
assert(H.dims[q] == golden(POINT_COH, "dim_H^" + str(q)).value, "dim H^" + str(q));
assert(H.dim_C(q) == golden(POINT_COH, "dim_C_eq^" + str(q)).value, "dim C^" + str(q));
}
...
print(H.dims, H.q_max, H.truncated, H.q_max_source);
\end{ocexcerpt}
\octranscript{001_dimensions}
So $\dim H^\bullet(\Gamma;\Sym^{10}\F_{691}^2)=(0,3,0)$: weight $k=12$, where $H^1$ is three-dimensional
(the Eisenstein class and the two-dimensional cuspidal part of $\Delta$).

\section{Local systems away from the modular curve}\label{ex:presets}

The presets put the equivariant machinery on $S^1$, the cylinder and the M\"obius band, where the answers
are theorems of topology:

\excerptfrom{013_presets.oc}
\begin{ocexcerpt}
S1 = s1_scalar(1, ELL_P);
H1 = cohomology(S1);
assert(H1.dims == [1, 1], "H^*(S^1; F_5) with trivial monodromy");
print("s1_scalar(1)", H1.dims);
...
CY = cohomology(cylinder_scalar(1, ELL_P));
print("cylinder_scalar(1)", CY.dims);
\end{ocexcerpt}
\octranscript{013_presets}
Reading it: on $S^1$ the scalar system with monodromy $1$ has the cohomology $(1,1)$ of the circle, monodromy $2$ (a
non-trivial scalar over $\F_5$) has none, and the rank-2 unipotent system keeps one invariant and one coinvariant line;
the cylinder and the M\"obius band retract onto a circle, hence the same answers with a vanishing $H^2$. \file{055_s24_language_tour.oc}
(Section~\ref{sec:tour}) adds the simplicial complexes themselves (\vb{sphere}, \vb{moebius}, \vb{projective\_plane},
\vb{complex}).

\section{Coefficients \texorpdfstring{$\Z/\ell^k$}{Z/ell^k}}\label{sec:ringcoh}

Under \texttt{set coefficient\_ring = "Z/ell\^{}2"} (the token \texttt{ell} is literal: the verb's own $\ell$) the same
complex is taken over $\Z/\ell^k$, and $H^q$ becomes a finite $\Z/\ell^k$-module, $H^q\cong\bigoplus_i\Z/\ell^{e_i}$
(\Ecite{sec:ring}). The \vb{Cohomology} handle then reports the \emph{lengths} $\sum_i e_i$ (\texttt{.lengths}), the
exponents (\texttt{.invariant\_factors(q)}), the reduction map to $\Fl$ (\texttt{.red(q)}), the Bockstein rank and, when the
precision suffices, the $\Z_{(\ell)}$-type of the integral lattice (\texttt{.ladic\_type(q)}, certified when
\texttt{.ladic\_type\_certified} is true). The computation is an echelon form over the local ring $\Z/\ell^k$: a Smith
form where the invariant factors are wanted, a Howell form (the canonical row echelon form of a submodule over
$\Z/\ell^k$) where a span must be compared. The script \file{024_e8_ring_points.oc} prints, per point, the lengths of
$H^0,H^1,H^2$, whether the $\ell$-adic type is certified, the Bockstein rank and the higher-torsion count:
\octranscript{024_e8_ring_points}
At $(30,5)$ the precision $\Z/5^2$ is not enough to certify the type (\texttt{false}); $\Z/5^3$ is. Hecke operators over
$\Z/\ell^k$ are in \file{025_e8_console_surface.oc} and \file{026_c5_leftovers_ring_members.oc}, $\ell^k$ beyond $2^{32}$ in
\file{028_wide_band.oc}, and the involution $\iota^*$ over $\Z/\ell^k$ in \file{049_lv3_iota_ring.oc}.

\section{A second model at \texorpdfstring{$\ell=2,3$}{l=2,3}: Moerdijk--Pronk}\label{sec:mp}

A second, independent model of the orbifold cohomology of $\SL_2(\Z)$ at the primes $\ell=2,3$ that divide the orders of the finite subgroups of $\SL_2(\Z)$: the
Moerdijk--Pronk (MP) simplicial model of {\ttfamily[C,~sec:mpmodel]}. It is reached \emph{only through validations}; there is no verb,
setting or record for it. Ten \texttt{OrbifoldMP\dots} names, all gating a C++ computation against a frozen
independent oracle (a Python script and its recorded output under \file{docs/implementation/slot22/mp/}):
\begin{center}\small
\begin{tabular}{@{}>{\raggedright\arraybackslash}p{0.40\linewidth}l>{\raggedright\arraybackslash}p{0.44\linewidth}@{}}
\toprule
\textbf{Validation} & \textbf{Tier} & \textbf{Checks}\\\midrule
\vb{OrbifoldMPRulesValidation} & fast & the triangulation descriptor, MP rules and the MP set\\
\vb{OrbifoldMPTeardropValidation} & fast & MP's teardrop, 9 cells\\
\vb{OrbifoldMPNerveValidation} & standard & nerve cochains, $H^{\le2}$ over $\F_\ell$ and $\Z/\ell^k$\\
\vb{OrbifoldMPComparisonValidation}, \vb{OrbifoldMPRefinementValidation} & standard & comparison and refinement maps\\
\vb{OrbifoldMPAllDegreesValidation} & standard & the all-degrees certificate\\
\vb{OrbifoldMPBoundaryValidation}, \vb{OrbifoldMPLevelConeValidation} & standard & boundary, restriction/corestriction, level pairs\\
\vb{OrbifoldMPHeckeValidation} & standard & Hecke $T_p$ on the MP model\\
\vb{OrbifoldMPNerveH3Slow} & slow & nerve $H^3$ at $n\in\{0,2\}$\\
\bottomrule
\end{tabular}
\end{center}
Run one by name, e.g.\ \cli{./hacs OrbifoldMPNerveValidation}. Each prints informational lines
(\texttt{MP\_RULES}, \texttt{NERVE}, \texttt{INTEGRAL}, \texttt{ORACLE <file> cells=<n> agree=<n>}, \texttt{NC \dots}
for negative controls that must fire) and ends with one summary line
\begin{shell}
ORBIFOLD_MP <Name> checks=<n> failing=0 wall_s=<t>
\end{shell}
Read \texttt{failing=0} and the \texttt{ORACLE \dots agree=} counts, not only the exit status.
\vb{OrbifoldMPAllDegreesValidation} additionally prints, per model (\texttt{M7}, \texttt{R\_Z}), a certificate line
\begin{shell}
MP_ALL_DEGREES model=<M7|R_Z> ring=F_ell nerve_to=2 periodic_on=C_I from=2 via=thm:mpcomparison certified=4/4 scope=bulk verdict=OK
\end{shell}
Its scope is $\Gamma=\SL_2(\Z)$, the bulk only: above degree 2 the cohomology is \emph{transported} by the
periodicity theorem, not computed. No MP validation writes a golden line, a record or a setting. The two
hooks are \env{ORBIFOLD_MP_CELLS} (Section~\ref{sec:env-hooks}) and the scratch emitter
\env{ORBIFOLD_MP_DUMP} (Section~\ref{sec:env-gated}).

%======================================================================
\chapter{Hecke operators}\label{ch:hecke}
%======================================================================

\section{The operators and their certificates}\label{sec:heckeops}

For a prime $p$ the Hecke operator $T_p$ is the double coset $\Gamma\,\mathrm{diag}(1,p)\,\Gamma$ acting on cohomology
({\ttfamily[C,~def:hecke]}); at $p=\ell$ the same code path gives $U_\ell$, for $\ell\le23170$ in the narrow tier
({\ttfamily[C,~lem:bounds]}). The family adds $T_{p^r}$ by the algebra recursion and $T(p,p)$
({\ttfamily[C,~thm:relations]}, {\ttfamily[C,~prop:diamond]}). The program computes an operator in three layers, which
reappear as the report levels (a)/(b)/(c) of a campaign: (a) the cochain-level matrix $\Theta^{(q)}$, (b) its matrix
$A^{(q)}(T_p)$ on $H^q$, (c) the primary decomposition of $H^q$ over a commuting family. Each layer is certified in
process: $\partial\tau=\tau\partial$ for the cochain map, two independent routes agree entry by entry, and a decomposition
block is certified \emph{local} (it has exactly one maximal ideal) by the locality certificate (setting \vb{locality}).
The verbs are \vb{hecke(H, p)} (one operator), \vb{family(H)} (the commuting family $T_2,T_3,\dots$), \vb{decompose(T, q)}
and \vb{report(D)}; relations between operators are checked in \file{004_family_identities.oc}, certificates in
\file{011_decompose_certify.oc}, the record schema in \file{012_report_schema.oc}, the strict locality mode in
\file{015_strict_locality.oc}, and the golden pins of $A(T_p)$ in \file{002_hecke_pins.oc} and
\file{003_bases_and_projectors.oc}. The grid scripts \file{007_grid_classical.oc} to \file{010_grid_edge_n.oc} walk the four
regimes of the golden grid (classical, Dickson-exceptional, $\ell\in\{2,3\}$, edge cases in $n$).

\section{The flagship: \texorpdfstring{$T_2$}{T2} on \texorpdfstring{$(10,691)$}{(10,691)}}\label{ex:flagship}

\excerptfrom{006_flagship.oc}
\begin{ocexcerpt}
T2 = hecke(H, P_HECKE);
A = T2.A(1);
print(charpoly(A));
D = decompose(T2, 1);
print(D.num_blocks);
for b in D.blocks {
print(b.dim, b.residue_degree, b.min_poly(2), b.jordan(2));
}
...
print(D.locality.complete, D.locality.splits, D.locality.phase_reached);
R = report(D);
print(R.q_max, R.truncated, R.q_max_source, R.degrees);
\end{ocexcerpt}
\octranscript{006_flagship}
Reading it: the characteristic polynomial of $T_2$ on $H^1$ is printed over $\F_{691}$; the decomposition
of $H^1$ over $\F_{691}[T_2]$ has one primary block of dimension 3 and residue degree 1, on which $T_2$ has
minimal polynomial $(X+24)^2$ and Jordan type $(2,1)$ -- the Ramanujan congruence
$\tau(2)\equiv 1+2^{11}\pmod{691}$ seen as a non-semisimple Hecke module. The block is certified local by
the linear phase without a split, and \vb{report(D)} produces the \texttt{hecke\_report/v2} record, which covers every degree
(\texttt{R.degrees} prints \texttt{[0, 1, 2]}) and carries the degree-1 decomposition.

\section{Linear algebra on Hecke operators}\label{ex:linalg}

\excerptfrom{005_linear_algebra.oc}
\begin{ocexcerpt}
H = cohomology(symmetric_power(MODEL, N_SYM, ELL));
F = family(H);
T2 = F.T(2)[1];
T3 = F.T(3)[1];
d = H.dims[1];
A = algebra([T2, T3]);
assert(in_span(A.basis, mul(T2, T3)), "the algebra is closed under multiplication");
C = commutant([T2, T3]);
...
N = matrix(7, [[0, 1, 0], [0, 0, 1], [0, 0, 0]]);
assert(jordan_type(N) == [3], "one Jordan block of size 3");
expect_error "not nilpotent" { jordan_type(identity(7, 2)); }
...
print(d, A.dim);
\end{ocexcerpt}
\octranscript{005_linear_algebra}
The algebra $\F_{691}[T_2,T_3]$ acting on the three-dimensional $H^1$ is two-dimensional, consistent with
the single block of Jordan type $(2,1)$ above.

%======================================================================
\chapter{Integral cohomology: the descent over \texorpdfstring{$\Z$}{Z}}\label{ch:integral}
%======================================================================

Over a field the answer is a list of dimensions. Over $\Z$ it is a list of finitely generated abelian groups, and the
interesting information is in their torsion: at $2$ and $3$, where $\Gamma$ has finite subgroups, and at finitely many
``bad'' primes $\ell\ge5$ tied to congruences. This chapter explains what the program computes over $\Z$, how it does it,
how to read the record, and what it does not (yet) do. The mathematics is [C]'s Part on the integral line
({\ttfamily[C,~sec:integral]}), the descent ({\ttfamily[C,~sec:descent]}, {\ttfamily[C,~sec:descentN]}), its Hecke operators
({\ttfamily[C,~sec:descenthecke]}) and the boundary pair ({\ttfamily[C,~sec:descentpair]}); the one statement that collects
what a record certifies is {\ttfamily[C,~thm:descentcert]}, and for the boundary pair {\ttfamily[C,~prop:descentpair]}
(the pair is outside {\ttfamily[C,~thm:descentcert]}(7)(c), a separate statement).

\section{What is computed}\label{sec:int-what}

For a level $N$ and $n\ge0$ let $A=\Gamma_0(N)$ and $V=\Sym^n(\Z^2)$. The group cohomology $H^q(A;V)$ is a finitely
generated abelian group, so by the structure theorem
\[
  H^q(A;V)\;\cong\;\Z^{b_q}\;\oplus\;\bigoplus_{\ell}\tors\,H^q(A;V)_{(\ell)},\qquad
  \tors\,H^q(A;V)_{(\ell)}\cong\bigoplus_i\Z/\ell^{a_i},
\]
a free part and, for each prime, an $\ell$-primary torsion part. Four facts make this computable
({\ttfamily[C,~thm:descentcert]}(1), (6)):
\begin{enumerate}[nosep]
\item \textbf{Ranks.} $b_q$ is the dimension over any \emph{good} prime field $\F_{\ell'}$, $\ell'\ge5$ with no
  $\ell'$-torsion; for $q\ge2$ it is $0$ ($\Gamma_0(N)$ has virtual cohomological dimension 1).
\item \textbf{Which primes.} The torsion lives at $2$, $3$ and at the \emph{bad primes} $\mathcal B$, a closed-form finite
  list: at level~1 and even $n>0$ these are the primes $5\le\ell\le n-1$ with
  $\#\{(a,b)\in\Z_{\ge0}^2:a(\ell+1)+b(\ell^2-\ell)=n\}>0$, and that count is the number of cyclic summands
  ({\ttfamily[C,~prop:badprimes]}; at level $N$, {\ttfamily[C,~prop:badprimesN]}). At $\ell\ge5$ the torsion sits in $H^1$
  only. For odd $n$, $2$ kills everything.
\item \textbf{Which degrees.} For $q\ge2$, $24$ kills $H^q$, and the cup product with an Euler class is an isomorphism
  $H^q\cong H^{q+2}$ of Hecke modules: the degrees $0\le q\le3$ determine all of them ($q_{\max}=3$, \texttt{truncated =
  true}, \texttt{q\_max\_source = derived}).
\item \textbf{No glueing problem.} The primes do not interact: the arithmetic fracture square
  $\Z=\Z[1/6]\times_{\Q_2\times\Q_3}(\Z_2\times\Z_3)$ gives, degree by degree, $H^q$ as a fibre product of its
  localisations, with vanishing connecting maps ({\ttfamily[C,~prop:fracture]}, {\ttfamily[C,~prop:cohglue]},
  {\ttfamily[C,~prop:descent]}). So the group is assembled from a rank and one $\ell$-primary part per prime, each read at
  that prime alone.
\end{enumerate}
On top of the groups the record carries the Hecke operators $T_p$, $p\nmid N$: the characteristic polynomial over $\Z$ of
$T_p$ on the free quotient $H^1/\tors$, and on each torsion group the isomorphism class of $(\tors H^q_{(\ell)}, T_p)$ as a
module, given by invariants that do not depend on any choice; and the boundary pair (bulk, boundary, compact support).

\section{Switching it on}\label{sec:int-switch}

The descent is the one consumer of \texttt{coefficient\_ring = Z}. In the console:
\begin{shell}
set coefficient_ring = "Z"
D = descent(N, n)            # or descent(N, n, [p1, p2, ...]): the Hecke primes
\end{shell}
or from the command line, the modifier after the console form:
\runoutput{descent_cli}
Under \texttt{Z} every other pipeline verb refuses, by name and with exit code~5, rather than compute over another ring:
\runoutput{descent_refused}
The verb's limits keep a console call interactive: $N\le300$, $n\le60$ and $\psi(N)\,(n+1)\le300$ (at the bound about
40\,s and 1.4\,GiB); the Hecke primes are 1 to 6 primes $\le97$ not dividing $N$, sum $\le100$ (default: the first three of
$2,3,5,7,11,13$ not dividing $N$). A call outside the limits is a type error (exit~3) naming the limit. Larger cells go
through a campaign (Section~\ref{sec:campaign-z}); their cost is discussed in Section~\ref{sec:env-descent}.

\section{The algorithm, step by step}\label{sec:int-algorithm}

The record is assembled from independent pieces, one per prime, each \emph{certified}: a piece either proves that what it
read is the true integral answer, or it is refused. Nothing is guessed, and the code uses no arbitrary-precision integers:
every computation is in machine words modulo a prime power.

\paragraph{1. The free rank, at a good prime.}
At a prime $\ell'\ge5$ with no $\ell'$-torsion (by default $\ell'=10007$) the model computes group cohomology, and
$b_q=\dim H^q(A;\Sym^n\F_{\ell'}^2)$ ({\ttfamily[C,~prop:certifiede]}). At level $N$ the program does not build the induced
module on the model; it uses the \emph{cokernel presentation} of $H^1$: generators indexed by the cosets, relations from the
two finite subgroups $\langle S\rangle$, $\langle U\rangle$, a sparse matrix over $\F_{\ell'}$ (route \texttt{sparse}).

\paragraph{2. Torsion at a bad prime \texorpdfstring{$\ell\ge5$}{l>=5}.}
The same presentation is reduced modulo $\ell^K$ and brought to \emph{Smith normal form}: invertible row and column
operations over $\Z/\ell^K$ make it diagonal, with pivots $u_i\ell^{\sigma_i}$, $u_i$ a unit, $\sigma_i<K$. The cokernel is
then visibly $\bigoplus\Z/\ell^{\sigma_i}$ plus a free part. The read is \emph{certified} when the number $s$ of pivots
equals the rank of the matrix over $\Q$ (known from step~1): then no pivot of the true integral matrix is hidden in
$\ell^K$, and $\tors H^1_{(\ell)}=\bigoplus_{\sigma_i>0}\Z/\ell^{\sigma_i}$ exactly ({\ttfamily[C,~prop:certifiedeell]}(6),
{\ttfamily[C,~prop:descentN]}). $K$ is chosen from an a-priori bound and recorded.

\paragraph{3. Torsion at \texorpdfstring{$2$}{2} and \texorpdfstring{$3$}{3}: a free model.}
Here the finite stabilisers matter and the seven-vertex model computes the wrong object, so the program uses an
\emph{$\ell$-local free model} of $\Gamma$: a complex of free $\Z_{(\ell)}[\Gamma]$-modules whose cohomology is group
cohomology. By default it is the \emph{tree model} ({\ttfamily[C,~prop:isomodeltree]}): $\SL_2(\Z)$ is the amalgam
$\Z/4*_{\Z/2}\Z/6$ of the stabilisers of its tree, so the cochains are the cone of
$P^{\langle-I\rangle}\to P^{\langle S\rangle}\oplus P^{\langle U\rangle}$ in closed form, periodic in the higher degrees
(route \texttt{tree}). Its differentials are brought to Smith form over $\Z/2^{30}$ (re-run in 61-bit words if not certified)
and over $\Z/3^{19}$, on two dedicated ring engines, \texttt{ring2k} and \texttt{ring3k} (\texttt{ring2k} also computes Howell forms, the
canonical echelon form of a submodule over $\Z/2^k$);
the certification rule is the one of step~2 ({\ttfamily[C,~prop:certifiedeell]}), and at~3 an uncertified differential is
refused. When $\Gamma_0(N)$ has no element of order~3 a shortcut applies at~3 ({\ttfamily[C,~prop:descentNthree]}).

\paragraph{4. Hecke operators on the free part: characteristic polynomials by residues.}
$\chi_p=\det(x-T_p\mid H^1/\tors)\in\Z[x]$ is read modulo several good primes, where it is the characteristic polynomial of
the shipped Hecke matrix, and recombined by the Chinese remainder theorem against a certified bound on its coefficients
({\ttfamily[C,~prop:crtcharpoly]}, {\ttfamily[C,~prop:crtcharpolyN]}). A coefficient that the residues do not determine is
reported as \texttt{"?"}, never guessed.

\paragraph{5. Hecke operators on torsion: an exact read.}
On a certified Smith form the generators of the torsion are known modulo $\ell^{K-\sigma_{\max}}$; applying $T_p$ to them
gives the module $(\tors H^q_{(\ell)},T_p)$ exactly provided $2\sigma_{\max}\le K$, and that condition is necessary
({\ttfamily[C,~prop:torsionhecke]}(3)). With a weaker precision only the layer polynomials and the Fitting types are
certified (\texttt{module\_certified = false}). The module is then described by \emph{choice-free invariants}
(Section~\ref{sec:int-reading}), never by a matrix, because a matrix depends on the generators chosen
({\ttfamily[C,~prop:torsionhecke]}(1), (2)).

\paragraph{6. The boundary pair.}
The Borel--Serre compactification adds a boundary (one circle per cusp). The record computes, by the same descent, the
cohomology of the bulk $H^q$, of the boundary $H^q(\partial)$ and with compact support $H^q_c$ (the cone of
restriction to the boundary), $q\le4$, glued in the long exact sequence
$\cdots\to H^q_c\to H^q\to H^q(\partial)\to H^{q+1}_c\to\cdots$ ({\ttfamily[C,~prop:descentpair]}).

\paragraph{Status at \texorpdfstring{$\ell=2,3$}{l=2,3}, and what is not computed.}
At $2$ and $3$ the groups, the Hecke torsion modules and the pair are computed for every level within the limits, and each
piece is either certified or refused, never wrong (at $3$ the precision cap is $K=19$ and an uncertified differential is
refused, Section~\ref{sec:int-algorithm}, step~3; a memory prediction can refuse too, Section~\ref{sec:int-reading}); at
$\ell\ge5$ likewise at every bad prime. \emph{Not} in the record ({\ttfamily[C,~thm:descentcert]}(7)): $U_p$ for
$p\mid N$ (deferred: a prime dividing $N$ is refused), a nebentypus, groups other than $\Gamma_0(N)$, integer Hecke matrices
in a basis, the Hecke-module structure of $H^1/\tors$ beyond its characteristic polynomial, and the extension class of
$0\to\tors\to H^1\to H^1/\tors\to0$.

\paragraph{The second route at level 1: the orbifold census.} The descent covers every level $N$ but only $n\le60$. For
level~1 and large $n$ at $\ell\in\{2,3\}$ a second route computes the same group cohomology of $\SL_2(\Z)$: the orbifold
census (Section~\ref{sec:orbifold}), on the isotropy model $C_I$, cross-checked against the amalgam resolution
$\Z/4*_{\Z/2}\Z/6$. Where both run they agree: \texttt{IntegralDescentGroupsValidation} compares the descent at level~1 with
an independent $C_I$ computation up to degree~4.

\section{Reading the record}\label{sec:int-reading}

\vb{descent(N, n, primes)} returns a lazy \vb{DescentRecord} (schema \texttt{integral\_descent\_record/v1},
{\ttfamily[C,~sec:recordschema]} item~7; members in Appendix~\ref{app:types}).

\paragraph{Groups.} \texttt{.groups} lists $H^0,\dots,H^3$ as strings: \texttt{"0"} for the trivial group, otherwise
\texttt{Z} (rank~1) or \texttt{Z\^{}r} ($r\ge2$) for the free part followed by the cyclic summands $\Z/\ell^a$ (written \texttt{Z/l\^{}a}, \texttt{Z/l} for
$a=1$), joined by \texttt{+}, ordered by prime and then by exponent; a summand that occurs several times is repeated.
\texttt{.degrees} is \texttt{[0, 1, 2, 3]}; the groups are empty when the record is refused. The same groups in machine form are \texttt{.groups\_inv}
(and \texttt{.pair\_inv} for the pair): one \texttt{\{free\_rank, torsion\}} per degree, \texttt{torsion} a list of
\texttt{[p, k, mult]} meaning \texttt{mult} summands $\Z/p^k$; the notebook renders these, scripts may prefer them to parsing
the strings.

\paragraph{Primes and precision.} \texttt{.bad\_primes} is $\mathcal B$ as a list of \texttt{[ell, count]} (the count of
cyclic summands at $\ell$); \texttt{.torsion\_primes} the primes whose torsion was read ($2$, $3$ for even $n$, and
$\mathcal B$); \texttt{.K} one entry per prime, \texttt{\{ell, route, K, W, certified\}}: the route (\texttt{tree} at $2$
and $3$; \texttt{shortcut} at $3$ when $\nu_3(N)=0$, i.e.\ when $\Gamma_0(N)$ has no element of order~3 --- the default
there, so the worked level~11 reads \texttt{shortcut} ({\ttfamily[C,~prop:descentNthree]}, Section~\ref{sec:int-algorithm},
step~3); \texttt{sparse} at $\ell\ge5$; \texttt{dense} and \texttt{bredon} only under the rollback arms of
Section~\ref{sec:env-descent}), the precision $K$ of $\Z/\ell^K$, the word width $W$ in bits, whether the read certified.
\texttt{.K} is provenance: it is not part of the record's hash.

\paragraph{Characteristic polynomials.} \texttt{.charpoly} has one entry per Hecke prime: \texttt{\{p, r, coefficients, exact,
determined, digest, refused\}}, $r=\operatorname{rank}H^1/\tors$, and the coefficients of $\chi_p$ \emph{from the leading one
down}, each a string: a decimal integer, \texttt{"[sign,D digits]"} for an integer beyond $2^{63}$ (its sign, and $D$ =
the number of mixed-radix digits of its CRT reconstruction, one per word prime, each $<2^{32}$ --- \emph{not} the number of
decimal digits), or \texttt{"?"} when not determined. \texttt{exact} (since 1.5.4) is the same list with \emph{every} coefficient
as an exact decimal integer (\texttt{"?"} when not determined): the record carries each coefficient beyond $2^{63}$ exactly
({\ttfamily[C,~sec:recordschema]} item~7~(l); in a campaign cell file the entry's list \texttt{big} of canonical signed hexadecimal
integers, decoded by \file{docs/integral\_census/icensus\_exact.py}). (The verbs of Section~\ref{sec:levels} print coefficient lists the other way
round, from the constant term up.) \texttt{.charpoly\_factored} gives the same $\chi_p$ factored over $\Z$: per prime
\texttt{\{p, status, expanded, factored, factors, complete, verified, note\}}, e.g.\ \texttt{factored = "(X + 24)\^{}2 (X - 2049)"}
at $(N,n)=(1,10)$, $p=2$, and \texttt{factors} a list of \texttt{\{coefficients, degree, multiplicity, irreducible\}}. The
factorisation is computed in the program (Yun's squarefree decomposition, then Zassenhaus: Hensel lifting and recombination),
exactly, and the product of the factors is checked against $\chi_p$; a coefficient beyond $2^{63}$ is read from the record's exact
value. The factorisation is presentation only: the record's JSON, digest and golden points do not contain it.

\paragraph{The torsion modules.} \texttt{.torsion} has one entry per $(q,\ell)$ with $q\in\{1,2,3\}$: \texttt{\{q, ell,
exponents, module\_certified, ops, refused\}}, where \texttt{exponents} are the $a_i$ and \texttt{ops} has one
\texttt{\{p, invariants\}} per Hecke prime. The \texttt{invariants} string of $(G,\phi)=(\tors H^q_{(\ell)},T_p)$ has up to six
fields separated by \texttt{;}, each defined in {\ttfamily[C,~prop:torsionhecke]}(1). The fields \texttt{ker=} and
\texttt{cok=} are \emph{present only} when \texttt{module\_certified} is true, $G$ has at most $400$ cyclic summands and the
ring's precision cap is at least $a_{\max}+2$; otherwise the string is \texttt{a=;P=;B=;fit=} (four fields). A
\texttt{B=} block of more than $300$ summands reads \texttt{$j$:$e$:-} (not computed, never guessed).
\begin{center}\small
\begin{tabular}{@{}l>{\raggedright\arraybackslash}p{0.78\linewidth}@{}}
\toprule
\textbf{Field} & \textbf{Meaning}\\\midrule
\texttt{a=}$a_1.a_2.\cdots$ & the type of $G$: $G\cong\bigoplus_i\Z/\ell^{a_i}$, exponents ascending\\
\texttt{P=}$j$\texttt{:}$c_0,c_1,\dots$\texttt{|}\dots & per layer $j$ (the summands of order $\ell^j$): the characteristic polynomial $P_j$ of $\phi$ on the layer $V_j=(\ell^{j-1}G)[\ell]/(\ell^jG)[\ell]$, over $\Fl$, coefficients from the leading one down ((a))\\
\texttt{B=}$j$\texttt{:}$e$\texttt{:}$c_0,\dots$\texttt{|}\dots & per layer $j$: the characteristic polynomial of the layer's diagonal block modulo $\ell^{e}$, $e=e_j=\min(j,d_j)$ the precision at which it is an invariant ($d_j$ the distance to the nearest other exponent); its reduction mod $\ell$ is $P_j$ ((e))\\
\texttt{ker=}$c$\texttt{:}type\texttt{|}\dots & for every $c=0,\dots,\ell^{a_{\max}}-1$: the type of $\ker(\phi-c)$ ((c))\\
\texttt{cok=}$c$\texttt{:}type\texttt{|}\dots & the same for $\operatorname{coker}(\phi-c)$\\
\texttt{fit=}$c$\texttt{:}type\texttt{|}\dots & for every $c=0,\dots,\ell-1$: the type of the Fitting component $G_{[c]}=\ker(\phi-c)^{\mathrm{len}\,G}$, the part on which $\phi\equiv c$ is nilpotent ((b))\\
\bottomrule
\end{tabular}
\end{center}
A type is written like \texttt{a}: exponents joined by \texttt{.}, and \texttt{-} for the trivial group. When the module is
not certified (Section~\ref{sec:int-algorithm}, step~5) the read certifies only the layer polynomials (\texttt{P=}), the Fitting types (\texttt{fit=}) and the blocks
modulo $\ell^{\min(e_j,K')}$, $K'=K-\sigma_{\max}$ ({\ttfamily[C,~prop:torsionhecke]}(3)). The same text is produced by the independent Python oracle
\file{docs/implementation/slot24/d3_hecke_oracle.py}, which is how it is tested.

\paragraph{The pair, status, identity.} \texttt{.pair} holds \texttt{bulk}, \texttt{boundary} and \texttt{compact}, each with
its own \texttt{groups} for $q=0,\dots,4$ (\texttt{compact} also \texttt{derived\_from = [bulk, boundary]}).
\texttt{.status} is \texttt{ok} or \texttt{refused}, \texttt{.refusal} the first refused part's reason (a precision cap or a
memory prediction: refused, never wrong); \texttt{.digest} is \texttt{record\_fnv1a64}, the FNV-1a-64 of the record's
invariants and the settings' canonical text; \texttt{.json} the record as one JSON line. Printing the record itself gives a
one-line summary.

\section{Worked examples}\label{sec:int-examples}

The script \file{054_s24_descent.oc} reads three cells and four refusals:

\excerptfrom{054_s24_descent.oc}
\begin{ocexcerpt}
set coefficient_ring = "Z";
D = descent(1, 12);
print("descent(1, 12):", D.status);
for q in D.degrees {
print("  H^" + str(q) + " =", D.groups[q]);
}
print("  bad primes [ell, count]:", D.bad_primes, " torsion primes:", D.torsion_primes);
...
t = D.torsion[0];
print("  tors H^" + str(t.q) + " at ell =", t.ell, ": exponents", t.exponents, " module certified:", t.module_certified);
print("  T_" + str(t.ops[0].p) + ":", t.ops[0].invariants);
...
E = descent(11, 0);
print("descent(11, 0):", E.groups);
print("  bulk    :", E.pair.bulk.groups);
print("  boundary:", E.pair.boundary.groups);
print("  compact :", E.pair.compact.groups);
...
F = descent(11, 2, [2, 3]);
...
expect_error "n must satisfy 0 <= n <= 60" { descent(1, 400) }
expect_error "must not divide N" { descent(4, 2, [2]) }
\end{ocexcerpt}
\octranscript{054_s24_descent}

\paragraph{The Hecke structure over $\Z$ (\texttt{hecke\_structure}).} The verb \texttt{hecke\_structure(N, n, algebra = "full")} (slot~25,
HZ-REC) returns the record \texttt{integral\_hecke\_structure/v1} of a cell ({\ttfamily[C,~sec:recordschema]} item~8): per prime $\ell$ whether
$\mathbb T\otimes\Z_\ell$ is certified to be generated by the Hecke operators (and by which route), and the congruence module $C_{\mathbb T}$
between the Eisenstein and the cuspidal part, \texttt{certified} only where its lower and upper bounds meet. Like \texttt{descent} it needs
\texttt{coefficient\_ring = "Z"}. The script \file{060_s25_hecke_structure.oc} reads $(1,10)$ ($C_{\mathbb T}=\Z/691$), $(1,22)$ and the
anemic algebra at $(11,0)$; a campaign writes the same records with \texttt{structure = hecke} (keys \texttt{algebra}, \texttt{hz\_cert}).

\paragraph{Level 1, \texorpdfstring{$n=12$}{n=12} (weight 14).} $S_{14}(\SL_2(\Z))=0$, so $H^1/\tors=\Z$ is spanned by the
Eisenstein class, on which $T_p$ acts by $1+p^{13}$: $\chi_2=x-8193$, $\chi_3=x-1594324$, $\chi_5=x-1220703126$. The bad
primes follow from the count of Section~\ref{sec:int-what}: at $\ell=5$, $6a+20b=12$ has the one solution $(2,0)$; at
$\ell=11$, $12a+110b=12$ has $(1,0)$; at $\ell=7$, $8a+42b=12$ has none. Hence one summand $\Z/5$ and one $\Z/11$ in $H^1$: the count is the dimension of the $\Gamma$-invariants of
$\Sym^{12}\Fl^2$ (Dickson invariants of degree~12), and each invariant line becomes a cyclic summand of $H^1$ through the
Bockstein ({\ttfamily[C,~prop:badprimes]}). The $2$- and $3$-parts come from the finite subgroups, and $H^2$, $H^3$ are killed by $24$. The \texttt{K} lines show which
route read each prime at which precision.

\paragraph{Reading an invariants string.} For $T_2$ on $G=\tors H^1_{(2)}\cong\Z/2\oplus\Z/4\oplus\Z/8$
(\texttt{a=1.2.3}): every layer is one-dimensional and \texttt{P=1:1,0|2:1,0|3:1,0} says $P_j=x$ on each, so $T_2$ is
nilpotent modulo~2 on every layer; \texttt{B} gives the same at precision $e_j=1$ (the exponents $1,2,3$ are at distance~1).
\texttt{ker=0:1.3|1:-|2:1.3|\dots} says $\ker(T_2-c)\cong\Z/2\oplus\Z/8$ for even $c$ and $0$ for odd $c$, and
\texttt{cok=} the cokernels $\Z/2\oplus\Z/2\oplus\Z/4$ (the orders agree, $2^4$, as they must for an endomorphism of a finite
group). \texttt{fit=0:1.2.3|1:-}: the whole group is the Fitting component at $c\equiv0$, i.e.\ $T_2$ is topologically
nilpotent on the 2-torsion, and $T_2-1$ is invertible there.

\paragraph{Level 11, \texorpdfstring{$n=0$}{n=0} (weight 2).} $H^0=\Z$, $H^1=\Z^3$: the two classes of the elliptic curve
$X_0(11)$ and one Eisenstein class. $\chi_2=x^3+x^2-8x-12=(x+2)^2(x-3)$: the eigenvalue $a_2=-2$ of the curve (twice) and
$1+2=3$ on the Eisenstein class; likewise $\chi_5=(x-1)^2(x-6)$. The pair shows the two cusps: $H^0(\partial)=H^1(\partial)=\Z^2$,
and in the long exact sequence the restriction $H^1\to H^1(\partial)=\Z^2$ has rank~1 (the Eisenstein class): its kernel is
the rank-2 cuspidal part, $H^1_c=\Z^3$ is that kernel plus the image of $H^0(\partial)/H^0=\Z$, and $H^2_c=\Z$ is the
fundamental class.

\paragraph{Level 11, \texorpdfstring{$n=2$}{n=2} (weight 4).} $\operatorname{rank}H^1=6=2\dim S_4(\Gamma_0(11))+2$ (two
cusp forms, two cusps); the characteristic polynomials have degree~6, their integer coefficients reconstructed from residues by the CRT step of
Section~\ref{sec:int-algorithm}. The only torsion of $H^1$ in this cell is one $\Z/2$.

\paragraph{Refusals.} The limits and $p\mid N$ are type errors; outside \texttt{coefficient\_ring = Z} the verb asks for the
ring; and under \texttt{Z} the field verbs refuse with exit code~5 (here caught by \vb{expect\_error}).

\section{From the command line and in campaigns}\label{sec:int-campaigns}

A grid of descent cells, with one JSON record per cell and a summary table, is a campaign over $\Z$
(Section~\ref{sec:campaign-z}). Its records are the same \texttt{integral\_descent\_record/v1} as the verb's. The golden reference
pins nine descent cells (156 tier-A entries), so a change to a record is a re-baseline.

%======================================================================
\chapter{Extensions, Hida theory and the forms algebra}\label{ch:ext}
%======================================================================

The extension layer of [E] studies $H^1$ over $\Fl$ as a module over its Hecke algebra $\mathbb T$, not just its
eigenvalues: congruences between an Eisenstein series and a cusp form, or between cusp forms, show up as \emph{non-split}
extensions, i.e.\ Jordan blocks of the Hecke operators. The verbs return records; the frozen scripts below show each one.

\section{The extension layer}\label{sec:extlayer}

The chain of objects is: the commuting family $T_2,T_3,\dots$ (with $U_\ell$ when the memory model admits it, setting
\vb{u\_ell}); the \emph{maximal ideals} $\mathfrak m$ of the algebra they generate (the mod-$\ell$ eigensystems, in a
canonical splitting field \vb{GF}); for each $\mathfrak m$ the \emph{local Hecke algebra} $\mathbb T_\mathfrak m$ and the
localised module $H^1_\mathfrak m$ with its Loewy and socle series, its presentation and its extension classes; and the
Sturm bound $B'$ that proves a finite family generates $\mathbb T$ (\Ecite{def:bprime}). \vb{ext\_record(n, ell)} assembles
all of it into one record per cell, optionally with the cone and boundary cohomology and the $\iota$-sign split ($\iota^*$
the involution of {\ttfamily[C,~def:iota]}; the plus part contains the Eisenstein side). The local-algebra invariants
(\vb{radical}, \vb{socle}, \vb{loewy}, \vb{tangent\_space}, \vb{nilpotency\_index}, \vb{is\_local}, \dots) apply to any
commutative algebra of operators (\file{047_c5_local_algebra_verbs.oc}).

\section{Sturm bounds and the \texorpdfstring{$U_\ell$}{U_ell} budget}\label{ex:ext}

\excerptfrom{016_extensions.oc}
\begin{ocexcerpt}
S = sturm_bound(24, 5);
print(S.k, S.ell, S.B_prime, S.dim_M_k, S.closure_dim, S.primes_done, S.cap);
CUT = sturm_bound(24, 5, 3);
print(CUT.outcome, CUT.closure_dim, CUT.dim_M_k, CUT.primes_done);
CONT = resume(CUT, 36);
print(CONT.B_prime, state(CONT).resumable.progress.budget_history);
V = u_ell_decision(16, 43867, 2048);
print(V.included, V.reason, V.predicted_mb, V.budget_mb, V.model_version);
W = u_ell_decision(10, 691, 2048);
print(W.included, W.reason, W.predicted_mb);
set u_ell = "off";
OFF = u_ell_decision(10, 691);
print(OFF.included, OFF.reason);
\end{ocexcerpt}
\octranscript{016_extensions}
The computed Sturm bound at weight 24 and $\ell=5$ is $B'=5$ (\Ecite{def:bprime}): $T_2$ and $T_3$ alone
generate a proper subalgebra, which is why the cap-3 run ends \texttt{none\_at\_cap}. The $U_\ell$ decision
is a frozen memory model: $U_{43867}$ at $n=16$ is refused against a 2048 MB budget, $U_{691}$ at $n=10$
fits, and \texttt{u\_ell = "off"} refuses without forming a prediction.

\section{The extension record of \texorpdfstring{$(10,691)$}{(10,691)}}\label{ex:extrecord}

\excerptfrom{040_c5_ext_record_10_691.oc}
\begin{ocexcerpt}
R = ext_record(N_SYM, ELL, cone = true, sign_split = true);
print(R);
m = R.ideals[0];
print(m);
T = local_algebra(m);
print(T);
P = module_structure(m);
print(P);
X = P.structure_map(0);
print(X);
Q = presentation(m);
print(Q);
plus = R.ideals_plus;
print(plus[0], R.ideals_minus[0]);
\end{ocexcerpt}
\octranscript{040_c5_ext_record_10_691}
One maximal ideal $\mathfrak{m}=(T_2+24)$, certified local, with $\mathbb{T}_\mathfrak{m}=\F_{691}[\varepsilon]/(\varepsilon^2)$;
$H^1_\mathfrak{m}$ has Loewy and socle series $(2,1)$ and is not cyclic; the extension class between the two
layers is non-split; and the $\iota$-sign records put a two-dimensional block in $H^1_+$ (the
Eisenstein-cuspidal congruence) and the cuspidal line in $H^1_-$. That this block is non-split is not printed
for the plus record; it follows because a Jordan block of size~2 cannot lie in the one-dimensional $H^1_-$.

\section{The Dickson sequence, cone, boundary and \texorpdfstring{$\iota$}{iota}}\label{sec:dickson}

At the \emph{Dickson primes} (those $\ell$ for which $\Sym^n\Fl^2$ has $\Gamma$-invariants, the Dickson invariants)
$H^0$ does not vanish, and the Dickson long exact sequence and the map $\theta$ relate $\Sym^n$ to $\Sym^{n-\ell-1}$ (\Ecite{sec:thetaobj}): verbs
\vb{dickson\_les}, \vb{theta}, \vb{dickson\_count}, script \file{027_dickson_theta.oc}. The cone and boundary cohomology
(\vb{cone}, \vb{boundary}) and the involution (\vb{iota}) appear in \file{042_c5_dickson_class_cone.oc},
\file{043_c5_boundary_coincidences.oc}, \file{044_c5_multiplicity_two.oc}, \file{045_c5_pd_loewy_reversal.oc} and
\file{046_c5_identities.oc}; the Kummer cells of [E]'s Stage~9 in \file{041_c5_stage9_kummer_cells.oc}; the canonical
fields in \file{047_c5_local_algebra_verbs.oc}; the Hecke algebra over $\Z/\ell^K$ with its certificate
({\ttfamily[C,~cor:stablecert]}) in \file{029_hecke_ring_algebra.oc}; the predictions in \file{017_predictions.oc}; the
phase scripts \file{018_e0_locality_grid.oc}, \file{019_e1_stage1_sturm.oc}, \file{020_e2_ramanujan_10_691.oc} and the duality
pairing in \file{023_e7_duality_10_691.oc}.

\section{Hida theory, \texorpdfstring{$\Lambda$}{Lambda}-adic objects and the forms algebra}\label{sec:hida}

Hida theory follows the \emph{ordinary} part (where $U_\ell$ is a unit) across weights $k\equiv k'\pmod{(\ell-1)\ell^{m-1}}$
and up the tower $\Gamma_1(\ell^r)$. \vb{hida\_tower(n, ell, m)} reads the ordinary rank and the Hecke polynomials at the
levels $\Z/\ell^e$, $e\le m$; \vb{hida\_lambda(ell, k, W)} the ordinary part of the forms over $\Z/\ell^W$ and its Eisenstein
component; \vb{forms\_algebra(k, ell)} the cusp forms $S_k$ over $\Fl$ by $q$-expansions, an independent check of the
cohomology; \vb{pow} and \vb{kummer\_unit} the matrix powers and the cyclotomic units behind the Kummer criterion
$\ell\mid B_i$; and \vb{lfunction} the $\ell$-adic $L$-function of a branch (\Ecite{sec:lvlfunction}).

\excerptfrom{051_s21_hida_tower.oc}
\begin{ocexcerpt}
T = hida_tower(10, 11, 2);
print("h =", T.h, " e_max =", T.e_max, " K =", T.K, " exps =", T.exps);
for L in T.levels {
print("e =", L.e, " ord_rank =", L.ord_rank, " readings agree:", L.readings_agree, " charpoly T_2 =", L.charpoly["2"]);
\end{ocexcerpt}
\octranscript{051_s21_hida_tower}
The ordinary part of $H^1(\Sym^{10})$ at $\ell=11$ has rank~3 at both levels, and the polynomial of $T_2$ over $\Z/11$ lifts
to one over $\Z/11^2$ (coefficients from the constant term up).

\excerptfrom{053_s22_forms_hida_lambda.oc}
\begin{ocexcerpt}
F = forms_algebra(12, 7);
print("S_12 over F_7: d =", F.d, " basis =", F.basis[0], " charpoly T_2 =", F.ops[0].charpoly);
...
H = hida_lambda(11, 12, 3, [2, 3]);
\end{ocexcerpt}
\octranscript{053_s22_forms_hida_lambda}
$S_{12}$ is spanned by $\Delta=q-24q^2+252q^3-\cdots$, whose reduction mod~7 is the printed basis, and $T_2$ acts by
$\tau(2)=-24\equiv4$ (the polynomial $x+3$). Over $\Z/11^3$ the ordinary cusp part has rank~1 and $T_2$ acts on the
Eisenstein part by $1+2^{11}$. The Kummer units are in \file{050_s21_pow_kummer_unit.oc}:
\octranscript{050_s21_pow_kummer_unit}
$37\mid B_{32}$ ($37$ is irregular): $v_\pi(\eta^{36}-1)=38>32$ and $\eta_{32}$ is not a $37$-th power; at $i=30$,
$37\nmid B_{30}$ and $v_\pi=30$, as Kummer's criterion predicts. The $\ell$-adic
$L$-function is in Section~\ref{sec:tour}.

\section{The residue atlas: \texttt{--atlas}}\label{sec:atlas}

\cli{--atlas} builds the residue atlas $\Sigma(\ell)$ (the anemic level-one eigensystems over weights
$\le\ell+1$, their $\theta$-twists and a separating prime set) outside the validation suite. It is not
listed by the top-level \cli{--help}.

% GENERATED by docs/latex/tools/gen_cli_manual_tables.sh, i.e.
% `CLI_MANUAL_EMIT=<dir> hacs CliManualCoverageValidation`, from the
% registries of the build it runs in.  Do not edit by hand: regenerate when a
% verb, type, setting, validation or command-line token changes.
\begin{clihelp}
usage: hacs --atlas <ell> [--n-max N] [--collect-only] [--quiet]
       hacs --atlas --close <store.json>
       hacs --atlas --verify <store.json>
       hacs --atlas --index <store.json> [--out <path>]
       hacs --atlas --query <store.json> '<request JSON>'   (OC_BUILD_NOTEBOOK=ON)
The store is written only under HECKE_ATLAS_OUT=<dir> (never under
docs/matrix-engine/results/).  --index writes the oc-atlas-index/1
sidecar (default <store.json>.ocidx; never the store itself, never
inside the source tree).  See HeckeExt/HeckeAtlasCli.h.
\end{clihelp}

% GENERATED by docs/latex/tools/gen_cli_manual_tables.sh, i.e.
% `CLI_MANUAL_EMIT=<dir> hacs CliManualCoverageValidation`, from the
% registries of the build it runs in.  Do not edit by hand: regenerate when a
% verb, type, setting, validation or command-line token changes.
\begin{tokentable}{Token}{Arguments}
\clitok{\textless{}ell\textgreater{}} & -- & Builds Sigma(ell) over every even n \textless{}= min(ell - 1, N); one ATLAS\_CELL line per cell. Requires 5 \textless{}= ell \textless{}= 100000 (otherwise usage, exit 1). \\
\clitok{-{}-{}n-{}max}\index{n-max@\texttt{-{}-{}n-{}max} (\texttt{-{}-{}atlas})} & \texttt{\textless{}N\textgreater{}} & Caps the weight range of the build. \\
\clitok{-{}-{}collect-{}only}\index{collect-only@\texttt{-{}-{}collect-{}only} (\texttt{-{}-{}atlas})} & -- & Stops after the collection (no closure, no S\_sep) and writes the \texttt{.collected.json} store. \\
\clitok{-{}-{}quiet}\index{quiet@\texttt{-{}-{}quiet} (\texttt{-{}-{}atlas})} & -- & Suppresses the per-cell lines. \\
\clitok{-{}-{}close}\index{close@\texttt{-{}-{}close} (\texttt{-{}-{}atlas})} & \texttt{\textless{}store.json\textgreater{}} & Re-derives the closure and S\_sep of a store without the pipeline. \\
\clitok{-{}-{}verify}\index{verify@\texttt{-{}-{}verify} (\texttt{-{}-{}atlas})} & \texttt{\textless{}store.json\textgreater{}} & Re-closes a closed store and compares byte for byte (exit 0 iff identical). \\
\clitok{-{}-{}index}\index{index@\texttt{-{}-{}index} (\texttt{-{}-{}atlas})} & \texttt{\textless{}store.json\textgreater{}} & Writes the \texttt{oc-{}atlas-{}index/1} sidecar of a store (default \texttt{\textless{}store.json\textgreater{}.ocidx}), keyed to the store's size and FNV; the store is opened read-only; an output that is the store itself, inside the source tree or under \texttt{docs/matrix-{}engine/results/} is refused (exit 1); a store that is not indexable exits 3. \\
\clitok{-{}-{}query}\index{query@\texttt{-{}-{}query} (\texttt{-{}-{}atlas})} & \texttt{\textless{}store.json\textgreater{} '\textless{}request JSON\textgreater{}'} & One-shot read-only query on a \texttt{hecke\_atlas/v1} store for scripts (slot-22 VERB-3, owner decision O13): the request handler of \texttt{-{}-{}serve-{}atlas} (\texttt{info}, \texttt{cells}, \texttt{cell}, \texttt{system}, \texttt{systems}; same validation and filters) run once, one \texttt{atlas\_reply} frame on stdout; the store is opened read-only and its path comes from the command line only. Exit 0 reply \texttt{ok}, 1 usage or a refused store, 3 an \texttt{error} / \texttt{rejected} reply; built only with \texttt{OC\_BUILD\_NOTEBOOK=ON} (otherwise exit 5). \\
\clitok{-{}-{}out}\index{out@\texttt{-{}-{}out} (\texttt{-{}-{}atlas})} & \texttt{\textless{}path\textgreater{}} & (\texttt{-{}-{}index}) the sidecar's path instead of \texttt{\textless{}store.json\textgreater{}.ocidx}. \\
\clitok{-{}-{}help}\index{help@\texttt{-{}-{}help} (\texttt{-{}-{}atlas})}, \clitok{-{}h}\index{h@\texttt{-{}h} (\texttt{-{}-{}atlas})} & -- & Prints the atlas usage. \\
\end{tokentable}


The store file is written only when \env{HECKE_ATLAS_OUT}\texttt{=<dir>} is set, and never under
\file{docs/matrix-engine/results/}. Every run prints an \texttt{ATLAS\_SUMMARY} line and an
\texttt{ATLAS\_STORE} line with the FNV-1a-64 of the store text. The atlas is an $\Fl$ object: under a
non-field \vb{coefficient\_ring} the command exits 5. In the console, \vb{atlas(ell)} builds
$\Sigma(\ell)$ in process for $\ell\in\{5,7,11,13\}$.

\paragraph{One-shot queries.} \cli{--atlas --query <store.json> '<request JSON>'} answers a single
read-only request against a \texttt{hecke\_atlas/v1} store and exits, for scripts that
do not want a server: the request is handled exactly as by \cli{--serve-atlas} (\texttt{info},
\texttt{cells}, \texttt{cell}, \texttt{system}, \texttt{systems}; same validation, same filters) and one
\texttt{atlas\_reply} frame is written to standard output. The store is opened read-only and its path comes
from the command line only. Exit 0 for an \texttt{ok} reply, 1 for usage or a refused store, 3 for an
\texttt{error} or \texttt{rejected} reply; without \texttt{OC\_BUILD\_NOTEBOOK=ON} it exits 5. A blank or
multi-line request is refused.

\paragraph{The index sidecar.} \cli{--atlas --index <store.json>} writes the byte-offset index
\texttt{oc-atlas-index/1} of a store beside it (\file{<store.json>.ocidx}, or \cli{--out <path>}): one line
per cell and per system with its offset, length and the attributes the notebook filters on ($r$, type,
collected, base, twist), the occurrences of the collected systems, the store's size and FNV-1a-64, and
a closing hash of the index itself ([P] v3.11 Section~18.16). The store is opened read-only; an output
that is the store itself, lies inside the source tree or under \file{docs/matrix-engine/results/} is
refused (exit~1); a file that is not a \texttt{hecke\_atlas/v1} store in the layout the build writes exits~3.
The sidecar is the owner's one-off for the large $\Sigma(691)$ store; \cli{--serve-atlas}
(Section~\ref{sec:serve-atlas}) reads it.

%======================================================================
\chapter{The console and its language}\label{ch:console}
%======================================================================

\section{The console}\label{sec:console}

The console is a small deterministic expression language, \emph{OC script}, over the whole pipeline.
Its binding design is [P]; \file{docs/console/console.md} is its generated Markdown reference. It has
four invocation forms and a few modifiers:

% GENERATED by docs/latex/tools/gen_cli_manual_tables.sh, i.e.
% `CLI_MANUAL_EMIT=<dir> hacs CliManualCoverageValidation`, from the
% registries of the build it runs in.  Do not edit by hand: regenerate when a
% verb, type, setting, validation or command-line token changes.
\begin{clihelp}
OC script console (console phase C1a).
Usage:
  hacs --console | -i          Interactive session (prompt 'oc> ').
  hacs --script <file.oc> [name=value]...
                                        Run a script; parameters are read as args.name.
  hacs -c '<code>'             Run one line of OC script.
  hacs -                       Run the script on standard input.
Modifiers:
  --set key=value                       Apply a session setting before the first statement.
  --log <file>                          Append the transcript (inputs prefixed '> ').
  --quiet | --verbose                   The pipeline's own chatter off / on.
  --script <file> --console             Run the script, then drop into the REPL.
  --help                                This text.
Exit codes: 0 ok, 1 usage, 2 parse error, 3 runtime error, 4 failed assert,
            5 a legal setting this build cannot execute.
Language and verb reference: `help` inside the session, `?name` for one name.
\end{clihelp}

% GENERATED by docs/latex/tools/gen_cli_manual_tables.sh, i.e.
% `CLI_MANUAL_EMIT=<dir> hacs CliManualCoverageValidation`, from the
% registries of the build it runs in.  Do not edit by hand: regenerate when a
% verb, type, setting, validation or command-line token changes.
\begin{tokentable}{Token}{Arguments}
\clitok{-{}-{}console}\index{console@\texttt{-{}-{}console} (console)}, \clitok{-{}i}\index{i@\texttt{-{}i} (console)} & -- & Interactive session; after \texttt{-{}-{}script \textless{}f\textgreater{}} it opens the REPL once the script has succeeded. \\
\clitok{-{}-{}script}\index{script@\texttt{-{}-{}script} (console)} & \texttt{\textless{}file.oc\textgreater{}} & Runs a script file (one of \texttt{-{}-{}script}, \texttt{-{}c}, \texttt{-{}}). \\
\clitok{-{}c}\index{c@\texttt{-{}c} (console)} & \texttt{\textless{}code\textgreater{}} & Runs one line of OC script. \\
\clitok{-{}}\index{-@\texttt{-{}} (console)} & -- & Runs the script read from standard input (stdin is the script, so no REPL is ever opened). \\
\clitok{-{}-{}set}\index{set@\texttt{-{}-{}set} (console)} & \texttt{key=value} & Sets a session setting before the first statement (repeatable); an unknown key or a bad value is exit 1, a legal value this build cannot execute exit 5. \\
\clitok{-{}-{}log}\index{log@\texttt{-{}-{}log} (console)} & \texttt{\textless{}file\textgreater{}} & Appends the transcript to a file, inputs prefixed \texttt{\textgreater{} }. \\
\clitok{-{}-{}quiet}\index{quiet@\texttt{-{}-{}quiet} (console)}, \clitok{-{}-{}verbose}\index{verbose@\texttt{-{}-{}verbose} (console)} & -- & The pipeline's own chatter off / on (default: on in the REPL, off in scripts). \\
\clitok{-{}-{}help}\index{help@\texttt{-{}-{}help} (console)}, \clitok{-{}h}\index{h@\texttt{-{}h} (console)} & -- & Prints the console usage. \\
\clitok{\textless{}name=value\textgreater{}} & -- & A script parameter, read as \texttt{args.name}: \texttt{true} / \texttt{false}, a decimal integer, a quoted string, otherwise a String. \\
\end{tokentable}


\paragraph{Interactive mode.} \cli{--console} (or \cli{-i}) prompts \texttt{oc> }, and \texttt{... }
while brackets are open or a line ends in \verb|\|. An expression statement prints its value; a trailing
\texttt{;} suppresses that. Every error is printed and the session stays usable; a verb-time error binds
nothing. \texttt{quit} ends the session. Lines are read with plain \texttt{getline}; run the console under
\texttt{rlwrap} for history and line editing. \texttt{Ctrl-C} during a statement stops it at its next
cancellation point (every statement and call, and the long verbs' loop steps: per Hecke prime, torsion prime,
level, stage) and prints \texttt{interrupted (at <site>)}; the prompt returns; the interrupted statement's own binding is not made, earlier effects of a compound statement stay as after any error (a lazy \texttt{descent} handle keeps its place and retries on the next read).
\texttt{Ctrl-C} at the idle prompt, or twice in a row, ends the process as before. Scripts and \cli{-c} are
not affected. \cli{--script f.oc --console} runs the script and then opens the REPL on its
bindings -- the natural way to explore a prepared state.

\paragraph{Script mode.} \cli{--script f.oc} runs a file; \cli{-c '<code>'} one line; \cli{-} standard
input. In a script \emph{every} error is fatal and sets the exit code (Section~\ref{sec:errors}); a parse
error in the script file itself is reported before anything runs, so a script with one has no side effect.
The exception is \vb{include("f.oc")}: the included file is parsed when the statement executes, so earlier
statements have already run when its parse error (exit~2) is reported. The pipeline's own
chatter is off by default in scripts (on in the REPL): a transcript must be a function of the script
alone.

\paragraph{Parameters.} \texttt{name=value} pairs after \cli{--script} become \vb{args.name}: \vb{true}
and \vb{false} are Bools, a decimal integer is an Int, a quoted string a String, anything else a String.
Reading an \vb{args.name} the command line did not supply is a runtime error (exit 3). For example
\file{Console/scripts/cli_params.oc} starts with

\excerptfrom{cli_params.oc}
\begin{ocexcerpt}
n = args.n;
ell = args.ell;
\end{ocexcerpt}
and is run as \cli{./hacs --script Console/scripts/cli\_params.oc n=10 ell=691}.

\paragraph{Logging and transcripts.} \cli{--log <file>} appends the session to a file, inputs prefixed
\texttt{> }. A line a script prints with the prefix \texttt{\#\textasciitilde{} } (the \vb{annotate} and
\vb{bench\_line} verbs) is an \emph{annotation}: timings and other machine-dependent text go there, and
every transcript comparison ignores it. The frozen transcripts of \file{Console/scripts/*.expected} are
the session output with annotations and trailing blanks removed, under one header line
\texttt{\#! oc-format 1}; they are regenerated only by the owner-gated emitter of
Section~\ref{sec:env-gated}.

\section{The OC language}\label{sec:language}

\subsection{Lexical structure}

A newline ends a statement unless brackets are open or the line ends in \verb|\|; \texttt{;} also ends
one and, after an expression statement, suppresses its printing. Comments run from \texttt{\#} or
\texttt{//} to the end of the line. Identifiers are \verb|[A-Za-z_][A-Za-z0-9_]*|, case-sensitive.
Reserved words:
\begin{center}
\texttt{for in if else set help quit true false and or not expect\_error include args def}
\end{center}
(\texttt{def} is reserved for a future user-function form and has no meaning yet.) Integers are decimal or
\texttt{0x} hexadecimal with optional \texttt{\_} separators (\texttt{134\_217\_757}); they are of type
\vb{Int} $=$ \texttt{int64\_t}, a literal outside $[-2^{63},2^{63}-1]$ is a parse error and arithmetic
overflow is an error, never a wrap. Strings are \texttt{"\dots"} with the escapes
\verb|\" \\ \n \t|; there is no interpolation -- use \vb{str(x)} and \texttt{+}.

\subsection{Grammar}

The normative grammar ([P] Sec.~3.2):
\begin{spec}
program    := { statement }
statement  := assign | exprstmt | for | if | setstmt | include | help | expect | quit
assign     := IDENT "=" expr TERM
quit       := "quit" TERM
exprstmt   := expr TERM
for        := "for" IDENT "in" expr block
if         := "if" expr block [ "else" ( block | if ) ]
setstmt    := "set" IDENT "=" expr TERM
include    := "include" "(" STRING ")" TERM
help       := ( "help" [ IDENT | "?" ] | "?" IDENT ) TERM
expect     := "expect_error" STRING block
block      := "{" { statement } "}"
TERM       := NEWLINE | ";"
expr       := orexpr
orexpr     := andexpr { "or" andexpr }
andexpr    := notexpr { "and" notexpr }
notexpr    := "not" notexpr | cmp
cmp        := arith [ ("==" | "!=" | "<" | "<=" | ">" | ">=") arith ]
arith      := term { ("+" | "-") term }
term       := power { ("*" | "/" | "%") power }
power      := unary [ "^" power ]                         (right-assoc)
unary      := "-" unary | postfix
postfix    := primary { "(" args ")" | "." IDENT | "[" expr "]" }
primary    := INT | STRING | "true" | "false" | IDENT | "args" | "(" expr ")"
            | "[" [ expr { "," expr } ] "]"
args       := [ arg { "," arg } ]           arg := expr | IDENT "=" expr
\end{spec}

\paragraph{Operators, from loosest to tightest.} \texttt{or}; \texttt{and}; \texttt{not}; one comparison
(\texttt{== != < <= > >=}, not chained); \texttt{+ -}; \texttt{* / \%}; \texttt{\^{}} (right-associative);
unary \texttt{-}; then calls \texttt{f(\dots)}, members \texttt{.name} and indexing \texttt{[i]}.

\begin{note}
\textbf{Two rules to know.}
(1) \emph{The unary minus binds inside the base of} \texttt{\^{}}: \texttt{power := unary [ "\^{}" power ]},
so \texttt{-2 \^{} 2 == 4} and \texttt{0 - 2 \^{} 2 == -4}. This is not the school convention.
(2) \emph{Strings print bare only at the top of} \vb{print} / \vb{str}: \vb{print("dim")} prints
\texttt{dim}, while a bare expression statement \texttt{"dim"}, and a string inside a list or record,
prints quoted.
\end{note}

\paragraph{Semantics.} \texttt{/} on \vb{Int} is exact and an error on a non-zero remainder; on \vb{Fp}
it is field division. \texttt{\^{}} on \vb{Int} is exact and overflow-checked. Comparison of objects is
mathematical equality (matrices entrywise, polynomials coefficientwise). Arguments are positional or
by name (\vb{bench("hecke", n = 48, ell = 134217757, p = 2)}); optional parameters have defaults.
Overloads are resolved in registration order: the first overload that accepts the arguments wins.

\subsection{Statements and control flow}
\begin{itemize}[nosep]
\item \textbf{Assignment} \texttt{x = expr} binds in the single global frame; there are no local scopes
  and no user functions.
\item \textbf{\texttt{if} / \texttt{else}}: \texttt{if cond \{ \dots \} else if cond2 \{ \dots \} else \{ \dots \}}; the
  condition must be a \vb{Bool}.
\item \textbf{\texttt{for}}: \texttt{for x in list \{ \dots \}} iterates a \vb{List} (\vb{range(a, b)} builds one).
\item \textbf{\texttt{include("f.oc")}} runs another file in the same global frame, resolved relative to
  the including file and cycle-checked.
\item \textbf{\texttt{set key = value}} changes a setting for the rest of the session
  (Section~\ref{sec:settings}).
\item \textbf{\texttt{expect\_error "substring" \{ \dots \}}} succeeds iff the block raises an error whose
  message contains the substring; it is the only error-consuming construct (there is no \texttt{try}),
  and a block that does \emph{not} fail is itself an error (exit 4).
\item \textbf{\texttt{help}}, \texttt{help <name>}, \texttt{help ?} and \texttt{?name} print the registry's
  documentation of a verb, type or setting; \texttt{?x} on a value describes it (type, state, settings
  snapshot, creation site, members) -- the same as \vb{describe(x)}.
\item \textbf{\texttt{quit}} ends a REPL session or a script (exit 0); \vb{exit(code)} ends a script with a
  chosen code; \vb{assert(cond, msg)} fails it with exit 4.
\end{itemize}

\subsection{Values and types}\label{sec:types}

The core value kinds are \vb{Nil}, \vb{Int}, \vb{Bool}, \vb{String}, \vb{List} (\texttt{[a, b]},
heterogeneous, indexed from 0, \vb{len}) and \vb{Record} (printed as \texttt{\{k: v, \dots\}}; read-only, members by
\texttt{.k}, \vb{keys}, \vb{json}). A Record has no literal: \texttt{x = \{a: 1\}} is a parse error
(exit~2). Records are produced by verbs and bindings (\vb{settings()}, \vb{state(x)}, \vb{golden(\dots)}, \vb{args}, the
pipeline's report objects). Everything else is an \emph{object} of a registered type: the
arithmetic types (\vb{Fp}, \vb{Poly}, \vb{Matrix}, \vb{GF}, \vb{GFElt}) and the pipeline types
(\vb{Model}, \vb{LocalSystem}, \vb{Cohomology}, \vb{HeckeOperator}, \vb{Family}, \vb{Decomposition},
\vb{Block}, \vb{Report}, \vb{ExtRecord}, \vb{MaximalIdeal}, \vb{HeckeAlgebra}, \vb{LoewyProfile},
\dots). This build registers \CMNumTypes{} types; Appendix~\ref{app:types} lists each with its members.

\paragraph{Lazy handles.} A pipeline verb returns at once with a \emph{handle} in state
\texttt{declared}; the work runs on the first access to a \emph{computed} member, under the settings
\emph{snapshot taken when the handle was created}. \emph{Structural} members are available while the
handle is still declared. A computation that fails leaves the handle in state \texttt{failed}: every later
access rethrows the stored error, naming both the access site and the creation site, and the handle is
never retried -- build a new one. \vb{compute(x)} forces a handle, \vb{state(x)} reports state, creation
site, settings snapshot and resumable state, and \vb{resume(x, budget)} continues a computation that a
budget cut off (Section~\ref{ex:settings}).

\subsection{Verbs}

Verbs are grouped by area; Appendix~\ref{app:verbs} lists all \CMNumOverloads{} overloads of the
\CMNumVerbs{} verbs with their documentation, in registration (= resolution) order:
\begin{description}[nosep,font=\normalfont\ttfamily]
\item[core] printing, introspection, lists and records, \vb{assert}, \vb{exit}, settings, \vb{compute} /
  \vb{state} / \vb{resume};
\item[arithmetic] $\F_p$ elements, polynomials (\vb{factor} is always certified), dense matrices over
  $\F_p$, linear algebra (\vb{rank}, \vb{kernel}, \vb{solve}, \vb{restrict}, \dots), commutative algebras
  generated by commuting operators, characteristic and minimal polynomials, seeded random data;
\item[pipeline] the $\SL(2,\Z)$ model and simplicial presets, local systems, \vb{cohomology}, \vb{hecke},
  \vb{family}, \vb{decompose}, \vb{report}, duality (\vb{dual}, \vb{pairing}), cross-checks;
\item[reproducibility] \vb{hash}, \vb{golden} (read-only access to the frozen reference),
  \vb{fingerprint}, \vb{math\_settings}, \vb{save\_state};
\item[harness] \vb{validations}, \vb{run\_validation}, \vb{bench}, \vb{micro}, the campaign verbs,
  \vb{golden\_subset}, \vb{timed}, \vb{rss};
\item[reference] \vb{help\_markdown} (the Markdown reference \file{docs/console/console.md});
\item[extension] the extension layer: \vb{sturm\_bound}, \vb{u\_ell\_decision}, \vb{ext\_record},
  \vb{local\_algebra}, \vb{module\_structure}, \vb{presentation}, the local-algebra invariants
  (\vb{radical}, \vb{socle}, \vb{loewy}, \dots), \vb{atlas}, the canonical fields \vb{GF}, \vb{iota},
  \vb{boundary}, \vb{cone}, \vb{dickson\_les}, \vb{theta}, \vb{hecke\_ring\_algebra}.
\end{description}
Inside a session, \vb{verbs()}, \vb{types()} and \vb{settings()} enumerate the registry, \texttt{help}
lists everything and \texttt{?name} documents one entry.

\subsection{Printing}\label{sec:printing}

Every value has a \emph{canonical} form -- what a transcript contains and what \vb{show}, \vb{repr} and an
expression statement print -- and a \emph{display} form, which differs only in that a top-level String
loses its quotes (\vb{print}, \vb{str}). Numbers are canonical residues in $[0,p)$; polynomials print from
the highest degree down with their modulus (\texttt{X\^{}3 + 72*X\^{}2 + 346*X + 4 (mod 691)}), although
\vb{coeffs} and \vb{poly} use the machine order low $\to$ high.
\begin{center}\small
\begin{tabular}{@{}>{\raggedright\arraybackslash}p{0.2\linewidth}>{\raggedright\arraybackslash}p{0.72\linewidth}@{}}
\toprule
\vb{print(a, b, ...)} & display forms on one line, separated by blanks\\
\vb{show(x)} & the full canonical form, ignoring \vb{print\_max\_entries}\\
\vb{display(a, b, ...)} & forces each declared handle, then shows each argument: as \vb{show} in a script or a terminal, as a rendered value in the notebook\\
\vb{describe(x)}, \texttt{?x} & type, state, settings snapshot, creation site, members\\
\vb{str(x)}, \vb{repr(x)} & the display / canonical form as a String\\
\vb{json(x)} & JSON for Int, Bool, String, List and Record (objects have their own \texttt{.json} members)\\
\vb{annotate(text)} & a \texttt{\#\textasciitilde{}} annotation line, ignored by transcript comparisons\\
\bottomrule
\end{tabular}
\end{center}
The engine setting \vb{print\_max\_entries} (default 400) caps what \vb{print} and expression statements
render: a larger object prints its summary form (\texttt{Matrix 1677 x 1419 over F\_134217757 (hash 0x\dots)}).

\paragraph{\texttt{display}: printing for the notebook.} A handle such as \vb{cohomology(L)} is \emph{declared}: nothing
is computed until a value is needed, and an expression statement prints \texttt{(declared, not yet computed)}.
\vb{display(x, ...)} first forces every declared handle among its arguments, exactly as \vb{compute(x)} does, and then shows
each argument on a line of its own. In a script or at the terminal that line is the full canonical form, as \vb{show}
prints it:
\excerptfrom{056_s24_display.oc}
\begin{ocexcerpt}
M = sl2z_model();
H = cohomology(symmetric_power(M, 10, 691));
H
display(H);
display(H.dims, 691, "dim H^1 = 2 dim S_12 + 1");
display(dickson_les(16, 7).dims);
\end{ocexcerpt}
prints (the frozen transcript)
\octranscript{056_s24_display}
In a served notebook (Section~\ref{sec:serve}) \vb{display} sends each argument as a \emph{value}, as an expression
statement does, so the page typesets it: $\dim_{\mathbb F_{691}}H^1=3$, a matrix as a matrix, an integral group as
$\Z^3\oplus\Z/2\oplus\Z/2^2\oplus\Z/3$, a character value as an element of the residue field. \vb{print} instead sends
plain text, so a notebook cell uses \vb{display} for the quantities and \vb{print} for the prose. In a served session the
text line of each argument is its canonical form under \vb{print\_max\_entries}, the echo an expression statement prints. That
line equals \vb{show}'s line, except for a value above the cap and for a type with a long form of its own, such as
\texttt{Model} and \texttt{Complex}. \vb{display} is served (\texttt{read\_only}) and changes no
value.

\subsection{Errors and exit codes}\label{sec:errors}

\begin{center}\small
\begin{tabular}{@{}c>{\raggedright\arraybackslash}p{0.86\linewidth}@{}}
\toprule
\textbf{Exit} & \textbf{Meaning}\\\midrule
0 & success\\
1 & usage: an unrecognised console argument, a script file not found, an undeclared \cli{--set} key, an invalid \cli{--set} value, or a \cli{--log} file that cannot be opened\\
2 & parse error -- in the script file itself reported before anything runs (an \vb{include}d file is parsed when the statement executes)\\
3 & runtime error: a type or arity error, a pipeline contract violation, every \texttt{ENGINE\_CHECK}, \texttt{hecke::HeckeOverflow}, out of memory, a failed handle\\
4 & a failed \vb{assert(cond, msg)} or \vb{expect\_error}\\
5 & a \emph{legal} setting this build cannot execute (\texttt{hecke::NotImplementedInThisBuild}): on a \texttt{set} statement, as a \cli{--set} argument, or at a verb that does not honour the setting in force; never reported as a bad value and never silently degraded to the default. Exceptions: \texttt{coefficient\_ring = Z/<m>} (reserved) is refused with exit~5 already at \texttt{set} / \cli{--set}, whereas \texttt{coefficient\_ring = Z} is \emph{accepted} there (stored; honoured by the integral descent only, {\ttfamily[C,~sec:recordschema]}, {\ttfamily[C,~def:ringsemantics]}(3)) and exits~5 at every verb other than \vb{descent}, naming the verb; and \texttt{field\_polynomial = "conway"} is \emph{accepted} there (the fingerprint moves) and only a verb that needs the field exits~5\\
\bottomrule
\end{tabular}
\end{center}
Every error is printed as \texttt{file:line:col: kind: message} followed by the offending source line; a
C++ exception names its class (\texttt{caused by: hecke::HeckeOverflow: \dots}), and an
\texttt{ENGINE\_CHECK} keeps its own \texttt{file:line}. Section~\ref{ex:errors} shows
\vb{expect\_error} over the documented error messages.

\section{A tour of the small verbs}\label{sec:tour}

The script \file{055_s24_language_tour.oc} runs once through the core, arithmetic and reproducibility verbs that the
larger examples take for granted: values and records, $\F_p$ elements and polynomials (printed from the highest degree
down, with their modulus), matrices, a finite field, the simplicial presets, hashes and the golden reference, and one
extension verb.
\octranscript{055_s24_language_tour}
Four verbs have no example by design: \vb{exit} (it ends the script), \vb{save\_state} (it writes a file),
\vb{help\_markdown} (its output is \file{docs/console/console.md}) and \vb{campaign\_spec} (it reads a spec file; the
same parser is shown by \vb{campaign\_spec\_text} in \file{048_c5c_ring_campaign.oc}).

\section{Settings}\label{sec:settings}

The console has one settings mechanism, \texttt{set key = value} (or \cli{--set key=value} on the
command line); \vb{settings()} shows every key with its value, source and family, and
\vb{reset\_settings()} returns to the session's baseline. There are two \emph{families}, and the
difference matters:

\begin{itemize}
\item \textbf{Engine settings} change how fast a value is produced, never the value: thread counts, the
  substrate, the device policy, the compiler passes of the tape engine. Their defaults come from the
  \texttt{ENGINE\_*} environment variables (Section~\ref{sec:env-engine}); some are \emph{env-only} and
  cannot be \texttt{set}.
\item \textbf{Mathematical settings} may change what a result \emph{means}: the search budget of the
  locality certificate, whether $U_\ell$ is included, the Sturm cap, the coefficient ring, \dots{} The
  engine owns their defaults and their validation; every record stores the values in force with their
  source; and a non-default value moves the \emph{fingerprint} (\vb{fingerprint()}, FNV-1a-64 of the
  non-default entries, \texttt{0x0000000000000000} at the defaults). The fingerprint enters every
  extension-record hash input and every campaign job id, and \vb{golden(\dots)} refuses to compare at a
  non-zero fingerprint: the frozen reference pins the defaults only.
\end{itemize}

A handle captures the settings in force \emph{when it is created}; changing a setting later does not
change an existing handle. The complete table (\CMNumSettings{} keys: \CMNumMathSettings{} mathematical,
\CMNumEngineSettings{} engine) is Appendix~\ref{app:settings}.

\subsection{The \CMNumMathSettings{} mathematical settings}\label{sec:mathsettings}

The legal values, from \file{Hecke/HeckeSettings.h}; the first one listed is the default.
{\small
\begin{longtable}{@{}>{\raggedright\arraybackslash}p{0.2\linewidth}>{\raggedright\arraybackslash}p{0.26\linewidth}>{\raggedright\arraybackslash}p{0.48\linewidth}@{}}
\toprule
\textbf{Key} & \textbf{Values} & \textbf{Meaning}\\\midrule\endhead
\vb{locality} & \vb{budget=64}, \vb{budget=<N>}, \vb{strict} & the locality certificate of a decomposition block (\Ecite{sec:localcert}): the witness-search budget, or the terminating strict mode (byte-identical wherever budget mode certifies)\\
\vb{u\_ell} & \vb{auto}, \vb{off}, \vb{budget\_mb=<M>} & whether the extension layer builds $U_\ell$: by a frozen memory model against 2048 MB, never, or against $M$ MB\\
\vb{sturm} & \vb{auto}, \vb{off}, \vb{cap=<B>} & the computed Sturm bound $B'(k,\ell)$ (\Ecite{def:bprime}): cap $12B(k)$, not attempted, or cap $B\ge2$\\
\vb{sturm\_unavailable} & \vb{nonfinal}, \vb{heuristic} & an ideal whose bound is not found at the cap stays non-final, or is swept heuristically\\
\vb{commutant\_budget} & \vb{20000000000}, \vb{<N>} & the work budget (modular multiplications) of the commutant certificate\\
\vb{final\_mode} & \vb{closure}, \vb{sturm\_heuristic} & \texttt{final = true} only on a proof, or \texttt{"heuristic"} on \Ecite{cert:sturm}'s heuristic range\\
\vb{field\_polynomial} & \vb{least}, \vb{conway} & the defining polynomial of the canonical field (\Ecite{def:canon}); \vb{conway} is legal but not executable (exit 5)\\
\vb{split\_field\_report} & \vb{canonical}, \vb{own\_field} & the primary form of reported eigenvalues (\Ecite{sec:field}); both forms are always emitted\\
\vb{presentation} & \vb{off}, \vb{on} & whether extension records carry module presentations\\
\vb{predictions} & \vb{annotate}, \vb{assert} & predictions are annotations, or a refuted prediction fails the run (a debug switch)\\
\vb{coefficient\_ring} & \vb{F\_ell}, \vb{Z/ell\^{}<k>} & coefficients in $\Fl$ or $\Z/\ell^k$, $2\le k\le4096$, the token \texttt{ell} literal (\Ecite{sec:ring}); \vb{Z} is stored and honoured by the integral descent only (the verb \vb{descent}, the campaign's descent cells, the record \texttt{integral\_descent\_record/v1}; {\ttfamily[C,~sec:recordschema]}, {\ttfamily[C,~def:ringsemantics]}(3)), every other verb refuses it by name with exit 5; the reserved value \vb{Z/<m>} is refused with exit 5 at \texttt{set} / \cli{--set}\\
\bottomrule
\end{longtable}
}

In the console a String setting takes a quoted value: \texttt{set locality = "budget=8"}; on the command
line \cli{--set locality=budget=8} (only the first \texttt{=} splits). A bad value is an ordinary error
naming the key (exit 3 on \texttt{set}, exit 1 on \cli{--set}); a legal value this build cannot execute
is exit 5. In a campaign spec the same keys appear as \texttt{key = value} lines.

\section{Errors on purpose: \texttt{expect\_error}}\label{ex:errors}

\excerptfrom{014_expect_error.oc}
\begin{ocexcerpt}
expect_error "is not a degree of this complex" { decompose(T2, 5) }
expect_error "two operators with the same prime" { decompose([T2, T2], 1) }
expect_error "no SL(2,Z) tessellation model" { hecke(cohomology(s1_scalar(1, 5)), 2).dims }
expect_error "unknown locality value" { set locality = "nonsense" }
expect_error "unknown setting 'no_such_key'" { set no_such_key = 1 }
...
set locality = "strict";
expect_error "the mathematical settings are not at their defaults" {
golden(POINT_COH, "dim_H^1")
}
set locality = "budget=64";
\end{ocexcerpt}
\octranscript{014_expect_error}
The transcript is one line because \vb{expect\_error} prints nothing when the expected error occurs. The
last block shows the reproducibility rule: under a non-default mathematical setting \vb{golden} refuses.

\section{Settings, the fingerprint and \texttt{resume}}\label{ex:settings}

\excerptfrom{021_settings_resume.oc}
\begin{ocexcerpt}
set threads = 1;
H1 = cohomology(symmetric_power(MODEL, N_SYM, ELL));
reset_settings();
print("threads = 1:", H1.dims);
...
set locality = "budget=8";
assert(fingerprint() != "0x0000000000000000", "a non-default key moves the fingerprint");
print(settings().locality.value, settings().locality.family);
set locality = "budget=64";
...
set field_polynomial = "conway";
expect_error "not_implemented_in_this_build" { canonical_field_polynomial(691, 2) }
set field_polynomial = "least";
print("C_{691,2} =", canonical_field_polynomial(691, 2));
...
CUT = sturm_bound(24, 5, 3);
...
CONT = resume(CUT, 36);
FRESH = sturm_bound(24, 5, 36);
...
print(CONT.B_prime, state(CONT).resumable.progress.budget_history,
state(FRESH).resumable.progress.budget_history);
\end{ocexcerpt}
\octranscript{021_settings_resume}
An engine setting (\vb{threads}) changes no value; a mathematical one (\vb{locality}) moves the
fingerprint; \vb{conway} is a legal value this build cannot execute (outside \vb{expect\_error} that is
exit 5); and a Sturm-bound computation cut off at cap 3 is resumed under cap 36 to the same $B'=5$ a fresh
run finds, with the budget history $[3, 36]$ recorded.

\section{The harness from the console}\label{ex:harness}

\excerptfrom{022_harness.oc}
\begin{ocexcerpt}
print("run_validation(mod_arithmetic_validation) =", run_validation("mod_arithmetic_validation"));
expect_error "no validation is registered" { run_validation("NoSuchValidation") }
expect_error "owner-gated emitter" { run_validation("GoldenEmit") }
B = bench("hecke", n = 48, ell = 134217757, p = 2);
bench_line(B);
print("bench(hecke, 48, 134217757, 2).dims =", B.dims);
...
G = golden_subset([[N_SYM, ELL]]);
print("golden_subset([[10, 691]]).mismatches =", G.mismatches);
\end{ocexcerpt}
\octranscript{022_harness}
\vb{bench\_line} prints the timing line as a \texttt{\#\textasciitilde{}} annotation, which is why no time
appears in the transcript; the owner-gated emitters cannot be run from a session. (The console can still write files, through
\vb{save\_state(x, path)} and \vb{run\_campaign(spec)}; only the emitters of Section~\ref{sec:env-gated} are out of reach.)

\section{The served console: \texttt{--serve}}\label{sec:serve}

\cli{--serve} runs \emph{one} console session driven by newline-delimited JSON on standard input and
output -- protocol \texttt{oc-serve/1}, specified in [P] v3.10 Section~18. It is the back end of the
notebook: a gateway (a separate program) starts the process, relays requests and frames, and owns the
network side; the binary never opens a socket. It is compiled only with the CMake option
\texttt{OC\_BUILD\_NOTEBOOK=ON}; in a default build the token is unknown (exit~1).

\paragraph{Page themes.} The notebook page's header carries a theme selector (\emph{Auto}, \emph{Light}, \emph{Dark}, two high-contrast presets, \emph{Custom}) and a \emph{Theme\ldots} editor. They change the page's colours only, in the browser: the choice is kept in the browser's \texttt{localStorage}, never sent to the gateway or the binary, and no verb, setting or command-line token is involved. A theme is exchanged as JSON of the format \texttt{oc-notebook-theme/1}; see \texttt{notebook/page/README.md} for the format and the tokens, and [P] Section~18.17 for the specification.

% GENERATED by docs/latex/tools/gen_cli_manual_tables.sh, i.e.
% `CLI_MANUAL_EMIT=<dir> hacs CliManualCoverageValidation`, from the
% registries of the build it runs in.  Do not edit by hand: regenerate when a
% verb, type, setting, validation or command-line token changes.
\begin{tokentable}{Token}{Arguments}
\clitok{-{}-{}set}\index{set@\texttt{-{}-{}set} (\texttt{-{}-{}serve})} & \texttt{key=value} & Sets a SETTABLE key before the \texttt{hello} frame (repeatable): \texttt{print\_max\_entries}, \texttt{lean}, \texttt{program\_cache} and the eleven mathematical settings; any other key is exit 1, a legal value this build cannot execute exit 5. \\
\clitok{-{}-{}rss-{}cap-{}mib}\index{rss-cap-mib@\texttt{-{}-{}rss-{}cap-{}mib} (\texttt{-{}-{}serve})} & \texttt{\textless{}N\textgreater{}} & The resident-set cap of the watchdog (polled every 250 ms and at every statement boundary): above it a \texttt{fatal} frame \texttt{rss\_cap} and exit 70; 0 (the default) is off. \\
\clitok{-{}-{}max-{}reply-{}bytes}\index{max-reply-bytes@\texttt{-{}-{}max-{}reply-{}bytes} (\texttt{-{}-{}serve})} & \texttt{\textless{}N\textgreater{}} & The frame cap (default 1048576, at least 4096): longer \texttt{out} text is split at line boundaries, a larger value is sent in its truncated form. \\
\end{tokentable}


\paragraph{Requests and frames.} Every request is one JSON object per line with an \texttt{op}:
\texttt{eval} (\texttt{src}, optional \texttt{args}) runs a cell as one program, parsed in full first;
\texttt{complete}, \texttt{help} and \texttt{describe} answer from the registry and the session without
evaluating; \texttt{interrupt} stops the running cell at its next cancellation point (a statement, a verb call, or a loop step inside a long verb); \texttt{close}
ends the session. The binary answers with frames: one \texttt{hello} first, then per request
\texttt{out} (what a script would print; annotations go to \texttt{progress}), \texttt{value} (the
JSON envelope \texttt{\{t, canonical, hash, v\}} of each echoed value; no \TeX, no floats, integers above
$2^{53}-1$ as decimal strings) and exactly one \texttt{done}. The \texttt{out} text of a cell is byte for
byte what \cli{--script} prints; \vb{ConsoleServeValidation} compares it with every frozen transcript.
A \vb{display(x, ...)} call (Section~\ref{sec:printing}) sends, per argument, the same canonical line and
\texttt{value} frame that the expression statement \texttt{x} would, after forcing a declared handle.

\paragraph{What the page renders.} The page turns the typed envelopes into \TeX{} (the binary emits none). It typesets
numbers, residues, polynomials, matrices, finite fields, local systems and cohomology, and the parts that render hints
mark as polynomials, $q$-series, abelian groups, Hilbert functions, Loewy layers, Jordan types, character values and the
maps of the Dickson long exact sequence. Records become tables. Each type's decision is listed in
\file{src/hacs/Console/serve/render_coverage.tsv}, and \vb{ConsoleReadOnlyTagValidation} keeps that list
complete. A declared value carries a \emph{compute} button, which evaluates \vb{display} of its statement. A matrix too
large for \vb{print\_max\_entries} carries a \emph{show all} button, which displays it once under a raised cap (up to
100\,000 entries) and then restores the setting. The tour \file{notebook/examples/notebook/} (ten short files, from the
model to the integral Hecke structure) is written for this mode.

\paragraph{Errors and exit codes.} A cell's error is a \texttt{done} with \texttt{status} and the exit code
of Section~\ref{sec:errors} (2 parse, 3 runtime, 4 assertion, 5 not implemented) and the session survives
it; 1 is a refusal of the request (nothing ran: \texttt{include}, \texttt{quit}, \vb{exit}, a verb that is
not served, a malformed request, a full queue). The process exits 0 on end of input, \texttt{close} or
\texttt{SIGTERM}, 1 when its launch is refused and 70 after a \texttt{fatal} frame (the RSS cap, a request
line over 8~MiB, an internal error).

\paragraph{Read-only.} A served session computes and displays; it writes no file and reaches no emitter.
Every verb overload carries a \texttt{read\_only} tag (Appendix~\ref{app:verbs} marks it
\textsf{[served]} / \textsf{[not served]}): \CMNumReadOnly{} of the \CMNumOverloads{} overloads, in
\CMNumReadOnlyVerbs{} verbs, are served; the other \CMNumNotReadOnly{} (the two overloads of \vb{exit},
\vb{report(D, level\_a, path)} --- the pathless \vb{report(D, level\_a)} is served and computes the record in
memory ---, \vb{save\_state(x, path)}, \vb{run\_validation}, \vb{bench}, \vb{micro}, \vb{campaign\_spec},
\vb{run\_campaign}) are refused. Only \texttt{print\_max\_entries}, \texttt{lean}, \texttt{program\_cache}
and the eleven mathematical settings may be \texttt{set}. The launch is refused (exit~1, nothing on
standard output) when the environment holds any \env{GOLDEN_*}, \env{CONSOLE_*}, \env{CLI_MANUAL_*},
\env{HECKE_*}, \env{ENGINE_*}, \env{OC_*} or sanitizer variable other than \env{ENGINE_THREADS},
\env{ENGINE_CUDA_DISABLE}, \env{ENGINE_CUDA_DEVICE} and \env{ENGINE_HOST_THREADS}; the process then sets
\texttt{RLIMIT\_FSIZE} and \texttt{RLIMIT\_CORE} to 0 and \texttt{PR\_SET\_NO\_NEW\_PRIVS}. This is a
safety rail against accidents, not a sandbox. The machine-readable twin of the tag column is
\file{docs/latex/cli_manual/read_only_verbs.json}.

\section{The atlas server: \texttt{--serve-atlas}}\label{sec:serve-atlas}

\cli{--serve-atlas} answers read-only queries on one \texttt{hecke\_atlas/v1} store -- protocol
\texttt{oc-serve-atlas/1}, [P] v3.11 Section~18.16 -- so that the notebook's atlas panel can browse the
$\Sigma(691)$ store (451\,875\,793 bytes) without ever sending it to a browser. Like \cli{--serve} it is
compiled only with \texttt{OC\_BUILD\_NOTEBOOK=ON}, speaks newline-delimited JSON on standard input and
output, and is started by the gateway (\texttt{--atlas-store}); the store's path comes from the command
line, never from a request.

% GENERATED by docs/latex/tools/gen_cli_manual_tables.sh, i.e.
% `CLI_MANUAL_EMIT=<dir> hacs CliManualCoverageValidation`, from the
% registries of the build it runs in.  Do not edit by hand: regenerate when a
% verb, type, setting, validation or command-line token changes.
\begin{tokentable}{Token}{Arguments}
\clitok{-{}-{}store}\index{store@\texttt{-{}-{}store} (\texttt{-{}-{}serve-{}atlas})} & \texttt{\textless{}store.json\textgreater{}} & The \texttt{hecke\_atlas/v1} store to serve (required): opened \texttt{O\_RDONLY}, mapped read-only, its FNV-1a-64 computed at open; the path comes from the command line only, never from a request. \\
\clitok{-{}-{}index}\index{index@\texttt{-{}-{}index} (\texttt{-{}-{}serve-{}atlas})} & \texttt{\textless{}path\textgreater{}} & The \texttt{oc-{}atlas-{}index/1} sidecar (default \texttt{\textless{}store.json\textgreater{}.ocidx} when it exists, else an index built in memory); compared field by field with the store at open, a stale or corrupt one refuses the launch (exit 1). \\
\clitok{-{}-{}expect-{}fnv}\index{expect-fnv@\texttt{-{}-{}expect-{}fnv} (\texttt{-{}-{}serve-{}atlas})} & \texttt{0x\textless{}hex\textgreater{}} & Refuses the launch unless the store's FNV-1a-64 equals this value (Sigma(691): \texttt{0x73c7210047a81b97}). \\
\clitok{-{}-{}max-{}reply-{}bytes}\index{max-reply-bytes@\texttt{-{}-{}max-{}reply-{}bytes} (\texttt{-{}-{}serve-{}atlas})} & \texttt{\textless{}N\textgreater{}} & The frame cap (default 1048576, at least 4096): a page is cut before the frame would exceed it (\texttt{next} continues). \\
\clitok{-{}-{}rss-{}cap-{}mib}\index{rss-cap-mib@\texttt{-{}-{}rss-{}cap-{}mib} (\texttt{-{}-{}serve-{}atlas})} & \texttt{\textless{}N\textgreater{}} & The resident-set cap of the watchdog (polled every 250 ms and after every request): above it an \texttt{atlas\_fatal} frame \texttt{rss\_cap} and exit 70; 0 (the default) is off. \\
\end{tokentable}


\paragraph{Requests and frames.} After one \texttt{atlas\_hello} (the store's name, size, FNV and counts,
the index source \texttt{sidecar} or \texttt{memory}) every request -- \texttt{info} (the store header,
the type and residue-degree profile), \texttt{cells} and \texttt{systems} (pages with \texttt{offset},
\texttt{limit} $\le 1000$ and, for systems, a \texttt{filter} on \texttt{r}, \texttt{type},
\texttt{collected}, \texttt{base}, \texttt{twist} or the cell \texttt{n}), \texttt{cell} (one cell and
its occurrences), \texttt{system} (one system by \texttt{sid}) -- gets exactly one \texttt{atlas\_reply}
carrying the store's FNV; \texttt{close} ends the process. A page is cut before a frame would exceed
\cli{--max-reply-bytes}, and \texttt{next} says where to continue. Errors (\texttt{protocol},
\texttt{not\_found}, \texttt{too\_large}) are replies; the process goes on.

\paragraph{Read-only.} The launch applies the environment rule and the hardening of \cli{--serve}
(Section~\ref{sec:serve}) and is refused (exit~1, nothing on standard output) on a forbidden variable, an
unknown token, a store that does not open or is not indexable, an FNV that differs from
\cli{--expect-fnv}, or a stale or corrupt sidecar: the index is re-derived from the store at every open
and compared field by field. The store is opened \texttt{O\_RDONLY} and mapped without write permission;
the process cannot extend any file.

\section{The scripts, by subject}\label{sec:morescripts}

Every numbered script is a frozen, tested example; read them as recipes.
\begin{center}\small
\begin{tabular}{@{}>{\raggedright\arraybackslash}p{0.26\linewidth}>{\raggedright\arraybackslash}p{0.66\linewidth}@{}}
\toprule
\textbf{Subject} & \textbf{Scripts}\\\midrule
cohomology & \file{001_dimensions.oc}, \file{013_presets.oc}, \file{055_s24_language_tour.oc}\\
Hecke operators & \file{002_hecke_pins.oc}, \file{003_bases_and_projectors.oc}, \file{004_family_identities.oc}, \file{005_linear_algebra.oc}, \file{006_flagship.oc}, \file{011_decompose_certify.oc}, \file{012_report_schema.oc}, \file{015_strict_locality.oc}\\
the golden grid & \file{007_grid_classical.oc}, \file{008_grid_dickson.oc}, \file{009_grid_small_ell.oc}, \file{010_grid_edge_n.oc}\\
$\Z/\ell^k$ & \file{024_e8_ring_points.oc}, \file{025_e8_console_surface.oc}, \file{026_c5_leftovers_ring_members.oc}, \file{028_wide_band.oc}, \file{029_hecke_ring_algebra.oc}, \file{048_c5c_ring_campaign.oc}, \file{049_lv3_iota_ring.oc}\\
over $\Z$ & \file{054_s24_descent.oc}\\
levels, Hida, forms & \file{050_s21_pow_kummer_unit.oc}, \file{051_s21_hida_tower.oc}, \file{052_s22_level_n_nebentypus.oc}, \file{053_s22_forms_hida_lambda.oc}\\
extension layer & \file{016_extensions.oc}, \file{017_predictions.oc}, \file{018_e0_locality_grid.oc}, \file{019_e1_stage1_sturm.oc}, \file{020_e2_ramanujan_10_691.oc}, \file{023_e7_duality_10_691.oc}, \file{027_dickson_theta.oc}, \file{040_c5_ext_record_10_691.oc} to \file{047_c5_local_algebra_verbs.oc}\\
the language & \file{014_expect_error.oc}, \file{021_settings_resume.oc}, \file{022_harness.oc}, \file{cli_params.oc}, \file{prelude.oc}\\
\bottomrule
\end{tabular}
\end{center}
(\file{041_c5_stage9_kummer_cells.oc}, \file{042_c5_dickson_class_cone.oc}, \file{043_c5_boundary_coincidences.oc},
\file{044_c5_multiplicity_two.oc}, \file{045_c5_pd_loewy_reversal.oc} and \file{046_c5_identities.oc} complete the range
040--047.)

%======================================================================
\chapter{Campaigns: \texttt{--campaign}}\label{ch:campaigns}
%======================================================================

\section{What a campaign is}\label{sec:campaign}

A campaign is a sweep over $(n,\ell,p,r,q)$ described by a small text file, the \emph{spec}. It runs
every job through the same public API as \cli{--bench hecke} and writes down the mathematics: one JSON
per job with the canonical objects and their hashes, the report levels (a)/(b)/(c), the certificates
that ran, timings and peak RSS, plus a campaign-level \file{index.json} and \file{summary.tsv}.

% GENERATED by docs/latex/tools/gen_cli_manual_tables.sh, i.e.
% `CLI_MANUAL_EMIT=<dir> hacs CliManualCoverageValidation`, from the
% registries of the build it runs in.  Do not edit by hand: regenerate when a
% verb, type, setting, validation or command-line token changes.
\begin{clihelp}
Usage: hacs --campaign <spec-file> [options]

Runs a parameter sweep of (n_sym, ell, p, r, q) Hecke jobs through the same
public API as `--bench hecke`, writing one streamed JSON per job plus
index.json and summary.tsv. Spec grammar and worked examples:
  docs/matrix-engine/campaign.md

Options:
  --output <dir>   override the spec's `output` directory
  --jobs <n>       override `jobs`: concurrent jobs per shape group
  --threads <n>    override `threads`: engine workers per concurrent job
  --dry-run        print the plan (groups and jobs) and exit
  --quiet          suppress per-job progress lines
  --help           this text

Spec example:
  name    = t2_sweep
  output  = out/t2_sweep
  certify = sampled(every 10)
  lean    = true
  n_sym   = 0..200 step 2
  ell     = 5,13,691,10007,134217757
  p       = 2
  r       = 1
\end{clihelp}

% GENERATED by docs/latex/tools/gen_cli_manual_tables.sh, i.e.
% `CLI_MANUAL_EMIT=<dir> hacs CliManualCoverageValidation`, from the
% registries of the build it runs in.  Do not edit by hand: regenerate when a
% verb, type, setting, validation or command-line token changes.
\begin{tokentable}{Token}{Arguments}
\clitok{-{}-{}output}\index{output@\texttt{-{}-{}output} (\texttt{-{}-{}campaign})} & \texttt{\textless{}dir\textgreater{}} & Overrides the spec's \texttt{output} directory. \\
\clitok{-{}-{}jobs}\index{jobs@\texttt{-{}-{}jobs} (\texttt{-{}-{}campaign})} & \texttt{\textless{}n\textgreater{}} & Overrides \texttt{jobs}: concurrent jobs per shape group. \\
\clitok{-{}-{}threads}\index{threads@\texttt{-{}-{}threads} (\texttt{-{}-{}campaign})} & \texttt{\textless{}n\textgreater{}} & Overrides \texttt{threads}: engine workers per concurrent job. \\
\clitok{-{}-{}dry-{}run}\index{dry-run@\texttt{-{}-{}dry-{}run} (\texttt{-{}-{}campaign})} & -- & Prints the plan (groups and jobs) and exits without computing. \\
\clitok{-{}-{}quiet}\index{quiet@\texttt{-{}-{}quiet} (\texttt{-{}-{}campaign})} & -- & Suppresses the per-job progress lines. \\
\clitok{-{}-{}help}\index{help@\texttt{-{}-{}help} (\texttt{-{}-{}campaign})}, \clitok{-{}h}\index{h@\texttt{-{}h} (\texttt{-{}-{}campaign})} & -- & Prints the campaign usage. \\
\end{tokentable}


\section{The spec format}\label{sec:spec}

The spec is line-oriented:
\begin{spec}
# or //          comment to end of line
---              on a line by itself: start a new SECTION
key = value      value is a comma-separated list whose elements are literals
                 or ranges `a..b [step c]` (step defaults to 1)
\end{spec}
\emph{Grid keys} (\vb{n\_sym}, \vb{ell}, \vb{p}, \vb{r}, optional \vb{q}) belong to the current section;
the campaign is the union of the sections' Cartesian products, and every section must give the four
mandatory axes. \emph{Campaign keys} may appear anywhere; the last assignment wins. Duplicate jobs are
removed (first occurrence wins). Every key the parser accepts:

% GENERATED by docs/latex/tools/gen_cli_manual_tables.sh, i.e.
% `CLI_MANUAL_EMIT=<dir> hacs CliManualCoverageValidation`, from the
% registries of the build it runs in.  Do not edit by hand: regenerate when a
% verb, type, setting, validation or command-line token changes.
\begin{tokentable}{Key}{Value / default}
\clitok{n\_sym}\index{n_sym@\texttt{n\_sym} (campaign spec key)}, \clitok{n}\index{n@\texttt{n} (campaign spec key)} & list (grid) & Symmetric-power degree n (d\_V = n + 1); every section must give it. \\
\clitok{ell}\index{ell@\texttt{ell} (campaign spec key)} & list (grid) & Coefficient prime; ell \textgreater{}= 2\textasciicircum{}32 selects the 64-bit pipeline. \\
\clitok{p}\index{p@\texttt{p} (campaign spec key)}, \clitok{p\_hecke}\index{p_hecke@\texttt{p\_hecke} (campaign spec key)} & list (grid) & Hecke prime; p = ell is U\_ell, the same code path. \\
\clitok{r}\index{r@\texttt{r} (campaign spec key)} & list (grid) & Exponent of T\_\{p\textasciicircum{}r\}; r \textgreater{}= 2 adds the algebra recursion (32-bit ell only). \\
\clitok{q}\index{q@\texttt{q} (campaign spec key)}, \clitok{degrees}\index{degrees@\texttt{degrees} (campaign spec key)} & list (grid), optional & Degrees to report (default: all of the model's); a non-default list is part of the job id. \\
\clitok{level}\index{level@\texttt{level} (campaign spec key)} & list (grid), coefficient\_ring = Z only & (slot-24 D6) The level N of Gamma\_0(N) of a descent cell (default 1); admitted only under \texttt{coefficient\_ring = Z}, where each (level, n\_sym) cell writes one integral\_descent\_record/v1 (campaign.md Sec. 17); a spec error everywhere else. \\
\clitok{parts}\index{parts@\texttt{parts} (campaign spec key)} & default \texttt{all}, coefficient\_ring = Z only & (slot-24 ICP2) \texttt{groups}, \texttt{groups,hecke} or \texttt{all}: which parts of the integral descent record a descent cell computes (the Hecke data char\_Z and the torsion modules; the boundary pair). A part not run is marked \textquotedbl{}ran\textquotedbl{}:false in the record; a non-default value is part of the cell id; a spec error outside \texttt{coefficient\_ring = Z}. \\
\clitok{structure}\index{structure@\texttt{structure} (campaign spec key)} & default \texttt{none}, coefficient\_ring = Z only & (slot-25 HZ-REC) \texttt{none} or \texttt{hecke}: under \texttt{coefficient\_ring = Z}, \texttt{hecke} makes each (level, n\_sym, algebra) cell write one integral\_hecke\_structure/v1 record instead of a descent record (campaign.md Sec. 18); it selects the record schema and is not a mathematical setting; a spec error outside \texttt{coefficient\_ring = Z}. \\
\clitok{algebra}\index{algebra@\texttt{algebra} (campaign spec key)} & list (grid), structure = hecke only & (slot-25 HZ-REC) \texttt{full}, \texttt{anemic} or both: the Hecke algebra of a hecke-structure cell (default \texttt{full} at level 1, \texttt{anemic} above; \texttt{anemic} at level 1 is a spec error). \\
\clitok{hz\_cert}\index{hz_cert@\texttt{hz\_cert} (campaign spec key)} & default off, structure = hecke only & (slot-25 HZ-EMIT) \texttt{on} also writes the certificate files of each record into \textless{}output\textgreater{}/hzcert/ for lean/tools/hzcert\_to\_lean.py; an output knob, no value depends on it. \\
\clitok{hz\_g8}\index{hz_g8@\texttt{hz\_g8} (campaign spec key)} & default \texttt{off}, structure = hecke only & (slot-25 H1) \texttt{off} or \texttt{hida}: the record's settings value heckez.G8. \texttt{off} writes no G8 field (the default records); \texttt{hida} adds H1's Hida block per maximal ideal at level 1 (ordinary, chain\_key, g, L\_k, v\_Bk, structure, I0\_generator), a non-default value with its own settings fingerprint and cell ids, never golden. \\
\clitok{name}\index{name@\texttt{name} (campaign spec key)} & default \texttt{campaign} & Echoed into \texttt{index.json} and every job file. \\
\clitok{output}\index{output@\texttt{output} (campaign spec key)}, \clitok{output\_dir}\index{output_dir@\texttt{output\_dir} (campaign spec key)} & default \texttt{campaign\_out} & The output directory (created). \\
\clitok{model}\index{model@\texttt{model} (campaign spec key)} & default \texttt{sl2z\_7vertex} & The model descriptor; an unknown name is a spec error. \\
\clitok{resume}\index{resume@\texttt{resume} (campaign spec key)} & default \texttt{false} & Skip jobs whose output file exists and carries the completion trailer. \\
\clitok{lean}\index{lean@\texttt{lean} (campaign spec key)}, \clitok{memory\_lean}\index{memory_lean@\texttt{memory\_lean} (campaign spec key)} & default \texttt{false} & Memory-lean mode, as \texttt{-{}-{}bench -{}-{}lean}. \\
\clitok{threads}\index{threads@\texttt{threads} (campaign spec key)} & default \texttt{0} & Engine workers per concurrent job; 0 shares the process pool (\texttt{ENGINE\_THREADS}). \\
\clitok{jobs}\index{jobs@\texttt{jobs} (campaign spec key)}, \clitok{concurrent\_jobs}\index{concurrent_jobs@\texttt{concurrent\_jobs} (campaign spec key)} & default \texttt{1} & Jobs of one shape group that run concurrently. \\
\clitok{certify}\index{certify@\texttt{certify} (campaign spec key)} & default \texttt{full} & \texttt{full}, \texttt{sampled(\textless{}fraction\textgreater{})}, \texttt{sampled(every \textless{}k\textgreater{})}, \texttt{sampled(\textless{}k\textgreater{})} or \texttt{none}: which jobs run the external cross-checks (the in-process certificates always run). \\
\clitok{level\_a}\index{level_a@\texttt{level\_a} (campaign spec key)} & default \texttt{auto} & \texttt{auto}, \texttt{always} or \texttt{never} (\texttt{off} is accepted as an alias of \texttt{never}): whether a job report carries level (a). \\
\clitok{level\_a\_max\_entries}\index{level_a_max_entries@\texttt{level\_a\_max\_entries} (campaign spec key)} & default \texttt{16000000} & The \texttt{auto} size gate of level (a), in matrix entries. \\
\clitok{theta\_rank\_max\_dim}\index{theta_rank_max_dim@\texttt{theta\_rank\_max\_dim} (campaign spec key)} & default \texttt{2048} & rank(Theta\_flat\textasciicircum{}q) is computed only below this dimension. \\
\clitok{ext}\index{ext@\texttt{ext} (campaign spec key)} & default \texttt{false} & Run the extension engine once per (n\_sym, ell) cell; \texttt{false} leaves the output byte-identical to the pre-extension driver. \\
\clitok{family}\index{family@\texttt{family} (campaign spec key)} & default \texttt{2,3,5,7} & The good primes of S\_0 for the ext pass; the literal token \texttt{ell} admits U\_ell through the \texttt{u\_ell} budget. \\
\clitok{sign}\index{sign@\texttt{sign} (campaign spec key)} & default \texttt{all} & \texttt{all}, \texttt{plus} or \texttt{minus}: which iota-sign record a cell reports. \\
\clitok{ext\_groups}\index{ext_groups@\texttt{ext\_groups} (campaign spec key)} & default \texttt{all} & \texttt{bulk}, \texttt{bulk+bd} or \texttt{all}; \texttt{bulk} leaves the cone unbuilt. \\
\clitok{ext\_jobs}\index{ext_jobs@\texttt{ext\_jobs} (campaign spec key)} & default \texttt{1} & Scheduling knob: cells of the ext pass run at once (changes no value, id or file). \\
\clitok{ext\_memory\_mb}\index{ext_memory_mb@\texttt{ext\_memory\_mb} (campaign spec key)} & default \texttt{40000} & Scheduling knob: the admission cap of the parallel ext pass, in MB. \\
\clitok{ext\_gpu}\index{ext_gpu@\texttt{ext\_gpu} (campaign spec key)} & default \texttt{off} & Scheduling knob: \texttt{off}, \texttt{on} or \texttt{force}, one extra device worker; never on by default. \\
\clitok{ext\_gpu\_cells}\index{ext_gpu_cells@\texttt{ext\_gpu\_cells} (campaign spec key)} & \texttt{n:ell} pairs & Scheduling knob: the cells the device worker takes; needs \texttt{ext\_gpu} on or force. \\
\clitok{orbifold}\index{orbifold@\texttt{orbifold} (campaign spec key)} & default absent (\texttt{off}) & \texttt{c\_i}: write one \texttt{orbifold\_cell/v2} record per (n\_sym, ell) cell, ell in \{2, 3\} -- group cohomology on C\_I plus the degree-1 amalgam reader -- instead of Hecke jobs; no \texttt{p} / \texttt{r} / \texttt{q} axis. Absent or \texttt{off} leaves the campaign byte-identical. \\
\clitok{orbifold\_primes}\index{orbifold_primes@\texttt{orbifold\_primes} (campaign spec key)} & default \texttt{2,3,5,7} & The T\_p of an orbifold cell (T\_ell is always included); a non-default list changes the cell id (\texttt{\_s\textless{}16 hex\textgreater{}}) and the record hash input. \\
\clitok{orbifold\_m\_target}\index{orbifold_m_target@\texttt{orbifold\_m\_target} (campaign spec key)} & default \texttt{auto} & The certified ell-adic precision every cell reaches (\texttt{auto} = max(floor((n+1)/(ell+1)) + 10, 40)); a divisor not certified below it is a \texttt{fallback}. Non-default: part of the cell id. \\
\clitok{orbifold\_k\_cap}\index{orbifold_k_cap@\texttt{orbifold\_k\_cap} (campaign spec key)} & default \texttt{0} & The largest working precision K the reader may climb to (0 = its instantiated range, 1024 at ell = 2, 312 at ell = 3). Non-default: part of the cell id. \\
\end{tokentable}


In addition, every registered \textbf{mathematical setting} (Section~\ref{sec:mathsettings}:
\vb{locality}, \vb{u\_ell}, \vb{sturm}, \dots, \vb{coefficient\_ring}) is accepted as a campaign-wide
key and validated by the engine itself; an unknown key is a spec error. A non-default mathematical key
changes every job id (a settings suffix) and the shape signature; at the defaults nothing is appended.
\vb{coefficient\_ring = Z/ell\^{}<k>} runs the job loop over $\Z/\ell^k$ (\Ecite{sec:ring}) and is
refused together with \vb{ext = true}, before any job runs.

A complete example, the one the usage text prints:
\begin{spec}
name    = t2_sweep
output  = out/t2_sweep
certify = sampled(every 10)
lean    = true
n_sym   = 0..200 step 2
ell     = 5,13,691,10007,134217757
p       = 2
r       = 1
\end{spec}
Run \cli{--campaign t2.spec --dry-run} first: it prints the shape groups and jobs and exits.

\section{What \texttt{certify} samples}\label{sec:certify}

The in-process certificates are never sampled: $\partial\tau=\tau\partial$, the agreement of the two
$\Theta$ routes on every entry, the exact $\Theta_{\mathrm{eq}}$ residual and the decomposition
certificate run on every job, and each job file records their executed counts. \vb{certify} only selects
the jobs that additionally run the \emph{external} cross-checks (the classical $\alpha/\beta/\omega$
factorisation and, for $r\ge2$, the direct double-coset route): \vb{full} every job;
\vb{sampled(every k)} or \vb{sampled(k)} the jobs with index $\equiv 0 \bmod k$; \vb{sampled(1/4)} the
same with $k=\mathrm{round}(1/\text{fraction})$; \vb{none} no job.

\section{Output}\label{sec:campaign-out}

\begin{spec}
<output>/
  index.json            campaign metadata, the shape groups, one row per job
  summary.tsv           one row per job -- the published table
  jobs/<job-id>.json    one streamed JSON per job
  ext/hecke_ext_ell<L>_n<N>.json      (only with ext = true) one record per cell
\end{spec}
A job id is \texttt{n<N>\_ell<L>\_p<P>\_r<R>}, followed by \texttt{\_q<d1>-<d2>-\dots} when the spec
names a non-default degree list and by a settings suffix when a mathematical key is not at its default.
The id is the file stem, the \texttt{id} column of \file{summary.tsv} and the \emph{resume key}:
with \vb{resume = true} a job is skipped iff \file{jobs/<id>.json} exists and carries the completion
trailer, so an interrupted campaign is restarted with the same command. Over $\Z/\ell^k$ the job file
schema is \texttt{hecke\_ring\_report/v1}; with \vb{ext = true} the cell records are
\texttt{hecke\_ext/v5}. The scheduling knobs (\vb{jobs}, \vb{threads}, \vb{ext\_jobs},
\vb{ext\_memory\_mb}, \vb{ext\_gpu}, \vb{ext\_gpu\_cells}) change only when and where a job runs, never
an id, a file or a value.

From the console the same driver is reached by \vb{campaign\_spec(path)}, \vb{campaign\_spec\_text(text)},
\vb{campaign\_plan(spec)} and \vb{run\_campaign(spec, output)}.

\section{Campaigns over \texorpdfstring{$\Z$}{Z}}\label{sec:campaign-z}

Under \texttt{coefficient\_ring = Z} a campaign is a grid of \emph{descent cells}, one per $(N,n)$, each writing one
\texttt{integral\_descent\_record/v1} (Chapter~\ref{ch:integral}). The grid keys are \vb{n\_sym} (required) and \vb{level}
(default~1; a spec error under any other ring); \vb{p} optionally lists the Hecke primes (none may divide a level of its
section; default as for the verb). \vb{ell}, \vb{r}, \vb{q}, \vb{orbifold}, \vb{ext = true}, \vb{jobs}${}>1$ and a non-prime
$p$ are refused by name before any cell runs. The spec of the examples below:
\VerbatimInput[fontsize=\small,frame=leftline,framesep=3mm,rulecolor=\color{black!40}]{cli_manual_runs/descent_small.spec}
\cli{--dry-run} lists the cells with their size $D=\psi(N)(n+1)$, the dimension of the induced coefficient module:
\runoutput{campaign_z_dry}
A cell id is \texttt{desc\_N<N>\_n<n>[\_p<p1>-<p2>-\dots]\_s<16 hex>}: \texttt{coefficient\_ring = Z} is a non-default
mathematical setting, so its fingerprint is part of every id. The run writes \file{cells/<id>.json} (one line,
\texttt{integral\_descent\_cell/v1}, with the record and a completion trailer), \file{summary.tsv} (\texttt{id N n primes
status H0 H1 H2 H3 record\_fnv1a64 seconds}) and \file{index.json}, and ends with the line \texttt{EFFECTIVE campaign
arm=descent cells= ran= resumed= refused= failed=}:
\runoutput{campaign_z_run}
A refused record (a precision cap or a memory prediction) is written with \texttt{status = refused} and counted; a cell
that throws is \texttt{failed} (exit~2); \vb{resume = true} skips a cell file that ends in the trailer. From the console,
\vb{campaign\_plan(spec)} of such a spec returns one group with the cell ids and \vb{run\_campaign(spec, output)} runs it.
Details: \file{docs/matrix-engine/campaign.md}~\S17.

\section{The orbifold census at \texorpdfstring{$\ell=2,3$}{l=2,3}}\label{sec:orbifold}

The pipeline of Chapters~\ref{ch:field}--\ref{ch:ext} computes \emph{Bredon} objects at $\ell=2,3$
(Section~\ref{sec:model}). To see the ordinary
group cohomology of $\SL_2(\Z)$ at those primes a campaign has a second mode, switched on by the one key
\vb{orbifold = c\_i} (``the isotropy model $C_I$''). Without the key nothing in this section is consulted and every
output is unchanged. The mode writes one record per \emph{cell} $(n,\ell)$ instead of one per Hecke
job. Details: \file{docs/matrix-engine/campaign.md}~\S16 and \file{docs/implementation/slot22/L23\_1\_report.md}.

\paragraph{What a cell holds.} The record schema is \texttt{orbifold\_cell/v2}; its hash is a function of the mathematics
only (timings are unhashed fields).\footnote{\texttt{v1} was the same record with the sum-rule oracle's wall time inside the
hashed gate detail; \texttt{v2} moves that time to an unhashed field. History:
\file{docs/matrix-engine/campaign.md}~\S16.} For one cell it contains
\begin{itemize}[nosep]
\item $H^q(\SL_2(\Z);\Sym^n\F_\ell^2)$ for $q\le3$ with the Hecke operators $T_p$ on each, the Jordan type of $T_\ell$
  and the cup product with the Euler class $e$ (computed over $\F_\ell$ on $C_I$);
\item the torsion of $H^q(\SL_2(\Z);\Sym^n\Z^2)$ (Smith invariants over $\Z/\ell^K$);
\item the $\ell$-adic lattice of $H^1(\SL_2(\Z);\Sym^n\Z^2)$: its rank, its torsion and the elementary divisors of
  $T_p-(1+p)$, $T_p-(1+p^{n+1})$ and $T_\ell$;
\item a cross-check against a second model of the same group, the amalgam $\Z/4*_{\Z/2}\Z/6$ (no model code shared with $C_I$; run for $\ell=2$, $n\ge1$ and for $\ell=3$, even $n\ge2$);
\item the \emph{gates}: every statement the cell checks, each tagged with where it comes from. A \texttt{G} gate is a
  theorem (of [C], of the review, or of the literature) or a consistency check between two readings of the same group
  (cross-model, UCT, the reader's self-check); a failure is a bug or a refutation; a \texttt{G-conj}
  gate tests a conjecture, and a mismatch is only recorded as a \emph{finding}.
\end{itemize}
A divisor is printed as \texttt{0} (a unit), \texttt{inf} (the known kernel, the Eisenstein line) or \texttt{>=<cap>}
(not certified below the precision cap). The record also lists the precision path the $\ell$-adic reader took
(\texttt{reader\_path}) and the provenance of the binary (commit, path, \texttt{sha256}); neither enters the
hash \texttt{record\_fnv1a64}.

\paragraph{The spec.} A section gives only \vb{n\_sym} and \vb{ell} (with $\ell\in\{2,3\}$); everything else is
optional:
\begin{center}\small
\begin{tabular}{@{}l>{\raggedright\arraybackslash}p{0.62\linewidth}@{}}
\toprule
\textbf{Key} & \textbf{Meaning}\\\midrule
\vb{orbifold = c\_i} & switch the campaign to orbifold cells; \texttt{off} or absent: Hecke jobs, unchanged\\
\vb{orbifold\_primes} & the Hecke primes $p$ whose $T_p$ a cell computes; default \texttt{2,3,5,7}; $T_\ell$ is always included\\
\vb{orbifold\_m\_target} & the $\ell$-adic precision every cell must reach; default \texttt{auto}, which is
  $\max(\lfloor (n+1)/(\ell+1)\rfloor+10,\,40)$\\
\vb{orbifold\_k\_cap} & the largest working precision the reader may climb to; default \texttt{0}, meaning its whole built range (1024 at $\ell=2$, 312 at $\ell=3$)\\
\bottomrule
\end{tabular}
\end{center}
\begin{spec}
name     = l23_first
orbifold = c_i
n_sym    = 0..700
ell      = 2
---
n_sym    = 0..700 step 2
ell      = 3
\end{spec}
Run \cli{--campaign l23.spec --dry-run} first. \emph{Refused before any cell runs} (exit~1, no output): $\ell\notin\{2,3\}$,
a \vb{p}, \vb{r} or \vb{q} axis, \vb{ext = true}, a non-default \vb{coefficient\_ring}, \vb{jobs}${}>1$ (shard by
writing several specs instead), a cell whose \texttt{m\_target}${}+16$ exceeds the reader's range, an unknown model, a
non-prime in \vb{orbifold\_primes}. The ordinary Hecke runner refuses an orbifold spec, so a Hecke grid is never run by
accident.

\paragraph{Cell ids and output.} A cell is named \texttt{orb\_ell<L>\_n<N>}. A non-default \vb{orbifold\_primes},
\vb{orbifold\_m\_target} or \vb{orbifold\_k\_cap} (each can change a recorded value) appends a settings suffix
\texttt{\_s<16 hex>} to the id, as for the Hecke jobs (Section~\ref{sec:campaign-out}). The output directory holds
\file{cells/<id>.json}, \file{summary.tsv} (one line per cell, 42 columns), \file{layers.tsv} (one line per degree $q$,
layer and $T_\ell$-eigenvalue), and \file{index.json}.

\paragraph{Reading the \texttt{EFFECTIVE} line.} The run starts with
\texttt{EFFECTIVE campaign arm=orbifold model=c\_i cells=<k> primes=\dots{} m\_target=\dots{} fingerprint=\dots} and ends with
\begin{shell}
EFFECTIVE campaign arm=orbifold cells=<k> ran=<r> skipped=<s> failed=<f> refused=0 fallback:<b> findings=<c> stopped=<0|1> wall_s=... peak_rss_mb=...
\end{shell}
Read it before trusting a shard. A clean census shard has \texttt{failed=0}, \texttt{refused=0}, \texttt{fallback:0} and
\texttt{stopped=0}. \texttt{ran} are the cells computed now, \texttt{skipped} those taken from an earlier run
(\vb{resume}). \texttt{fallback:<b>} counts the cells whose status is \texttt{fallback}: a divisor that should be finite
stayed \texttt{>=cap}, or the reader hit \vb{orbifold\_k\_cap} unresolved; it is recorded and counted, not a failure,
and a census shard must read \texttt{fallback:0}. \texttt{findings} counts failed \texttt{G-conj} gates.
In the calibration runs (cells $(2,240)$, $(2,500)$, $(2,700)$, $(3,700)$, $(2,701)$) every cell read
\texttt{ran=1 failed=0 refused=0 fallback:0 findings=0 stopped=0}.

\paragraph{Resume, and when a campaign stops.} With \vb{resume = true} a cell is skipped iff its file ends with the
trailer \texttt{"orbifold\_cell\_complete":true} and its status is \texttt{ok} or \texttt{fallback}. A complete record
whose status is \texttt{failed} \emph{stops the campaign there}: the program prints
\texttt{STOP (resume on a failed cell; [P3] R7): <id>: <first failure> -- remove or rename cells/<id>.json to re-run it},
nothing after that cell runs, \texttt{stopped=1} and \texttt{failed=1} appear in the \texttt{EFFECTIVE} line, and the
exit code is~2. A failed \texttt{G} gate means a bug or a refutation, so it is to be reported, not re-run until it passes.
A cell killed while writing leaves no trailer and is simply recomputed on the next start.

\paragraph{Memory and time.} The orbifold driver has no memory limit of its own, deliberately: a hard limit inside the
process would turn a failed allocation into a false \texttt{failed} cell. Limit memory from outside, one shard at a time.
The calibration measured, with one cell at a time (peak resident memory, wall time): $(2,240)$ 1.7\,GiB, 2.7\,min;
$(2,500)$ 6.8\,GiB, 25\,min; $(2,700)$ 13.2\,GiB, 63\,min; $(3,700)$ 12.2\,GiB, 32\,min; the odd cell $(2,701)$
13.5\,GiB, 74\,min. Memory grows like $n^{1.95}$, time like $n^{2.7}$ to $n^{3.1}$. For many cells at once, plan
$13.5\,\mathrm{GiB}\cdot(n/701)^{1.95}$ per cell and keep the cells with $n\ge600$ to a few at a time
(\file{docs/implementation/slot22/L23\_4\_calib\_report.md}).

\paragraph{A small run.} The spec \file{docs/latex/cli_manual_runs/orbifold_small.spec} asks for the cells $(n,\ell)=(10,2)$,
$(11,2)$, $(10,3)$:
\runoutput{orbifold_dry}
\runoutput{orbifold_run}
Each cell line prints the dimensions $h^q$ over $\Fl$, the rank of the lattice, the type of $\tors H^1$, the Jordan type of
$T_\ell$ on $H^1$, the precisions $K$, $m$ and the gate counts; the \texttt{EFFECTIVE} lines name the routes in force (the
L23 levers of Section~\ref{sec:env-l23}); the summary has 42 columns, of which the first 16 are shown.

\section{Long census runs}\label{sec:localcensus}\label{sec:second-machine}

The large census cells (the campaign specs with \vb{ext = true}, and the saturation-defect cases) need no other front end
than the ones above: a cell is a \cli{--campaign} spec (Section~\ref{sec:campaign}), and a saturation-defect case is a
run of \vb{HeckeSaturationDefectValidation} with its cell list in the environment (Section~\ref{sec:env-ln23}). Batch
drivers that schedule many such cells over one or more machines are site-specific and are not part of this release; what
any such driver should do is the following.
\begin{itemize}[nosep]
\item \textbf{One build, one roster.} Build the binary once (CPU-only is enough for every census value; the CUDA layer
  changes no value), record its \texttt{sha256} and its \cli{--list} count, and refuse to resume a run with a different
  binary. Records computed by two machines are comparable when the binaries come from the same source commit; recompute a
  few cells already known (\emph{anchor} cells) on a new machine first and compare their records byte for byte before
  the bulk run.
\item \textbf{Predict, then cap memory.} Predict every cell's peak resident memory before it starts (from the
  calibrations of Section~\ref{sec:orbifold}, or the \texttt{predicted\_bytes} line where a validation prints one); keep
  the sum of the predictions of the running cells below the free memory, and limit memory from outside the process,
  never inside it.
\item \textbf{Threads.} \env{ENGINE_THREADS} per cell; one thread per cell for the saturation-defect cases.
\item \textbf{Resume, never re-run blindly.} A finished cell keeps its record; a resumed run starts only the unfinished
  ones. A cell that fails twice with the same symptom is a finding, not a reason to change its parameters.
\item \textbf{Read the \texttt{EFFECTIVE} lines} of every cell (Section~\ref{sec:validations}) before trusting a run:
  they show the cells, the device route and any fallback.
\end{itemize}
For the saturation-defect target the per-cell environment is
\texttt{ENGINE\_THREADS=1 HECKE\_SAT\_DEFECT\_CELLS=n:ell:kmax HECKE\_SAT\_DEFECT\_STRICT=1}.
\texttt{HECKE\_SAT\_DEFECT\_STRICT=1} makes the validation fail every block whose reading was cut short by \texttt{kmax}
(\texttt{not\_stable\_at\_KMAX}) as well, so a too-small \texttt{kmax} is loud; without it a cell fails only when it has no
\texttt{ok} block and some block is not a torsion refusal.

%======================================================================
\chapter{Performance and engine switches}\label{ch:perf}
%======================================================================

Every value the program prints is a function of the mathematics and the mathematical settings only. How fast and on which
hardware it is produced is governed separately: by the build type (Section~\ref{sec:build}), the engine settings of the
console (Section~\ref{sec:settings}) and the environment variables of this chapter. The rule that makes this safe:
\textbf{no environment variable of Sections~\ref{sec:env-engine}--\ref{sec:env-l23} can change a value}; each is gated by a
validation that compares the routes entry by entry, and most print an \texttt{EFFECTIVE} line naming the route in force.

\section{Benchmarks: \texttt{--bench}}\label{sec:bench}

\cli{--bench} drives the shipped pipeline through its public API -- the same code paths as the
validations -- for one point, and prints phase-split timings, peak RSS and one machine-readable line
\begin{center}
\texttt{BENCH mode=\dots{} n=\dots{} ell=\dots{} [p=\dots] [r=\dots] total\_s=\dots{} rss\_mb=\dots{} <mode-specific keys>}
\end{center}
(\texttt{rss\_mb} is the peak resident set of the whole process). The usage text:

% GENERATED by docs/latex/tools/gen_cli_manual_tables.sh, i.e.
% `CLI_MANUAL_EMIT=<dir> hacs CliManualCoverageValidation`, from the
% registries of the build it runs in.  Do not edit by hand: regenerate when a
% verb, type, setting, validation or command-line token changes.
\begin{clihelp}
Usage: hacs --bench <mode> <args...> [--repeat k] [--cross-checks]

Modes:
  cohom  <n_sym> <ell>            EquivariantCohomologyOnK on the SL(2,Z) model, all q
  hecke  <n_sym> <ell> <p>        T_p end to end (cosets -> tau -> Theta_g -> A^(q))
  family <n_sym> <ell> <p> <r>    T_{p^r} by recursion AND by the direct double-coset route
  decomp <n_sym> <ell> <p>        + primary module decomposition of the largest H^q
  ext    <n_sym> <ell> [p,q,...]  + the extension layer, Stages 1-5 (default family 2,3,5)
  micro  <op> <args...>           Phase-2 arithmetic micro-benchmarks; `micro --help`
                                  lists the ops (gemm/gemv/rref/kernel/transpose/suite)

Options:
  --repeat k          run the whole pipeline k times; print min/median of the totals
  --backend b         (mode `micro`) legacy | engine; default engine
                      (`legacy` = the pre-Phase-2 substrate; ENGINE_BACKEND
                      sets the default)
  --cross-checks      (mode `hecke`) additionally run the alpha/beta/omega route
                      (off by default: it roughly doubles the measured time)
  --lean              memory-lean configuration: drop the cohomology objects no
                      production consumer reads and skip the Theta_eq certification
                      (A^(q) unchanged; reported as lean=1 on the BENCH line)
  --count-allocs      report heap allocations (count and requested bytes) made
                      during each measured repeat, via a counting replacement
                      global operator new; adds allocs=/alloc_bytes= to the
                      BENCH line. With -DENGINE_COUNT_OBJECTS=ON it also reports
                      FlatMatrix / BlockFlatMatrix constructions.

Every run ends with one machine-readable line:
  BENCH mode=... n=... ell=... [p=...] [r=...] total_s=... rss_mb=... <mode-specific keys>
(rss_mb is peak RSS over the whole process, from /proc/self/status VmHWM.)

NOTE: on this model H^*(K, R; Sym^n) vanishes for ODD n (the -I / odd-n base-model
gap); use an even n_sym for a meaningful measurement.
\end{clihelp}


All modes and options the parser accepts, including the flags that are deliberately kept out of the
frozen usage text:

% GENERATED by docs/latex/tools/gen_cli_manual_tables.sh, i.e.
% `CLI_MANUAL_EMIT=<dir> hacs CliManualCoverageValidation`, from the
% registries of the build it runs in.  Do not edit by hand: regenerate when a
% verb, type, setting, validation or command-line token changes.
\begin{tokentable}{Token}{Arguments}
\clitok{cohom}\index{cohom@\texttt{cohom} (\texttt{-{}-{}bench})} & \texttt{\textless{}n\_sym\textgreater{} \textless{}ell\textgreater{}} & EquivariantCohomologyOnK on the SL(2,Z) model, all degrees. \\
\clitok{hecke}\index{hecke@\texttt{hecke} (\texttt{-{}-{}bench})} & \texttt{\textless{}n\_sym\textgreater{} \textless{}ell\textgreater{} \textless{}p\textgreater{}} & T\_p end to end: cosets, tau, Theta\_g, A\textasciicircum{}(q). \\
\clitok{family}\index{family@\texttt{family} (\texttt{-{}-{}bench})} & \texttt{\textless{}n\_sym\textgreater{} \textless{}ell\textgreater{} \textless{}p\textgreater{} \textless{}r\textgreater{}} & T\_\{p\textasciicircum{}r\} by the algebra recursion and by the independent direct double-coset route. \\
\clitok{decomp}\index{decomp@\texttt{decomp} (\texttt{-{}-{}bench})} & \texttt{\textless{}n\_sym\textgreater{} \textless{}ell\textgreater{} \textless{}p\textgreater{}} & \texttt{hecke} plus the primary module decomposition of the largest H\textasciicircum{}q. \\
\clitok{ext}\index{ext@\texttt{ext} (\texttt{-{}-{}bench})} & \texttt{\textless{}n\_sym\textgreater{} \textless{}ell\textgreater{} [p,q,...]} & \texttt{hecke} plus the extension layer, Stages 1-5; the family defaults to 2,3,5 (2,3,5,7 under \texttt{-{}-{}campaign-{}cell}). \\
\clitok{micro}\index{micro@\texttt{micro} (\texttt{-{}-{}bench})} & \texttt{\textless{}op\textgreater{} \textless{}args...\textgreater{}} & Phase-2 arithmetic micro-benchmarks (the \texttt{micro} table). \\
\clitok{-{}-{}repeat}\index{repeat@\texttt{-{}-{}repeat} (\texttt{-{}-{}bench})} & \texttt{\textless{}k\textgreater{}} & Runs the whole pipeline k \textgreater{}= 1 times and prints min / median of the totals. \\
\clitok{-{}-{}backend}\index{backend@\texttt{-{}-{}backend} (\texttt{-{}-{}bench})} & \texttt{legacy\textbar{}engine\textbar{}gpu} & The substrate of \texttt{micro} (default \texttt{engine}; \texttt{ENGINE\_BACKEND} sets the default). \\
\clitok{-{}-{}cross-{}checks}\index{cross-checks@\texttt{-{}-{}cross-{}checks} (\texttt{-{}-{}bench})}, \clitok{-{}-{}no-{}cross-{}checks}\index{no-cross-checks@\texttt{-{}-{}no-{}cross-{}checks} (\texttt{-{}-{}bench})} & -- & (\texttt{hecke}) also run, or explicitly skip, the classical alpha/beta/omega route; off by default (it roughly doubles the time). \\
\clitok{-{}-{}lean}\index{lean@\texttt{-{}-{}lean} (\texttt{-{}-{}bench})} & -- & Memory-lean configuration: drop objects no production consumer reads and skip the Theta\_eq certification; A\textasciicircum{}(q) unchanged, \texttt{lean=1} on the BENCH line. \\
\clitok{-{}-{}count-{}allocs}\index{count-allocs@\texttt{-{}-{}count-{}allocs} (\texttt{-{}-{}bench})} & -- & Reports the heap allocations of each measured repeat (\texttt{allocs=}, \texttt{alloc\_bytes=} on the BENCH line). \\
\clitok{-{}-{}campaign-{}cell}\index{campaign-cell@\texttt{-{}-{}campaign-{}cell} (\texttt{-{}-{}bench})} & -- & (\texttt{ext}) the campaign cell's own driver options: cone and sign split built, default family 2,3,5,7; not in the frozen help text. \\
\clitok{-{}-{}u-{}ell}\index{u-ell@\texttt{-{}-{}u-{}ell} (\texttt{-{}-{}bench})} & \texttt{auto\textbar{}off\textbar{}budget\_mb=\textless{}M\textgreater{}} & (\texttt{ext}) scopes the MATHEMATICAL setting \texttt{u\_ell} for the run: it decides whether U\_ell enters the record, so it can change a value; not in the frozen help text. \\
\clitok{-{}-{}help}\index{help@\texttt{-{}-{}help} (\texttt{-{}-{}bench})}, \clitok{-{}h}\index{h@\texttt{-{}h} (\texttt{-{}-{}bench})} & -- & Prints the bench usage (after \texttt{micro}: the micro usage). \\
\end{tokentable}


\paragraph{Value-preserving and value-changing flags.} \cli{--repeat}, \cli{--lean},
\cli{--count-allocs}, \cli{--cross-checks}, \cli{--backend} and the three lever pairs never change a
computed value; equal checksums across them are part of what the validations assert. The one exception
is \cli{--u-ell}, which sets the \emph{mathematical} setting \vb{u\_ell} (Section~\ref{sec:settings})
for an \texttt{ext} run and so decides whether $U_\ell$ enters the extension record.

\paragraph{The \texttt{--backend} values.}
% GENERATED by docs/latex/tools/gen_cli_manual_tables.sh, i.e.
% `CLI_MANUAL_EMIT=<dir> hacs CliManualCoverageValidation`, from the
% registries of the build it runs in.  Do not edit by hand: regenerate when a
% verb, type, setting, validation or command-line token changes.
\begin{tokentable}{Token}{Arguments}
\clitok{engine}\index{engine@\texttt{engine} (\texttt{-{}-{}backend})} & -- & The Phase-2 ModCtx / DenseBuf / Kernels layer; the default and the shipped substrate. \\
\clitok{legacy}\index{legacy@\texttt{legacy} (\texttt{-{}-{}backend})} & -- & The pre-Phase-2 substrate end to end (an engine::elim::ScopedLegacy guard); equal checksums against \texttt{engine} are a bit-identity statement. \\
\clitok{gpu}\index{gpu@\texttt{gpu} (\texttt{-{}-{}backend})} & -- & The Phase-5 device kernels (gemm / gemv / transpose); rref / kernel and every op without a usable device report \texttt{not available}, exit 2. \\
\end{tokentable}



A run on the reference machine:
\runoutput{bench_hecke}

\section{Micro-benchmarks}\label{sec:micro}

\cli{--bench micro <op> <args...>} measures single arithmetic kernels on fixed-seed inputs; the
reported FNV-1a-64 checksum of the result is a bit-identity check between backends.

% GENERATED by docs/latex/tools/gen_cli_manual_tables.sh, i.e.
% `CLI_MANUAL_EMIT=<dir> hacs CliManualCoverageValidation`, from the
% registries of the build it runs in.  Do not edit by hand: regenerate when a
% verb, type, setting, validation or command-line token changes.
\begin{clihelp}
Usage: hacs --bench micro <op> <args...> [--repeat k] [--backend legacy|engine|gpu]

Ops (every <ell> defaults to 134217757, the benchmark prime of doc 03):
  gemm      <n> [<k>] [<ell>]                 C = A*B, A: n x k, B: k x n; work = n*n*k
  gemv      <n> [<ell>]                       y = A*x, A: n x n; work = n*n
  rref      <rows> <cols> [<rank>] [<ell>]    RREF of a matrix of prescribed rank
  kernel    <rows> <cols> [<rank>] [<ell>]    kernel() of the same matrix family
  transpose <n> [<ell>]                       transpose of an n x n matrix
  suite                                       the Phase-2 exit-criterion set + floors
  cuda                                        the Phase-5 GPU throughput table (needs a device)

<rank> omitted, 0, or >= min(rows,cols) means FULL rank (a plain random matrix,
which is what doc 03 Sec. 7.3's reference elimination uses); a smaller <rank>
builds L(rows x rank) * R(rank x cols) so rank-deficient inputs are covered.

All inputs are fixed-seed (splitmix64 over the shape), so two runs -- and the two
backends -- see byte-identical inputs; the reported FNV-1a-64 checksum of the
result therefore doubles as a bit-identity check between backends.

--backend engine selects the Phase-2 ModCtx/DenseBuf/Kernels layer (the default).
--backend gpu selects the Phase-5 engine::cuda device kernels (gemm / gemv /
transpose; rref / kernel stay CPU-driven and report "not available", exit 2).
Without a usable CUDA device every gpu op reports "not available", exit 2.
\end{clihelp}

% GENERATED by docs/latex/tools/gen_cli_manual_tables.sh, i.e.
% `CLI_MANUAL_EMIT=<dir> hacs CliManualCoverageValidation`, from the
% registries of the build it runs in.  Do not edit by hand: regenerate when a
% verb, type, setting, validation or command-line token changes.
\begin{tokentable}{Token}{Arguments}
\clitok{gemm}\index{gemm@\texttt{gemm} (\texttt{-{}-{}bench micro})} & \texttt{\textless{}n\textgreater{} [\textless{}k\textgreater{}] [\textless{}ell\textgreater{}]} & C = A*B with A: n x k, B: k x n; work n*n*k. \\
\clitok{gemv}\index{gemv@\texttt{gemv} (\texttt{-{}-{}bench micro})} & \texttt{\textless{}n\textgreater{} [\textless{}ell\textgreater{}]} & y = A*x with A: n x n. \\
\clitok{rref}\index{rref@\texttt{rref} (\texttt{-{}-{}bench micro})} & \texttt{\textless{}rows\textgreater{} \textless{}cols\textgreater{} [\textless{}rank\textgreater{}] [\textless{}ell\textgreater{}]} & RREF of a fixed-seed matrix of prescribed rank (omitted, 0 or \textgreater{}= min(rows, cols): full rank). \\
\clitok{kernel}\index{kernel@\texttt{kernel} (\texttt{-{}-{}bench micro})} & \texttt{\textless{}rows\textgreater{} \textless{}cols\textgreater{} [\textless{}rank\textgreater{}] [\textless{}ell\textgreater{}]} & kernel() of the same matrix family. \\
\clitok{transpose}\index{transpose@\texttt{transpose} (\texttt{-{}-{}bench micro})} & \texttt{\textless{}n\textgreater{} [\textless{}ell\textgreater{}]} & Transpose of an n x n matrix. \\
\clitok{suite}\index{suite@\texttt{suite} (\texttt{-{}-{}bench micro})} & -- & The Phase-2 exit-criterion set with its floors. \\
\clitok{cuda}\index{cuda@\texttt{cuda} (\texttt{-{}-{}bench micro})} & -- & The Phase-5 GPU throughput table (needs a device). \\
\clitok{-{}-{}help}\index{help@\texttt{-{}-{}help} (\texttt{-{}-{}bench micro})}, \clitok{-{}h}\index{h@\texttt{-{}h} (\texttt{-{}-{}bench micro})}, \clitok{help}\index{help@\texttt{help} (\texttt{-{}-{}bench micro})} & -- & Prints the micro usage. \\
\end{tokentable}


\paragraph{Reference figures.} The standing performance gates (\texttt{--bench hecke 256 134217757 2},
\texttt{--bench decomp 128 134217757 2}) and the measured tables live in
\file{docs/performance-reference.md}; timings depend on the machine and its load and are never part of
a result.

\paragraph{The standing gates.} Two points guard the pipeline's speed from slot to slot: \cli{--bench hecke 256 134217757 2}
must stay at or below 0.470\,s and \cli{--bench decomp 128 134217757 2} at or below 0.109\,s on the reference machine, and the
median of 13 alternating rounds against the previous slot's binary must not be slower by more than 4\,\%
(\file{docs/matrix-engine/tools/standing_gate_h2h.sh}); the fitted exponent of the total time in $n$ over $\{64,96,128\}$ stays
$\le3.2$. Compare two binaries only when they place the global thread pool alike
(\file{docs/matrix-engine/tools/pool_placement_check.sh}).

\section{Environment variables: overview}\label{sec:env}

Environment variables fall into five groups: the engine switches (this section to Section~\ref{sec:env-l23}), which can
never change a value; the owner-gated emitters, which write files; the interlocks and the strict acceptance mode; and the
read-only hooks of individual validations. Unless a variable of Section~\ref{sec:env-gated} is set, \cli{--all} and
\cli{--all-full} are read-only. Three conventions hold throughout:
\begin{itemize}[nosep]
\item A \emph{rollback arm} is a variable whose value \texttt{0} restores an older route verbatim (the new route is the
  default); any other value prints a \texttt{WARNING} and keeps the default. It exists to compare and to recover, not to
  tune.
\item An \emph{opt-in route} is off unless set to exactly \texttt{1}.
\item Most variables are read once per process; set them on the command line of the run.
\end{itemize}

\section{Threads, substrate and the tape engine}\label{sec:env-engine}

{\small
\begin{longtable}{@{}>{\raggedright\arraybackslash}p{0.33\linewidth}>{\raggedright\arraybackslash}p{0.61\linewidth}@{}}
\toprule
\textbf{Variable} & \textbf{Effect}\\\midrule\endhead
\env{ENGINE_THREADS} & worker threads of the engine pool (console: \vb{threads}); the descent's Smith forms use it too (an integer in $[1,256]$; default the CPU affinity count, at most 64)\\
\env{ENGINE_HOST_THREADS} & the host-side fan-out of independent tasks (one thread per task); \texttt{0} runs them serially (env-only)\\
\env{ENGINE_BACKEND} & \texttt{engine} (default) or \texttt{legacy}, the pre-engine substrate (console: \vb{backend})\\
\env{ENGINE_BATCH}, \env{ENGINE_KSPLIT}, \env{ENGINE_ARENA}, \env{ENGINE_CSE}, \env{ENGINE_FUSE}, \env{ENGINE_LAYOUT} & the tape compiler's passes (console: \vb{batch}, \vb{ksplit}, \vb{arena}, \vb{cse}, \vb{fuse}, \vb{layout})\\
\env{ENGINE_TILE_MULT}, \env{ENGINE_TILE_COLS}, \env{ENGINE_SPIN}, \env{ENGINE_ELIM_BASE}, \env{ENGINE_ELIM_PANEL}, \env{ENGINE_ELIM_TRSM_BASE} & tuning of the tiles, the spin-wait and the blocked elimination (panel, base case, triangular-solve base), read once at start-up (env-only settings)\\
\env{ENGINE_HOST_ELIM_HOOK} & \texttt{0}: rollback arm; the cone and $\iota^*$ certificates no longer bind the elimination's inner parallel hook to the pool\\
\env{ENGINE_LADIC_WIDTH} & \texttt{32}, \texttt{64} or \texttt{128}: the digit width of the $\ell$-adic route for $\ell^k\ge2^{64}$; the records are width-independent\\
\env{ENGINE_BIGNTT_MIN_LIMBS} & a positive integer (default \texttt{1024}): the size in 32-bit limbs from which the untrusted big-integer factor finder multiplies modulo $n$ by an NTT product with Barrett's reduction instead of BigInt's Karatsuba and Knuth-D route; a route knob, it cannot change a value (slot~27 BIGNTT)\\
\bottomrule
\end{longtable}
}
The console's engine settings of the same names override the environment for a session; the \emph{env-only} ones cannot
be changed after start-up (\texttt{set tile\_mult = "2"} is refused with a message naming \texttt{ENGINE\_TILE\_MULT}).

\section{CPU kernels and native small fields}\label{sec:env-cpu}

The matrix product over $\F_p$ (``GEMM'') is the inner loop of almost everything. It has several exact implementations,
chosen at run time from what the CPU offers; all give the same residues.
{\small
\begin{longtable}{@{}>{\raggedright\arraybackslash}p{0.33\linewidth}>{\raggedright\arraybackslash}p{0.61\linewidth}@{}}
\toprule
\textbf{Variable} & \textbf{Effect}\\\midrule\endhead
\env{ENGINE_KERNEL} & \texttt{auto} (default: AVX-512 if the CPU has it, else AVX2, else portable), \texttt{avx512}, \texttt{avx2} or \texttt{portable}: the 32-bit GEMM microkernel. \texttt{avx2} on an AVX-512 machine imitates a CPU without AVX-512 (e.g.\ Zen~3). An unavailable value warns and stays \texttt{auto}; a set value prints \texttt{EFFECTIVE ENGINE\_KERNEL path= value= cpu=avx2:<0|1>,avx512:<0|1>}\\
\env{ENGINE_AVX2_PAIR} & \texttt{0} (or \texttt{off}): rollback arm of the AVX2 pair route (16-bit products, $p\le4096$); the AVX2 path then runs its 64-bit microkernel for every $p$\\
\env{ENGINE_VNNI}, \env{ENGINE_VNNI_SHADOW} & the AVX-512 VNNI route for $p\le128$ and its two-limb form for $128<p\le2^{14}$, on by default where the CPU has \texttt{avx512\_vnni}; \texttt{ENGINE\_VNNI=0} (or \texttt{off}, the former rollback arm) is retired (S26 LEVERS-26): refused in a default build, a comparison arm in a legacy build; \texttt{ENGINE\_VNNI\_SHADOW=1} recomputes every routed product on the 32-bit kernel and stops on a difference\\
\env{ENGINE_VNNI2_FLOOR} & the two-limb VNNI route's work-per-thread floor ($mnk\ge$ floor $\times$ threads, default $2^{20}$ (S26 LEVERS-26; was $2^{18}$), threads = the configured pool size); \texttt{0} restores the slot-24 rule; a value that is not a string of digits (or overflows 64 bits) warns and keeps the default; a set value prints \texttt{EFFECTIVE ENGINE\_VNNI2\_FLOOR floor= value= threads=} when the floor is first consulted (under \texttt{auto}, at the first two-limb product that passes the large-$p$ test); it cannot change a value\\
\env{ENGINE_BINARY_F2} & \texttt{1}: opt-in; the dense eliminations over $\F_2$ run bit-packed (Four Russians), Section~\ref{sec:env-slot22}\\
\env{ENGINE_TERNARY_F3} & \texttt{1}: opt-in; the dense eliminations over $\F_3$ run bit-sliced (two bit planes per row), proved entrywise by \vb{TernaryF3Validation}\\
\bottomrule
\end{longtable}
}
Example: the AVX2 kernel without the pair route on the reference machine (an AVX-512 CPU):
\runoutput{env_kernel}

\section{The device: \texttt{ENGINE\_DEVICE} and the CUDA layer}\label{sec:env-device}

The default device policy is \texttt{cpu}: the golden path, which never creates a CUDA context. A GPU is used only when
asked for. The reference machine's card is (temporarily) an NVIDIA RTX PRO 6000 Blackwell with 96\,GB (compute capability
12.0, like the RTX 5070 Ti it replaces).
{\small
\begin{longtable}{@{}>{\raggedright\arraybackslash}p{0.33\linewidth}>{\raggedright\arraybackslash}p{0.61\linewidth}@{}}
\toprule
\textbf{Variable} & \textbf{Effect}\\\midrule\endhead
\env{ENGINE_DEVICE} & \texttt{cpu} (default, the golden path), \texttt{gpu}, \texttt{hybrid} (console: \vb{device}). On the Hecke pipeline's shapes the GPU is slower than the CPU and stays opt-in; for the integral descent's 2-adic Smith form, \texttt{gpu} is the fast route (Section~\ref{sec:env-descent})\\
\env{ENGINE_CUDA_DISABLE}, \env{ENGINE_CUDA_DEVICE}, \env{ENGINE_CUDA_PROBE} & disable the CUDA layer, pick a device, probe; \texttt{ENGINE\_CUDA\_DISABLE=1} is required for sanitizer runs\\
\env{ENGINE_CUDA_GRAPH}, \env{ENGINE_CUDA_MERGE_LEVELS}, \env{ENGINE_CUDA_PROGRAM_CACHE}, \env{ENGINE_ELIM_SESSION} & GPU tuning knobs, only under a GPU device policy. \texttt{ENGINE\_CUDA\_GRAPH}: \texttt{0} never capture a CUDA graph, \texttt{1} (default) capture on the second run of a program, \texttt{2} on the first. \texttt{ENGINE\_CUDA\_MERGE\_LEVELS=0} stops merging consecutive GPU-only wavefronts into one graph. \texttt{ENGINE\_CUDA\_PROGRAM\_CACHE=<k>} sets the number of cached programs (default 64). \texttt{ENGINE\_ELIM\_SESSION=0} turned off the device-resident elimination session; retired at the slot-24 close (CL-R2): refused in a default build (\texttt{OC\_ENABLE\_LEGACY=OFF}), a comparison arm only in a legacy build\\
\env{ENGINE_CUDA_PINNED_MB}, \env{ENGINE_CUDA_PINNED_CHUNK_MB}, \env{ENGINE_CUDA_COPY_THREADS}, \env{ENGINE_CUDA_ARENA_SLACK_MB} & the pinned host staging pool (size, chunk), the threads of the copy pool, and the idle bytes the device arena keeps before trimming\\
\env{ENGINE_GPU_LEAN} & by default a program whose GPU nodes are all memory-bound runs on the host, and the arena keeps a sixteenth of the card as slack; \texttt{0\dots} (the former rollback arm) is retired (S25 LEVERS-25): refused in a default build, the old route in a legacy build\\
\env{ENGINE_GPU_IMMA} & the int8 tensor-core GEMM for $p\le127$, always on; \texttt{0\dots} (the former rollback arm) is retired (CL-R2): refused in a default build, a comparison arm in a legacy build (Section~\ref{sec:env-slot22})\\
\env{ENGINE_GPU_ELIM}, \env{ENGINE_GPU_ELIM_MIN}, \env{ENGINE_GPU_ELIM_VERIFY} & the device-resident elimination for $p\le127$: \texttt{0\dots} retired (CL-R2; refused in a default build), size floor in matrix entries (default $2^{20}$), host cross-check (Section~\ref{sec:env-slot22})\\
\env{ENGINE_GPU_SMITH2K}, \env{ENGINE_GPU_SMITH2K_MIN}, \env{ENGINE_GPU_SMITH2K_VRAM_MB}, \env{ENGINE_GPU_SMITH2K_HOST_MB}, \env{ENGINE_GPU_SMITH2K_VERIFY} & the device 2-adic Smith LU of the descent, under \texttt{ENGINE\_DEVICE=gpu} only (Section~\ref{sec:env-descent}); the switches of this opt-in route (not rollback levers, S26 LEVERS-26: \texttt{0} is the host fallback of the opt-in device route, kept while that route exists; a card decision, unlike the retired \texttt{ENGINE\_GPU\_ELIM}/\texttt{\_IMMA}): host LU under \texttt{gpu} (\texttt{0}), size floor $m\cdot n$ (default $2^{24}$), device budget (default the free memory less 512\,MiB and $1/32$ of the card), the host store of the tiled route per phase (default 4096\,MiB), and a host re-run of every device LU (\texttt{1})\\
\bottomrule
\end{longtable}
}

\section{Native \texorpdfstring{$\F_2$}{F2} arithmetic and the GPU elimination}\label{sec:env-slot22}

Three routes sit beside the tape and block code. All three are \emph{value-preserving}:
the reduced row echelon form, the rank and the pivots of a matrix over a field are unique functions of
the matrix, and each route computes them in exact integer arithmetic, so a record is byte-identical with
the route on or off. None is a mathematical setting, none enters a record's hash or a campaign job id,
and none is on by default on the golden path (\texttt{ENGINE\_DEVICE=cpu}). The contract is
\file{ENGINE\_API.md}, Sections~8.1, 11.8 and~11.9.

\paragraph{Native $\F_2$ arithmetic: \env{ENGINE_BINARY_F2}\texttt{=1}.} Opt-in, off by default, and read
once per process; only the value \texttt{1} switches it on. The in-place dense eliminations over a field
with $p=2$ (\texttt{uint32\_t} or \texttt{uint64\_t}) -- \texttt{rref}, \texttt{rref\_with\_transform},
\texttt{inverse}, \texttt{solve\_multi\_rhs}, \texttt{complement\_basis} and
\texttt{column\_rank\_profile} -- then pack the matrix into one bit per entry and eliminate with the
Four-Russians method (M4RI, strips of eight pivots; AVX-512 or AVX2 row XORs chosen at run time, a
portable body on other CPUs), single-threaded. The tape elimination and the block model are not routed,
and no other modulus is affected. Measured on the reference machine at $N\times N$ with
\texttt{--bench micro rref}: $12\times$ at $N=1024$, $46\times$ at $4096$, $70\times$ at $8192$ (identical
checksums). Gate: \vb{BinaryF2Validation} (\texttt{standard}); \texttt{GOLDEN\_STRICT=1 ENGINE\_BINARY\_F2=1
./hacs GoldenRegressionValidation} is clean.

\paragraph{GPU routes for $p\le127$: \env{ENGINE_GPU_IMMA}, \env{ENGINE_GPU_ELIM}.} Both are reached only
under a GPU device policy (\texttt{ENGINE\_DEVICE=gpu}, or \texttt{hybrid} for the GEMM); the default
\texttt{cpu} never creates a CUDA context, and without a device of compute capability $\ge8.0$ the program
silently keeps the CPU kernels.
\begin{itemize}[nosep]
\item \env{ENGINE_GPU_IMMA}: modular GEMM on the int8 tensor cores for moduli $2\le p\le127$ (balanced
  residues, exact \texttt{int32} accumulation). Always on for the device GEMMs; a value starting with
  \texttt{0} was the rollback arm until the slot-24 close (CL-R2) and is now refused in a default build
  (\texttt{OC\_ENABLE\_LEGACY=OFF}); a legacy build still honours it. Gate: \vb{CudaImmaBitIdentityValidation} (\texttt{standard}).
\item \env{ENGINE_GPU_ELIM}: the whole elimination on the device for a \texttt{uint32\_t} field with
  $p\le127$, replacing the tape elimination and the dense RREFs, and only under
  \texttt{ENGINE\_DEVICE=gpu}. A value starting with \texttt{0} turns it off in a legacy build only; a default build refuses
  it (retired, CL-R2), and \texttt{ENGINE\_DEVICE=cpu} remains the way to stay off the GPU. \env{ENGINE_GPU_ELIM_MIN}\texttt{=<m*n>}
  is a size floor (default $1\,048\,576$ entries; tuning only). \env{ENGINE_GPU_ELIM_VERIFY}\texttt{=1}
  re-runs each device elimination on the host and compares the result entrywise (a difference raises an
  error); \env{ENGINE_TRACE_IMMA_ELIM}\texttt{=1} prints the phase times of each call. Gate:
  \vb{CudaImmaElimValidation} (\texttt{standard}). \texttt{Z/ell\^{}k} is not covered.
\end{itemize}
A device elimination that fails before writing output (for instance a memory allocation) falls back to the
host route with the same result. The fallback is never silent: one stderr line
\texttt{[imma\_rref] fallback to the host route (<route>, <m>x<n>): <reason>} per fallback, and one summary line
\begin{shell}
EFFECTIVE DEVICE_ELIM calls=<n> device=<n> fallbacks=<n>
\end{shell}
on stderr when the process exits, if any device elimination was attempted (\texttt{device} counts the
eliminations that ran on the GPU, \texttt{calls} $=$ \texttt{device} $+$ \texttt{fallbacks}). Read it
after a GPU run: \texttt{fallbacks=0} means every elimination ran on the device. A process that never
attempted one prints no such line. \env{ENGINE_GPU_ELIM_SIMULATE_CUDA_ERROR}\texttt{=1} is a test hook that
forces the fallback path.

\section{The integral descent: routes, threads and the GPU}\label{sec:env-descent}

The descent of Chapter~\ref{ch:integral} spends almost all of its time in Smith forms over $\Z/2^k$ and $\Z/3^k$ of large
dense matrices (the 2-adic piece at level $N$ has about $\psi(N)(n+1)\cdot\tfrac43$ rows). Its routes, all value-preserving,
are on by default; each lever below had a rollback arm (\texttt{0}) gated by the validation named in
\file{docs/validation-catalogue.md}. Since S26 LEVERS-26 these arms are \emph{retired}: a default build refuses them
(\texttt{REFUSED <NAME>=0: compiled out (OC\_ENABLE\_LEGACY=OFF)}, exit 1), a legacy build still reads them, and the validations
reach the earlier routes in process. \texttt{HECKE\_DESCENT\_SMITH\_PANEL} and \texttt{HECKE\_DESCENT\_SMITH\_SIMD} are tuning knobs and
accept every documented value in every build.
{\small
\begin{longtable}{@{}>{\raggedright\arraybackslash}p{0.33\linewidth}>{\raggedright\arraybackslash}p{0.61\linewidth}@{}}
\toprule
\textbf{Variable} & \textbf{Effect}\\\midrule\endhead
\env{HECKE_DESCENT_TREE}, \env{HECKE_DESCENT_TREE_PAIR} & the tree models of the 2- and 3-adic pieces and of the pair (closed-form cochains); \texttt{0} restores the dense isotropy model (same groups and maps)\\
\env{HECKE_DESCENT_GOOD_SPARSE}, \env{HECKE_DESCENT_GOOD_DENSE} & the good-prime step at level~1: the sparse cokernel (rollback: the seven-vertex pipeline), read as dense ranks of the averaging idempotents (rollback: the sparse elimination)\\
\env{HECKE_DESCENT_SMITH_MT} & the threaded 2-/3-adic Smith forms on \texttt{ENGINE\_THREADS} threads; \texttt{0} (or \texttt{ENGINE\_THREADS=1}) the serial code. Prints \texttt{EFFECTIVE HECKE\_DESCENT\_SMITH\_MT phase=start|exit \dots} when set\\
\env{HECKE_DESCENT_SMITH_PANEL} & the LU panel width of the threaded Smith form. \textbf{Unset = auto}: the 2-adic host LU takes \textbf{512} columns when the matrix it factors has $\ge$\textbf{12\,000} rows, else 128; the 3-adic LU 128; the device LU always 128. An explicit width in $[1,4096]$ wins for every size (a tuning knob, not a lever); \texttt{128} is the earlier default, \texttt{64} the first threaded one\\
\env{HECKE_DESCENT_SMITH_F2_MT} & the lazy blocked $\F_2$/$\F_3$ rank profile of each Smith level and the vectorised level passes\\
\env{HECKE_DESCENT_SMITH_BLOCKED}, \env{HECKE_DESCENT_SMITH_SIMD} & the blocked, threaded integer Smith form of the maps and presentations (rollback: the scalar oracle); \texttt{\_SIMD=scalar|avx2|avx512} forces its kernel body\\
\env{HECKE_DESCENT_PAR_ASM}, \env{HECKE_DESCENT_PAR_ELL}, \env{HECKE_DESCENT_PAR_LN}, \env{HECKE_DESCENT_EMBED_MAP} & concurrency of the assembly: the pair's complex assembly, the $\ell\ge5$ pieces, the level-$N$ stages (each under a joint memory prediction), and the cone maps without a dense cochain map\\
\bottomrule
\end{longtable}
}
\paragraph{The GPU route.} Under \texttt{ENGINE\_DEVICE=gpu} the 2-adic Smith LU runs device-resident (tiled when the matrix
does not fit) above the size floor \texttt{ENGINE\_GPU\_SMITH2K\_MIN}; the result is bit-identical, and the run prints
\texttt{EFFECTIVE GPU\_SMITH2K \dots{} device=<calls> resident= tiled= fallbacks=} at exit. On the 96\,GB card every LU of
the largest cells is resident: $\Gamma_0(1000)$, $n=20$ takes 142.8\,s on the GPU against 903.9\,s on 16 CPU threads, with a
predicted peak of about 36.7\,GiB. Below the floor nothing changes; with the floor at~0 even a small cell uses the device:
\runoutput{env_descent_gpu}
And the panel lever in action (an explicit width; the start and exit lines show the width in force and how many LU calls
took each width):
\runoutput{env_descent_panel}

\section{The Hecke module and forms routes}\label{sec:env-module}

The $U_\ell$ operator, the $\Z/\ell^k$ module route (saturation and Loewy cells of the census), and the forms algebra each
have memory and speed levers. All are value-preserving. Where a lever is still a rollback arm, \texttt{0} restores the earlier route;
a \emph{retired} arm is refused in a default build (\texttt{REFUSED <NAME>=0: compiled out (OC\_ENABLE\_LEGACY=OFF)}, exit 1) and
works as before in a legacy build (\texttt{-DOC\_ENABLE\_LEGACY=ON}); the earlier routes themselves stay, as the comparisons of the validations.
{\small
\begin{longtable}{@{}>{\raggedright\arraybackslash}p{0.33\linewidth}>{\raggedright\arraybackslash}p{0.61\linewidth}@{}}
\toprule
\textbf{Variable} & \textbf{Effect}\\\midrule\endhead
\env{HECKE_U_ELL_LOW_RANK}, \env{HECKE_COSET_STREAM}, \env{HECKE_PARALLEL_SETUP}, \env{HECKE_RANK1_CLOSED_FORM} & the four memory levers of $U_\ell$: low-rank factorisation, streamed cosets, parallel set-up, the closed-form rank-1 terms; all four rollback arms retired (S24 CL-R1: refused in a default build)\\
\env{HECKE_TAU_WIDE} & the wide tier of $U_\ell$ for $\ell>23170$ (\texttt{0}: such $\ell$ are refused as before)\\
\env{HECKE_RHO2_STREAM} & the cross-check route of an operator streamed key by key (no per-degree memo); \texttt{0} retired (S25 LEVERS-25)\\
\env{HECKE_MODULE_STREAM}, \env{HECKE_MODULE_EVAL}, \env{HECKE_RHO_SIMD}, \env{HECKE_ROUTE2_SHARED_MEMO} & the streamed $\Z/\ell^k$ module route: read-offs at unit pivots, fast $\rho_n$ evaluation (SIMD: \texttt{0|scalar|avx2|avx512}), one memo per coset for the cross-check; \texttt{0} (and \texttt{scalar} for \texttt{HECKE\_RHO\_SIMD}) retired (S25 LEVERS-25)\\
\env{HECKE_MODULE_PARALLEL}, \env{HECKE_MODULE_MT}, \env{HECKE_MODULE_TILE_WORK}, \env{HECKE_MODULE_LEAN_PEAK} & its parallel assembly, threaded decompositions, the minimum work per parallel tile, and the lean peak memory; \texttt{HECKE\_MODULE\_PARALLEL=0} retired (S25 LEVERS-25)\\
\env{HECKE_PRECISION_REUSE} & one module run per cell with the lower precisions derived (\texttt{0}, one run per level: retired, S25 LEVERS-25)\\
\env{HECKE_SAT_PERF} & per-cell memos of the saturation certificates\\
\env{HECKE_FIELD_DIET} & the memory diet of a field cell (dense programs in parts, early release, \texttt{malloc\_trim}); \texttt{0} retired (S26 LEVERS-26)\\
\env{HECKE_FORMS_FASTMUL}, \env{HECKE_FORMS_NTT_CACHE}, \env{HECKE_FORMS_NTT3} & the forms algebra's $q$-series products by number-theoretic transforms, their cache, and the three-prime route for $\ell\ge2^{16}$; \texttt{HECKE\_FORMS\_FASTMUL=0} retired (S25 LEVERS-25)\\
\env{HECKE_LEVEL_N_GEN_ORDER} & \texttt{sorted} (default) or \texttt{as\_built}: the generator order of the level-$N$ cokernel (any other value is an error)\\
\bottomrule
\end{longtable}
}

\section{The orbifold census routes (\texttt{L23\_*})}\label{sec:env-l23}

The orbifold census at $\ell=2,3$ (Section~\ref{sec:orbifold}) prints one \texttt{EFFECTIVE l23\_\dots=on|off} line per
route (see the example there). Each is on by default; \texttt{0} restores the earlier route, except for the two retired arms
(\texttt{L23\_TERNARY\_F3=0}, \texttt{L23\_BIN3O1=0}, S25 LEVERS-25: refused in a default build, the earlier route in a legacy build).
{\small
\begin{longtable}{@{}>{\raggedright\arraybackslash}p{0.33\linewidth}>{\raggedright\arraybackslash}p{0.61\linewidth}@{}}
\toprule
\textbf{Variable} & \textbf{Effect}\\\midrule\endhead
\env{L23_TERNARY_F3} & the degree-1 amalgam reader's $\F_3$ route (bit-sliced eliminations); \texttt{0} retired\\
\env{L23_RING2K}, \env{L23_MOD3K} & the Smith readings of an $\ell=2$ (resp.\ $3$) cell over $\Z/2^{31}$ (resp.\ $\Z/3^{19}$) on the native ring engines\\
\env{L23_F3L} & $\rho_n$ by the $O(d^2)$ column recurrence and the $\F_3$ eliminations of an $\ell=3$ cell\\
\env{L23_BIN3O1} & the amalgam reader's $O(d^2)$ relation matrix, parity split and vectorised Taylor shifts; \texttt{0} retired\\
\bottomrule
\end{longtable}
}

\section{Diagnostics and test-only variables}\label{sec:env-diag}

{\small
\begin{longtable}{@{}>{\raggedright\arraybackslash}p{0.33\linewidth}>{\raggedright\arraybackslash}p{0.61\linewidth}@{}}
\toprule
\textbf{Variable} & \textbf{Effect}\\\midrule\endhead
\env{ENGINE_TRACE_NODES}, \env{ENGINE_TRACE_RSS}, \env{ENGINE_TRACE_ELIM}, \env{ENGINE_TRACE_LEVELS}, \env{ENGINE_TRACE_PASSES}, \env{ENGINE_TRACE_DEVICE}, \env{ENGINE_TRACE_DEVICE_RELEASE}, \env{ENGINE_TAPE_TRACE} & diagnostic traces on standard output\\
\env{ENGINE_TRACE_VRAM}\texttt{=<ms>} & samples the device memory every \texttt{<ms>} and prints a summary at exit\\
\env{ENGINE_TRACE_IMMA_ELIM}, \env{ENGINE_KERNEL_TIMING} & per-call phase times of the device elimination; the timing table of \vb{KernelBitIdentityValidation}\\
\env{ENGINE_GPU_ELIM_SIMULATE_CUDA_ERROR} & test-only error injection for the counted host fallback\\
\env{ENGINE_SCALING_REPS}, \env{ENGINE_TRAP_REPS} & repetitions of two engine validations\\
\bottomrule
\end{longtable}
}

\section{Owner-gated emitters, interlocks and strict mode}\label{sec:env-gated}

\begin{danger}
The variables in this box make an otherwise read-only command \textbf{write files}. Two of them
rewrite tracked reference files of the repository (a \emph{re-baseline}, which only the project owner
performs, explaining every changed line). Never set one of them to ``make a validation pass''.
\begin{itemize}[nosep,leftmargin=1.2em]
\item \env{GOLDEN_EMIT}\texttt{=1 ./hacs GoldenEmit} (optionally \env{GOLDEN_OUT}\texttt{=<path>})
  regenerates the golden reference \file{Golden/GoldenReference.inc}.
\item \env{CONSOLE_EMIT_EXPECTED}\texttt{=1 ./hacs ConsoleScriptEmit} regenerates the frozen
  \file{.expected} transcripts of \file{Console/scripts/}.
\item \env{GOLDEN_DUMP}\texttt{=<dir> ./hacs GoldenDump} writes the before/after object dump of
  the re-baseline protocol into \texttt{<dir>}.
\item \env{HECKE_DUMP_TAU}\texttt{=<dir>} (\vb{HeckeTauDumpValidation}) and \env{HECKE_DUMP_SAT_CERT}\texttt{=<dir>}
  (\vb{HeckeSatCertDumpValidation}) write data files for the Lean checkers into a scratch directory, which
  must exist and lie outside the source tree. Their sub-knobs select what is written:
  \env{HECKE_DUMP_TAU_RHO_CELLS}\texttt{="n:ell,\dots"} (the $\rho_n$ cells, default $(2,5),(10,691),(16,13)$),
  \env{HECKE_DUMP_TAU_PAIR} (the $\tau$-table pairs) and \env{HECKE_DUMP_SAT_CERT_CELLS}\texttt{="n:ell:kmax,\dots"}
  (default \texttt{168:13:2,220:17:2,258:37:2}); a set-but-empty or malformed value is refused (exit~1).
\item \env{CLI_MANUAL_EMIT}\texttt{=<dir>} (\vb{CliManualCoverageValidation}) writes this manual's generated
  tables into a scratch directory outside the tree (Appendix~\ref{app:regen}).
\item \env{ORBIFOLD_MP_DUMP}\texttt{=<dir>} (\vb{OrbifoldMPRulesValidation}) writes six
  \texttt{mp\_*.tsv} data files of the Moerdijk--Pronk model (orbits, points, top simplices, arrows,
  predicates, $|S_0|,|S_1|,|S_2|$) for the Lean cross-checks, and prints
  \texttt{EFFECTIVE OrbifoldMPRulesValidation arm=dump dir=<dir> files=6 fnv=0x\dots}. The directory must
  exist and lie outside the source tree, otherwise the validation prints \texttt{REFUSED} and fails;
  unset, nothing is written. It changes no value.
\item \env{HECKE_ATLAS_OUT}\texttt{=<dir>} lets \cli{--atlas} write its store;
  \env{HECKE_EXT_JSON_OUT}\texttt{=<file>} makes \vb{HeckeExtRamanujanValidation} write the $(10,691)$
  extension record to a file.
\end{itemize}
\end{danger}

\begin{itemize}
\item \env{GOLDEN_STRICT}\texttt{=1} is not an emitter: it makes \vb{GoldenRegressionValidation} gate its
  tier B (bit-exact objects) as well as tier A.
\item \env{HECKE_LOCALITY_PROBE}\texttt{=1} and \env{HECKE_HIDA_HEAVY}\texttt{=1} are \emph{interlocks on a
  run}: without them \vb{HeckeLocalityProbe} and \vb{HeckeHidaHeavyCellProbe} print \texttt{[SKIP]} even
  under \cli{--all-full}. The second can need tens of GiB of memory -- predict the peak RSS before
  enabling it.
\end{itemize}

\section{Read-only validation hooks}\label{sec:env-hooks}

Many validations accept a variable that widens or replaces the list of cells an arm runs over; none of
them writes or changes a value. They share a naming scheme:
{\small
\begin{longtable}{@{}>{\raggedright\arraybackslash}p{0.33\linewidth}>{\raggedright\arraybackslash}p{0.61\linewidth}@{}}
\toprule
\textbf{Pattern} & \textbf{Meaning}\\\midrule\endhead
\texttt{HECKE\_<ARM>\_ALL=1} & run an arm over its whole published grid (e.g.\ \env{HECKE_CENSUS_SPLIT_ALL}, \env{HECKE_THETA_IMAGE_ALL})\\
\texttt{HECKE\_<ARM>\_CELLS="n:ell,\dots"} & run an arm over the listed cells instead of its default list (e.g.\ \env{HECKE_RING_LOEWY_CELLS}, \env{HECKE_SAT_DEFECT_CELLS}, \env{HECKE_NEBENTYPUS_CELLS}, \env{HECKE_LEVEL_N_SPARSE_CELLS}, \env{HECKE_HIDA_COMPONENTS_CELLS}); a malformed cell is refused by its position\\
\texttt{\dots\_U\_ELL=<mb>}, \texttt{\dots\_RECORDS=<dirs>}, \texttt{\dots\_TSV=<file>} & the $U_\ell$ budget, the record roots or the census table an arm reads\\
\texttt{\dots\_LARGE=1} & add a few expensive cells to a coverage validation\\
\env{ORBIFOLD_MP_CELLS}\texttt{=4,6} (or \texttt{4}, \texttt{6}, \texttt{all}) & adds the nerve-$H^2$ cells $n\in\{4,6\}$ to \vb{OrbifoldMPNerveValidation} (12 cells, about 5\,s); prints \texttt{EFFECTIVE OrbifoldMPNerveValidation arm=h2 cells=<k> n=<list>}; any other token is refused (exit 1)\\
\env{HECKE_TAMEFAMILY_HEAVY}, \env{HECKE_TAME_FAMILY_CELLS}\texttt{="p:l0:l1,\dots"} & the heavy cell, respectively a survey list, of \vb{HeckeTameFamilyValidation} (the $\Gamma_H(N)$ tame-family check)\\
\env{HECKE_LAMBDA_TOWER_CELLS}\texttt{="l:a:r[:K],\dots"}, \env{HECKE_LAMBDA_TOWER_R3}\texttt{=1} & the cell list of \vb{HeckeLambdaTowerValidation}; the $r=3$ cells, refused by name above the caps \env{HECKE_LAMBDA_TOWER_RSS_CAP_GIB} (default 4) and \env{HECKE_LAMBDA_TOWER_OPS_CAP}\\
\env{HECKE_NEB_WILD_ALL}\texttt{=1} & adds the three-class cell $(7,r=3)$ to \vb{HeckeNebentypusWildValidation} (about $2.6\cdot10^{11}$ operations)\\
\env{HECKE_SHAPIRO_CELLS}\texttt{="N:n:ell:m,\dots"} & extra cells of \vb{HeckeLValueShapiroValidation} (ranks on $\Gamma_0(N)$ by Shapiro against Eichler--Shimura)\\
\env{HECKE_LN23_STURM}\texttt{=1} & \vb{HeckeLevelNCurveValidation} runs only the Sturm separating-family arms\\
\env{NF_PRANK_CE}\texttt{=CE1,CE2,CE1C,CE2C|all}, \env{NF_PRANK_CELLS}\texttt{="N:p,\dots"}, \env{NF_PRANK_DENSE} & cells and dense bound of \vb{NumberFieldPRankValidation}\\
\env{CUDA_IMMA_ELIM_BENCH}\texttt{=m,n[,p]} & times one random and one sparse $m\times n$ device elimination in \vb{CudaImmaElimValidation} (arm \texttt{bench})\\
\env{HECKE_NEBENTYPUS_NB3_FULL}\texttt{=1} & the full 132-cell nebentypus grid (run at every slot close)\\
\env{HECKE_LEVEL_N_HOOK_RSS_CAP_GIB} & the pre-run RSS cap of the level-$N$ hooks (default 8)\\
\env{CONSOLE_DETERMINISM_ARMS}\texttt{=threads} & restrict \vb{ConsoleDeterminismValidation} to its thread arms\\
\bottomrule
\end{longtable}
}
{\raggedright The descent validations take the same kind of cell lists: \env{HECKE_DESCENT_CELLS}, \env{HECKE_DESCENT_REPEATS},
\env{HECKE_DESCENT_LEVELN_CELLS}, \env{HECKE_DESCENT_N_CELLS}, \env{HECKE_DESCENT_PAIR_CELLS}, \env{HECKE_DESCENT_BDRY_CELLS},
\env{HECKE_DESCENT_TREE_CELLS}, \env{HECKE_DESCENT_TREEPAIR_CELLS}, \env{HECKE_DESCENT_HECKE_CELLS},
\env{HECKE_DESCENT_PARLN_CELLS} replace the cost cells of the descent's slow validations, and
\env{HECKE_DESCENT_SMITH_MT_BENCH}, \env{HECKE_DESCENT_SMITH_MT_CELLS}, \env{HECKE_DESCENT_SMITH_MT_THREADS},
\env{HECKE_DESCENT_SMITH_MT_REPEATS}, \env{HECKE_DESCENT_SMITH_MT_PANEL}, \env{HECKE_DESCENT_SMITH_GPU_CELLS} and
\env{HECKE_DESCENT_SMITH_GPU_ARMS} add timing arms to \vb{DescentSmithMtValidation} and \vb{DescentSmithGpuValidation}.\par}
Every such arm prints its \texttt{EFFECTIVE} line. The complete table, with the cost and memory of each
hook, is in \file{docs/validation-catalogue.md}.

\section{Hooks of the level-\texorpdfstring{$N$}{N} curve and the \texorpdfstring{$\Lambda$}{Lambda}-adic tower}\label{sec:env-ln23}

Two validations take a family of read-only variables. As for every hook of this section they only choose which
cells run (or how far a certificate climbs), never write, and never change a value; each arm prints an
\texttt{EFFECTIVE} line, which is the proof that the variable was read.

\paragraph{\vb{HeckeLevelNCurveValidation} (LN23).} The validation computes, for the weight-2 Eisenstein
block at prime level $N$ and $p=2,3$, the Hecke module $W$ of the \emph{curve} $X_0(N)$ (its dimension, Loewy length and
layers) and compares it with the published results of Calegari--Emerton (CE) on the Eisenstein ideal. It is a
\texttt{standard} validation (about 19\,s, 1177 checks). A reading is taken with the Hecke primes $q\le Q$ and is
\emph{certified} when the family contains a Mazur good prime, the reading stays the same under the next two widenings, and
the layers have the curvilinear shape $d,d-2,\dots,3,1$ (or $\dim W=1$); the validation then
compares the certified reading, and also prints the one at $q\le13$. Its lines are \texttt{LN23\_CELL},
\texttt{LN23\_THM11\_SUMMARY}, \texttt{LN23\_THM12I\_SUMMARY}, \texttt{LN23\_STURM\_SUMMARY} and so on; the last line is
\texttt{HeckeLevelNCurveValidation: <c> checks, <f> failing, <t> s}.
\begin{center}\small
\begin{tabular}{@{}>{\raggedright\arraybackslash}p{0.36\linewidth}>{\raggedright\arraybackslash}p{0.58\linewidth}@{}}
\toprule
\textbf{Variable} & \textbf{Effect}\\\midrule
\env{HECKE_LN23_QMAX}\texttt{=<q>} & the first rung of the Hecke primes $q\le Q$ every reading uses (default 13; an integer in $[2,997]$); the certificate climbs along $13,19,31,47,71,101,\dots$ until it holds\\
\env{HECKE_LN23_Z1_NMAX}\texttt{=<n>} & the torsion gate over all primes $N\le n$, plus every composite $N\le100$ (default 600; at most 20000; \texttt{2000}: 377 levels, about 1\,s)\\
\env{HECKE_LN23_THM11_NMAX}\texttt{=<n>} & the grid of CE's Theorem~1.1 (primes $N\equiv9\bmod16$; default 2000; at most 100000; \texttt{30000}: 397 cells, 72\,min, 0.77\,GiB per cell at \texttt{ENGINE\_THREADS=4})\\
\env{HECKE_LN23_THM12I_NMAX}\texttt{=<n>} & the grid of CE's Theorem~1.2(i) (primes $N\equiv10,19\bmod27$; default 2500; at most 100000; \texttt{10000}: 133 cells, about 3.5\,min)\\
\env{HECKE_LEVEL_N_R3_SMALLP_CELLS}\texttt{="N:p,\dots"} & replaces all arms by readings at the listed cells: $N$ prime, $p\in\{2,3\}$, $p\mid\mathrm{num}((N-1)/12)$; any other cell is refused by its position\\
\env{HECKE_LN23_STURM}\texttt{=1} & only the Sturm separating-family arm (Hecke primes $q\le\lfloor(N+1)/6\rfloor$); 76 cells, 228 checks, about 27\,min, 0.45\,GiB. Combined with the previous variable it runs on those cells. Any value other than \texttt{1} is an error\\
\env{HECKE_LN23_HEAVY}\texttt{=1} & adds the cell $(21929,2)$ ($\dim W=255$, Loewy length 128; about 34\,s, 0.34\,GiB at \texttt{ENGINE\_THREADS=4})\\
\bottomrule
\end{tabular}
\end{center}
A value that is not a positive integer, or lies outside its range, is never ignored: the validation stops with an \texttt{[EXCEPTION]} line naming the variable and counts as \texttt{[FAIL]}. The same holds for every token ``refused by its position'' in this section.
At the default $q\le13$ the Theorem~1.1 grid to $3\cdot10^4$ reads 369 of 397 cells without the certificate; with it,
397 of 397 agree --- which is why the certified reading is the one compared. Example:
\begin{shell}
HECKE_LN23_STURM=1 HECKE_LEVEL_N_R3_SMALLP_CELLS="41:2,73:2" ENGINE_THREADS=4 ./hacs HeckeLevelNCurveValidation
\end{shell}
prints \texttt{EFFECTIVE HeckeLevelNCurveValidation arm=STURM cells=2 (HECKE\_LN23\_STURM=1; every other arm skipped)} and, in a
fraction of a second, \texttt{HeckeLevelNCurveValidation: 6 checks, 0 failing}. Read the \texttt{cells=} count: it shows
which cells were selected.

\paragraph{\vb{HeckeLambdaTowerValidation} ($\Lambda$-adic tower).} The validation reads the Hecke modules of the
$\Lambda$-adic tower for $\ell=5,7,11,13$ ([U], \texttt{sec:u-tower}) on the cells $(\ell,a,r)$: $a$ even with $a<\ell-1$ is the
weight class and $r$ the tower level ($\Gamma_1(\ell^r)$). There are 23 default cells:  $\ell\in\{5,7\}$ with all even $a$ and $r=1,2,3$, and
$(11,2),(11,4),(13,0),(13,4)$ for $r=1,2$ (\texttt{EFFECTIVE HeckeLambdaTowerValidation arm=DEFAULT cells=23}).
\begin{center}\small
\begin{tabular}{@{}>{\raggedright\arraybackslash}p{0.36\linewidth}>{\raggedright\arraybackslash}p{0.58\linewidth}@{}}
\toprule
\textbf{Variable} & \textbf{Effect}\\\midrule
\env{HECKE_LAMBDA_TOWER_CELLS}\texttt{="l:a:r[:K],\dots"} & replaces the cell list ($K$, default 3, is the precision exponent, $1\le K\le3$; $5\le\ell$, $a$ even, $a<\ell-1$, $1\le r\le3$). A malformed or out-of-range token is refused by its position, an empty list fails, and $\ell>13$ is refused with \texttt{tower\_edge\_not\_built}. Prints \texttt{EFFECTIVE HeckeLambdaTowerValidation arm=CELLS cells=<k>}\\
\env{HECKE_LAMBDA_TOWER_R3}\texttt{=1} & adds the two $r=3$ cells $(11,2,3)$ and $(13,0,3)$, subject to the two caps below (any other value is an error)\\
\env{HECKE_LAMBDA_TOWER_RSS_CAP_GIB}\texttt{=<g>} & the memory cap of an $r=3$ cell in GiB (default 4)\\
\env{HECKE_LAMBDA_TOWER_OPS_CAP}\texttt{=<x>} & the cap on the predicted number of Krylov operations (default $10^{12}$)\\
\bottomrule
\end{tabular}
\end{center}
\textbf{How the caps work.} Before an $r=3$ cell allocates anything, the validation predicts its memory and its Krylov
work from the $r=2$ cell of the same $(\ell,a)$ measured in the same run (with \env{HECKE_LAMBDA_TOWER_CELLS} set, list
$(11,2,2)$ and $(13,0,2)$ too, or both $r=3$ cells are refused as \texttt{tower\_rss\_cap}) and prints
\texttt{EFFECTIVE HeckeLambdaTowerValidation arm=R3 cell=<l:a:r> predicted\_bytes=\dots{} cap\_bytes=\dots{} predicted\_krylov\_ops=\dots{} ops\_cap=\dots}.
If the prediction exceeds a cap, the cell is \emph{refused}, not run: \texttt{REFUSED tower\_rss\_cap <cell>} or
\texttt{REFUSED tower\_cost\_cap <cell>}, each marked ``counted, not a failure''. At the default caps both $r=3$ cells are
refused (predicted 33.5\,GB / $1.5\cdot10^{14}$ operations and 86\,GB / $6.4\cdot10^{14}$). A cell runs only when it is
under \emph{both} caps; raising them is the user's decision about the machine: read the \texttt{predicted\_bytes} line
first, and do not raise the memory cap above the free memory. The second example raises only the memory cap: it prints
both predictions and still refuses both cells ($(11,2,3)$ as \texttt{tower\_cost\_cap}, $(13,0,3)$ as \texttt{tower\_rss\_cap}),
which is the safe way to read them. Example:
\begin{shell}
HECKE_LAMBDA_TOWER_CELLS="5:0:2,7:2:3:3" ./hacs HeckeLambdaTowerValidation
HECKE_LAMBDA_TOWER_R3=1 HECKE_LAMBDA_TOWER_RSS_CAP_GIB=40 ./hacs HeckeLambdaTowerValidation   # predictions only
\end{shell}
The first line reports \texttt{arm=CELLS cells=2}; the record per cell is \texttt{HECKE\_LAMBDA\_TOWER ell= a= r= K= \dots{}
digest=0x\dots} (schema \texttt{hecke\_lambda\_tower/v1}, not pinned) and the last line is
\texttt{HECKE\_LAMBDA\_TOWER checks=<c> failing=<f> wall\_s=<t> peak\_rss\_mib=<m> verdict=OK} (50 checks, about 10\,s and
0.6\,GiB for this example).

\paragraph{Saturation defect.} \env{HECKE_SAT_DEFECT_CELLS}\texttt{="n:ell:kmax,\dots"} selects the cells of
\vb{HeckeSaturationDefectValidation}; \env{HECKE_SAT_DEFECT_STRICT}\texttt{=1} additionally fails every
\texttt{not\_stable\_at\_KMAX} block, so that a too-small \texttt{kmax} is loud. A census driver sets both for each case
(Section~\ref{sec:localcensus}).

%======================================================================
\chapter{Glossary and troubleshooting}\label{ch:reference}
%======================================================================

\section{Glossary}\label{sec:glossary}

\begin{description}[style=nextline,font=\normalfont\bfseries,nosep]
\item[cell $(n,\ell)$] a weight $n$ (the symmetric power $\Sym^n$) together with a coefficient prime $\ell$; the unit of every census and campaign. The \emph{flagship cell} is $(10,691)$.
\item[golden reference / point / object] the committed file \file{Golden/GoldenReference.inc}: for each \emph{point} (a cell and operator) the hashed \emph{objects} (matrices, polynomials, decompositions) that every build must reproduce bit for bit; \vb{golden(\dots)} compares against it. A change to it is a \emph{re-baseline} and is \emph{owner-gated}: only the project owner performs one, explaining each changed line.
\item[fingerprint] a hash of the mathematical settings in force; a non-default setting moves it, so results of different meaning never compare equal.
\item[handle] a lazy pipeline object, \texttt{declared}, \texttt{computed} or \texttt{failed} (Section~\ref{sec:types}).
\item[arm, \texttt{EFFECTIVE} line] an \emph{arm} is one environment-selected variant of a validation; the \texttt{EFFECTIVE} line says which arm actually ran and over how many cells.
\item[tier] \texttt{fast}, \texttt{standard} or \texttt{slow}: which of \cli{--all} / \cli{--all-full} runs a validation.
\item[lean] the memory-lean mode of a computation (smaller peak RSS, same values).
\item[report levels (a)/(b)/(c)] the three layers of a campaign job's report: (a) the flat cochain matrix $\Theta^{(q)}$, (b) the matrix $A^{(q)}(T_p)$ of the Hecke operator, (c) the certified decomposition (\texttt{level\_a = auto|always|never} controls (a)).
\item[shape group] the campaign jobs with the same matrix shapes (same $n$, $\ell$, degrees), batched together.
\item[\texttt{BENCH} line] the one machine-readable timing line a \cli{--bench} run prints.
\item[transcript, annotation] a frozen \file{.expected} output of a script; an annotation is a \texttt{\#\textasciitilde{}} line that comparisons ignore (Section~\ref{sec:console}).
\item[$S_0$, $S_{\mathrm{sep}}(\ell)$, ``anemic''] atlas words: $S_{\mathrm{sep}}(\ell)$ is the finite set of primes $q\ne\ell$ whose $T_q$ separate the mod-$\ell$ eigensystems of the atlas ([E]); $S_0$ the small base prime set, and ``anemic'' the eigenvalue vector $(a_q)_q$ taken over primes outside the level only (Section~\ref{sec:atlas}).
\end{description}

\section{Troubleshooting}\label{sec:trouble}

\begin{description}[style=nextline,font=\normalfont\bfseries]
\item[Everything is zero.] An odd $n$: $H^\bullet(\Gamma;\Sym^n)=0$ for odd $n$ on this model
  ({\ttfamily[C,~rmk:negI]}). Use an even $n$.
\item[\texttt{[FAIL]} for a name I did not mean to run.] The first argument was not a recognised token,
  so it and everything after it were taken as validation names -- typically a console modifier given
  before the console form (\cli{--set k=v --script f.oc}). Put the form first.
\item[\texttt{golden(\dots)}: the mathematical settings are not at their defaults.] The frozen reference
  pins the defaults only. \vb{reset\_settings()}, or start a fresh session; \vb{fingerprint()} must be
  \texttt{0x0000000000000000}.
\item[Exit code 5.] A legal setting this build cannot execute: \vb{field\_polynomial = "conway"},
  \vb{coefficient\_ring = "Z/<m>"}, or a verb that does not honour the coefficient ring in force
  (the atlas and the extension record are $\Fl$ objects; \vb{coefficient\_ring = "Z"} is legal and is
  honoured by \vb{descent} only, every other verb refuses it by name). Not a bad value -- the value is recorded, never
  silently replaced by the default.
\item[\texttt{CliManualCoverageValidation} fails with ``STALE''.] A verb, type, setting, validation or
  command-line token changed after this manual was generated. Regenerate (Appendix~\ref{app:regen});
  a ``no row'' failure means a parser gained a flag: add a row to the token table in
  \file{Console/CliManualCoverageValidation.cpp} and regenerate.
\item[\texttt{--bench micro \dots{} --backend gpu}: not available, exit 2.] No usable CUDA device, or an op
  (\texttt{rref}, \texttt{kernel}) that has no device implementation.
\item[The process was killed or the machine became unresponsive.] Memory. The large cells of the
  extension layer, the census hooks and the level-$N$ hooks can need many GiB; several hooks refuse
  cells whose predicted peak exceeds a cap. Estimate the peak RSS first (the \texttt{--bench} line reports
  \texttt{rss\_mb}; \vb{rss()} in the console), use \cli{--lean} or \texttt{lean = true} for sweeps, and
  run one heavy job at a time.
\item[An env hook seems to have no effect.] Look for its \texttt{EFFECTIVE} line. No line: the variable
  was not read (wrong name, or the validation was not run). \texttt{cells=0}: the list was empty.
\item[A run is slower than the figures in the documentation.] Timings depend on the build type (use
  \texttt{RelWithDebInfo} or \texttt{Bench}, not \texttt{Debug}), on \env{ENGINE_THREADS} and on the load
  of the machine; they are never part of a result.
\item[\texttt{cmake} fails under WSL with a missing \texttt{/run/user/\dots}.] The snap wrapper needs
  \texttt{XDG\_RUNTIME\_DIR}; call \file{/snap/cmake/current/bin/cmake} directly.
\item[A configuration change has no effect.] The build directory's \file{CMakeCache.txt} pins the first
  configuration; remove it and \file{CMakeFiles/}, or use a fresh build directory.
\item[\texttt{cmake .} fails with ``In-source builds are not supported''.] Configure into a separate
  directory (\verb|cmake -B build|) and remove the \file{CMakeCache.txt} and \file{CMakeFiles/} the attempt left
  in the repository root.
\item[AddressSanitizer reports a leak of 50\,960 bytes.] The WSL CUDA driver's one-time probe; run sanitizer
  builds with \env{ENGINE_CUDA_DISABLE}\texttt{=1}. ThreadSanitizer additionally needs
  \texttt{setarch x86\_64 -R} on this kernel.
\item[A validation reads the wrong files, or fails after a checkout was removed.] The binary reads tracked
  files through the source path baked in at configure time (Section~\ref{sec:build}).
\end{description}

%======================================================================
\appendix
%======================================================================

\chapter{Verb reference}\label{app:verbs}

All \CMNumOverloads{} overloads of the \CMNumVerbs{} verbs in \CMNumVerbGroups{} groups, in registration
order, which is also overload-resolution order. \texttt{x: Kind = v} is an optional parameter with its
default; \texttt{...} a variadic tail.

% GENERATED by docs/latex/tools/gen_cli_manual_tables.sh, i.e.
% `CLI_MANUAL_EMIT=<dir> hacs CliManualCoverageValidation`, from the
% registries of the build it runs in.  Do not edit by hand: regenerate when a
% verb, type, setting, validation or command-line token changes.
\begin{verbgroup}{core}
\ocverb{print}{\texttt{print(value: any,\allowbreak{} ...)}}{prints its arguments on one line, separated by spaces (a String prints bare)}{\index{print@\texttt{print} (verb)}}{ro}
\ocverb{show}{\texttt{show(value: any)}}{prints the full canonical form, ignoring print\_max\_entries}{\index{show@\texttt{show} (verb)}}{ro}
\ocverb{describe}{\texttt{describe(value: any)}}{prints type, state, settings snapshot, creation site and members (same as ?x)}{\index{describe@\texttt{describe} (verb)}}{ro}
\ocverb{str}{\texttt{str(value: any)}}{the display form as a String}{\index{str@\texttt{str} (verb)}}{ro}
\ocverb{repr}{\texttt{repr(value: any)}}{the canonical form as a String (strings quoted)}{\index{repr@\texttt{repr} (verb)}}{ro}
\ocverb{json}{\texttt{json(value: any)}}{JSON for Int / Bool / String / List / Record, with a proper escaper}{\index{json@\texttt{json} (verb)}}{ro}
\ocverb{annotate}{\texttt{annotate(text: String)}}{prints a '\#\textasciitilde{} ' line, which every transcript comparison ignores}{\index{annotate@\texttt{annotate} (verb)}}{ro}
\ocverb{len}{\texttt{len(value: any)}}{length of a List, String or Record}{\index{len@\texttt{len} (verb)}}{ro}
\ocverb{keys}{\texttt{keys(record: Record)}}{the keys of a Record, in order}{\index{keys@\texttt{keys} (verb)}}{ro}
\ocverb{range}{\texttt{range(from: Int,\allowbreak{} to: Int)}}{the half-open range [from, to) as a List of Int}{\index{range@\texttt{range} (verb)}}{ro}
\ocverb{range}{\texttt{range(from: Int,\allowbreak{} to: Int,\allowbreak{} step: Int)}}{the half-open range [from, to) with the given step}{}{ro}
\ocverb{abs}{\texttt{abs(value: Int)}}{absolute value}{\index{abs@\texttt{abs} (verb)}}{ro}
\ocverb{min}{\texttt{min(a: Int,\allowbreak{} b: Int)}}{the smaller of two Ints}{\index{min@\texttt{min} (verb)}}{ro}
\ocverb{max}{\texttt{max(a: Int,\allowbreak{} b: Int)}}{the larger of two Ints}{\index{max@\texttt{max} (verb)}}{ro}
\ocverb{type}{\texttt{type(value: any)}}{the type name as a String}{\index{type@\texttt{type} (verb)}}{ro}
\ocverb{assert}{\texttt{assert(condition: Bool)}}{fails the script (exit code 4) when the condition is false}{\index{assert@\texttt{assert} (verb)}}{ro}
\ocverb{assert}{\texttt{assert(condition: Bool,\allowbreak{} message: String)}}{fails the script (exit code 4) with the given message}{}{ro}
\ocverb{exit}{\texttt{exit()}}{ends the script with exit code 0}{\index{exit@\texttt{exit} (verb)}}{rw}
\ocverb{exit}{\texttt{exit(code: Int)}}{ends the script with the given exit code}{}{rw}
\ocverb{vars}{\texttt{vars()}}{every bound variable as a Record}{\index{vars@\texttt{vars} (verb)}}{ro}
\ocverb{verbs}{\texttt{verbs()}}{every registered verb name}{\index{verbs@\texttt{verbs} (verb)}}{ro}
\ocverb{types}{\texttt{types()}}{every registered type name}{\index{types@\texttt{types} (verb)}}{ro}
\ocverb{settings}{\texttt{settings()}}{every setting with its value, source (session / env) and family (engine / mathematical)}{\index{settings@\texttt{settings} (verb)}}{ro}
\ocverb{reset\_settings}{\texttt{reset\_settings()}}{re-applies the settings snapshot taken when the session started}{\index{reset_settings@\texttt{reset\_settings} (verb)}}{ro}
\ocverb{script\_args}{\texttt{script\_args()}}{the script parameters, same as \texttt{args}}{\index{script_args@\texttt{script\_args} (verb)}}{ro}
\ocverb{compute}{\texttt{compute(object: Object)}}{forces a lazily declared handle and returns it}{\index{compute@\texttt{compute} (verb)}}{ro}
\ocverb{state}{\texttt{state(object: Object)}}{state, creation site, settings snapshot and resumable state of a handle}{\index{state@\texttt{state} (verb)}}{ro}
\ocverb{resume}{\texttt{resume(object: Object)}}{continues a computation that a budget cut off, with the same budget}{\index{resume@\texttt{resume} (verb)}}{ro}
\ocverb{resume}{\texttt{resume(object: Object,\allowbreak{} budget: any = nil)}}{continues a budget-cut computation with a new budget, recorded in the resumed handle's settings snapshot}{}{ro}
\ocverb{display}{\texttt{display(value: any,\allowbreak{} ...)}}{forces each declared handle, then shows each argument: in a served notebook as a rendered value (one value frame each, like a bare expression), elsewhere as the full canonical form (like show)}{\index{display@\texttt{display} (verb)}}{ro}
\end{verbgroup}
\begin{verbgroup}{arithmetic}
\ocverb{Fp}{\texttt{Fp(p: Int,\allowbreak{} a: Int)}}{the element a of F\_p, as a canonical residue}{\index{Fp@\texttt{Fp} (verb)}}{ro}
\ocverb{poly}{\texttt{poly(p: Int,\allowbreak{} coeffs: List)}}{the polynomial with the given coefficients LOW -\textgreater{} HIGH over F\_p}{\index{poly@\texttt{poly} (verb)}}{ro}
\ocverb{X}{\texttt{X(p: Int)}}{the indeterminate X over F\_p}{\index{X@\texttt{X} (verb)}}{ro}
\ocverb{coeffs}{\texttt{coeffs(f: Poly)}}{the coefficients LOW -\textgreater{} HIGH (the machine order; the print form is HIGH -\textgreater{} LOW)}{\index{coeffs@\texttt{coeffs} (verb)}}{ro}
\ocverb{degree}{\texttt{degree(f: Poly)}}{the degree, -1 for the zero polynomial}{\index{degree@\texttt{degree} (verb)}}{ro}
\ocverb{factor}{\texttt{factor(f: Poly)}}{the certified factorisation of f (always passed through CertifyFactorization)}{\index{factor@\texttt{factor} (verb)}}{ro}
\ocverb{gcd}{\texttt{gcd(f: Poly,\allowbreak{} g: Poly)}}{the monic gcd}{\index{gcd@\texttt{gcd} (verb)}}{ro}
\ocverb{powmod}{\texttt{powmod(g: Poly,\allowbreak{} e: Int,\allowbreak{} f: Poly)}}{g\textasciicircum{}e mod f}{\index{powmod@\texttt{powmod} (verb)}}{ro}
\ocverb{eval}{\texttt{eval(f: Poly,\allowbreak{} x: Int)}}{f(x) in F\_ell}{\index{eval@\texttt{eval} (verb)}}{ro}
\ocverb{matrix}{\texttt{matrix(p: Int,\allowbreak{} rows: List)}}{a matrix over F\_p from a list of equal-length rows}{\index{matrix@\texttt{matrix} (verb)}}{ro}
\ocverb{identity}{\texttt{identity(p: Int,\allowbreak{} n: Int)}}{the n x n identity over F\_p}{\index{identity@\texttt{identity} (verb)}}{ro}
\ocverb{zero}{\texttt{zero(p: Int,\allowbreak{} rows: Int,\allowbreak{} cols: Int)}}{the zero matrix over F\_p}{\index{zero@\texttt{zero} (verb)}}{ro}
\ocverb{mul}{\texttt{mul(a: any,\allowbreak{} b: any)}}{the product of two matrices or two polynomials}{\index{mul@\texttt{mul} (verb)}}{ro}
\ocverb{add}{\texttt{add(a: any,\allowbreak{} b: any)}}{the sum of two matrices or two polynomials}{\index{add@\texttt{add} (verb)}}{ro}
\ocverb{sub}{\texttt{sub(a: any,\allowbreak{} b: any)}}{the difference of two matrices or two polynomials}{\index{sub@\texttt{sub} (verb)}}{ro}
\ocverb{scale}{\texttt{scale(A: Matrix,\allowbreak{} c: Int)}}{c * A}{\index{scale@\texttt{scale} (verb)}}{ro}
\ocverb{rank}{\texttt{rank(A: Matrix)}}{the rank of A}{\index{rank@\texttt{rank} (verb)}}{ro}
\ocverb{nullity}{\texttt{nullity(A: Matrix)}}{cols(A) - rank(A)}{\index{nullity@\texttt{nullity} (verb)}}{ro}
\ocverb{rref}{\texttt{rref(A: Matrix)}}{the reduced row echelon form of A}{\index{rref@\texttt{rref} (verb)}}{ro}
\ocverb{kernel}{\texttt{kernel(A: Matrix)}}{a basis of ker A, one column per basis vector}{\index{kernel@\texttt{kernel} (verb)}}{ro}
\ocverb{image}{\texttt{image(A: Matrix)}}{a basis of the column space of A, one column per basis vector}{\index{image@\texttt{image} (verb)}}{ro}
\ocverb{column\_space}{\texttt{column\_space(A: Matrix)}}{a basis of the column space of A (the name plan Sec. 3.5 uses)}{\index{column_space@\texttt{column\_space} (verb)}}{ro}
\ocverb{inverse}{\texttt{inverse(A: Matrix)}}{A\textasciicircum{}\{-1\}}{\index{inverse@\texttt{inverse} (verb)}}{ro}
\ocverb{transpose}{\texttt{transpose(A: Matrix)}}{A\textasciicircum{}T}{\index{transpose@\texttt{transpose} (verb)}}{ro}
\ocverb{pow}{\texttt{pow(A: Matrix,\allowbreak{} e: Int)}}{A\textasciicircum{}e for a square A and e \textgreater{}= 0 (A\textasciicircum{}0 = I)}{\index{pow@\texttt{pow} (verb)}}{ro}
\ocverb{solve}{\texttt{solve(A: Matrix,\allowbreak{} B: Matrix)}}{the exact X with A*X = B (A of full column rank; typed error otherwise)}{\index{solve@\texttt{solve} (verb)}}{ro}
\ocverb{hstack}{\texttt{hstack(parts: List)}}{the parts side by side}{\index{hstack@\texttt{hstack} (verb)}}{ro}
\ocverb{vstack}{\texttt{vstack(parts: List)}}{the parts stacked}{\index{vstack@\texttt{vstack} (verb)}}{ro}
\ocverb{submatrix}{\texttt{submatrix(A: Matrix,\allowbreak{} r0: Int,\allowbreak{} r1: Int,\allowbreak{} c0: Int,\allowbreak{} c1: Int)}}{the half-open block A[r0:r1, c0:c1]}{\index{submatrix@\texttt{submatrix} (verb)}}{ro}
\ocverb{columns}{\texttt{columns(A: Matrix,\allowbreak{} idx: List)}}{the listed columns of A, in the order given}{\index{columns@\texttt{columns} (verb)}}{ro}
\ocverb{vec}{\texttt{vec(A: Matrix)}}{A as one column, ROW-MAJOR (SubspaceOps' convention)}{\index{vec@\texttt{vec} (verb)}}{ro}
\ocverb{unvec}{\texttt{unvec(v: Matrix,\allowbreak{} rows: Int,\allowbreak{} cols: Int)}}{the inverse of vec}{\index{unvec@\texttt{unvec} (verb)}}{ro}
\ocverb{sum\_spaces}{\texttt{sum\_spaces(U: Matrix,\allowbreak{} V: Matrix)}}{a basis of the column span of [U \textbar{} V]}{\index{sum_spaces@\texttt{sum\_spaces} (verb)}}{ro}
\ocverb{intersect}{\texttt{intersect(U: Matrix,\allowbreak{} V: Matrix)}}{a basis of the intersection of the two column spans}{\index{intersect@\texttt{intersect} (verb)}}{ro}
\ocverb{contains}{\texttt{contains(U: Matrix,\allowbreak{} v: Matrix)}}{whether v lies in the column span of U}{\index{contains@\texttt{contains} (verb)}}{ro}
\ocverb{in\_span}{\texttt{in\_span(list: List,\allowbreak{} A: Matrix)}}{whether A lies in the F\_ell-span of the listed matrices}{\index{in_span@\texttt{in\_span} (verb)}}{ro}
\ocverb{restrict}{\texttt{restrict(A: Matrix,\allowbreak{} basis: Matrix)}}{the matrix of A on the invariant subspace spanned by the basis columns}{\index{restrict@\texttt{restrict} (verb)}}{ro}
\ocverb{algebra}{\texttt{algebra(ops: List)}}{the unital F\_ell-algebra generated by pairwise-commuting operators: \{dim, basis\}}{\index{algebra@\texttt{algebra} (verb)}}{ro}
\ocverb{commutant}{\texttt{commutant(ops: List)}}{\{X : X A = A X for every A in ops\}: \{dim, basis\}}{\index{commutant@\texttt{commutant} (verb)}}{ro}
\ocverb{is\_cyclic}{\texttt{is\_cyclic(ops: List,\allowbreak{} v: Matrix)}}{whether F\_ell[ops] . v is the whole space}{\index{is_cyclic@\texttt{is\_cyclic} (verb)}}{ro}
\ocverb{cyclic\_vector}{\texttt{cyclic\_vector(ops: List,\allowbreak{} seed: Int)}}{a cyclic vector for F\_ell[ops], from a MANDATORY seed}{\index{cyclic_vector@\texttt{cyclic\_vector} (verb)}}{ro}
\ocverb{random\_vector}{\texttt{random\_vector(p: Int,\allowbreak{} n: Int,\allowbreak{} seed: Int)}}{a SplitMix64 vector; the seed is mandatory, so it is reproducible}{\index{random_vector@\texttt{random\_vector} (verb)}}{ro}
\ocverb{random\_matrix}{\texttt{random\_matrix(p: Int,\allowbreak{} rows: Int,\allowbreak{} cols: Int,\allowbreak{} seed: Int)}}{a SplitMix64 matrix; the seed is mandatory}{\index{random_matrix@\texttt{random\_matrix} (verb)}}{ro}
\ocverb{jordan\_type}{\texttt{jordan\_type(A: Matrix)}}{the Jordan partition of a NILPOTENT A, from the nullities of its powers}{\index{jordan_type@\texttt{jordan\_type} (verb)}}{ro}
\ocverb{charpoly}{\texttt{charpoly(A: Matrix)}}{the characteristic polynomial (Berkowitz; division-free)}{\index{charpoly@\texttt{charpoly} (verb)}}{ro}
\ocverb{minpoly}{\texttt{minpoly(A: Matrix)}}{the minimal polynomial (Krylov)}{\index{minpoly@\texttt{minpoly} (verb)}}{ro}
\ocverb{trace}{\texttt{trace(A: Matrix)}}{the trace in F\_p}{\index{trace@\texttt{trace} (verb)}}{ro}
\ocverb{eval\_at}{\texttt{eval\_at(f: Poly,\allowbreak{} A: Matrix)}}{f(A)}{\index{eval_at@\texttt{eval\_at} (verb)}}{ro}
\end{verbgroup}
\begin{verbgroup}{pipeline}
\ocverb{sl2z\_model}{\texttt{sl2z\_model()}}{the 7-vertex SL(2,Z) orbifold model with its exact tessellation tables}{\index{sl2z_model@\texttt{sl2z\_model} (verb)}}{ro}
\ocverb{sl2z\_model}{\texttt{sl2z\_model(name: String)}}{the named model descriptor (this build ships 'sl2z\_7vertex')}{}{ro}
\ocverb{models}{\texttt{models()}}{the model names this build can build}{\index{models@\texttt{models} (verb)}}{ro}
\ocverb{sphere}{\texttt{sphere(n: Int)}}{the simplicial n-sphere (n = 1, 2, 3)}{\index{sphere@\texttt{sphere} (verb)}}{ro}
\ocverb{cylinder}{\texttt{cylinder()}}{the simplicial cylinder}{\index{cylinder@\texttt{cylinder} (verb)}}{ro}
\ocverb{moebius}{\texttt{moebius()}}{the simplicial Moebius band}{\index{moebius@\texttt{moebius} (verb)}}{ro}
\ocverb{projective\_plane}{\texttt{projective\_plane()}}{the simplicial projective plane}{\index{projective_plane@\texttt{projective\_plane} (verb)}}{ro}
\ocverb{complex}{\texttt{complex(simplices: List,\allowbreak{} name: String)}}{a complex from its simplex lists, grouped by dimension}{\index{complex@\texttt{complex} (verb)}}{ro}
\ocverb{symmetric\_power}{\texttt{symmetric\_power(model: Model,\allowbreak{} n: Int,\allowbreak{} ell: Int)}}{the local system Sym\textasciicircum{}n(F\_ell\textasciicircum{}2) on the model}{\index{symmetric_power@\texttt{symmetric\_power} (verb)}}{ro}
\ocverb{cohomology}{\texttt{cohomology(L: LocalSystem)}}{H\textasciicircum{}*\_eq(K, R; L) -- lazy: the work runs on the first computed access}{\index{cohomology@\texttt{cohomology} (verb)}}{ro}
\ocverb{dickson\_exceptional}{\texttt{dickson\_exceptional(n: Int,\allowbreak{} ell: Int)}}{whether (n, ell) is on the Dickson-invariant exceptional locus}{\index{dickson_exceptional@\texttt{dickson\_exceptional} (verb)}}{ro}
\ocverb{dickson\_count}{\texttt{dickson\_count(n: Int,\allowbreak{} ell: Int)}}{\#\{(a, b) \textgreater{}= 0 : a(ell+1) + b(ell\textasciicircum{}2-ell) = n\}}{\index{dickson_count@\texttt{dickson\_count} (verb)}}{ro}
\ocverb{hecke}{\texttt{hecke(H: Cohomology,\allowbreak{} p: Int)}}{T\_p on H\textasciicircum{}* (this IS U\_ell when p == ell -- the same code path)}{\index{hecke@\texttt{hecke} (verb)}}{ro}
\ocverb{hecke}{\texttt{hecke(H: Cohomology,\allowbreak{} p: Int,\allowbreak{} r: Int = 1)}}{the primitive double-coset pipeline for N = p\textasciicircum{}r}{}{ro}
\ocverb{family}{\texttt{family(H: Cohomology)}}{the Hecke operator family on H\textasciicircum{}* (32-bit only)}{\index{family@\texttt{family} (verb)}}{ro}
\ocverb{dual}{\texttt{dual(L: LocalSystem)}}{the contragredient D\_n of a model local system (E7; lazy)}{\index{dual@\texttt{dual} (verb)}}{ro}
\ocverb{pairing}{\texttt{pairing(L: LocalSystem,\allowbreak{} D: ContragredientLocalSystem)}}{the Poincare-duality pairing between H\_c(L) and H(D), D = dual(L) (E7; lazy)}{\index{pairing@\texttt{pairing} (verb)}}{ro}
\ocverb{s1\_scalar}{\texttt{s1\_scalar(scalar: Int,\allowbreak{} ell: Int)}}{the rank-1 local system on S\textasciicircum{}1 with scalar monodromy (a preset: no Hecke action)}{\index{s1_scalar@\texttt{s1\_scalar} (verb)}}{ro}
\ocverb{s1\_unipotent}{\texttt{s1\_unipotent(ell: Int)}}{the rank-2 unipotent local system on S\textasciicircum{}1 (a preset: no Hecke action)}{\index{s1_unipotent@\texttt{s1\_unipotent} (verb)}}{ro}
\ocverb{cylinder\_scalar}{\texttt{cylinder\_scalar(scalar: Int,\allowbreak{} ell: Int)}}{the scalar-monodromy local system on the cylinder}{\index{cylinder_scalar@\texttt{cylinder\_scalar} (verb)}}{ro}
\ocverb{cylinder\_unipotent}{\texttt{cylinder\_unipotent(ell: Int)}}{the rank-2 unipotent local system on the cylinder}{\index{cylinder_unipotent@\texttt{cylinder\_unipotent} (verb)}}{ro}
\ocverb{moebius\_scalar}{\texttt{moebius\_scalar(scalar: Int,\allowbreak{} ell: Int)}}{the scalar-twist local system on the Moebius band}{\index{moebius_scalar@\texttt{moebius\_scalar} (verb)}}{ro}
\ocverb{moebius\_unipotent}{\texttt{moebius\_unipotent(ell: Int)}}{the rank-2 unipotent local system on the Moebius band}{\index{moebius_unipotent@\texttt{moebius\_unipotent} (verb)}}{ro}
\ocverb{decompose}{\texttt{decompose(ops: any,\allowbreak{} q: Int)}}{the certified primary decomposition of H\textasciicircum{}q over a commuting Hecke family (one operator or a list of them); q is mandatory}{\index{decompose@\texttt{decompose} (verb)}}{ro}
\ocverb{report}{\texttt{report(D: Decomposition,\allowbreak{} level\_a: Bool = false)}}{the hecke\_report/v2 record of the decomposition's degree}{\index{report@\texttt{report} (verb)}}{ro}
\ocverb{report}{\texttt{report(D: Decomposition,\allowbreak{} level\_a: Bool = false,\allowbreak{} path: String = \textquotedbl{}\textquotedbl{})}}{the hecke\_report/v2 record of the decomposition's degree, also written to a file}{}{rw}
\ocverb{crosschecks}{\texttt{crosschecks(T: HeckeOperator)}}{the classical alpha/beta/omega factorization checks (i)/(ii)/(iii) against the fused pipeline -- expensive in p}{\index{crosschecks@\texttt{crosschecks} (verb)}}{ro}
\end{verbgroup}
\begin{verbgroup}{reproducibility}
\ocverb{hash}{\texttt{hash(x: any)}}{the per-kind FNV-1a-64 of the golden harness, as canonical 0x... text}{\index{hash@\texttt{hash} (verb)}}{ro}
\ocverb{hash\_int}{\texttt{hash\_int(x: any)}}{the same hash as an Int, when it fits one}{\index{hash_int@\texttt{hash\_int} (verb)}}{ro}
\ocverb{golden}{\texttt{golden(point: String,\allowbreak{} object: String)}}{the frozen GoldenReference.inc entry \{tier, rows, cols, value, hash\} -- READ-ONLY}{\index{golden@\texttt{golden} (verb)}}{ro}
\ocverb{golden\_has}{\texttt{golden\_has(point: String,\allowbreak{} object: String)}}{whether the reference has that (point, object)}{\index{golden_has@\texttt{golden\_has} (verb)}}{ro}
\ocverb{golden\_size}{\texttt{golden\_size()}}{the number of entries in the reference}{\index{golden_size@\texttt{golden\_size} (verb)}}{ro}
\ocverb{fingerprint}{\texttt{fingerprint()}}{the mathematical-settings fingerprint (0x0... at the defaults)}{\index{fingerprint@\texttt{fingerprint} (verb)}}{ro}
\ocverb{math\_settings}{\texttt{math\_settings()}}{the mathematical settings in force, with their sources and the fingerprint (read-only in C1a; \texttt{set} on them is phase C2)}{\index{math_settings@\texttt{math\_settings} (verb)}}{ro}
\ocverb{save\_state}{\texttt{save\_state(x: any)}}{the R3 resumable-state object of a handle, as JSON text}{\index{save_state@\texttt{save\_state} (verb)}}{ro}
\ocverb{save\_state}{\texttt{save\_state(x: any,\allowbreak{} path: String)}}{the same, written to a file (a state over 64 KiB goes to a sidecar referenced by relative path and FNV-1a-64)}{}{rw}
\end{verbgroup}
\begin{verbgroup}{harness}
\ocverb{validations}{\texttt{validations()}}{every registered validation as [name, tier], in registration order}{\index{validations@\texttt{validations} (verb)}}{ro}
\ocverb{validations}{\texttt{validations(tier: String)}}{the registered validations of one tier (\textquotedbl{}fast\textquotedbl{} / \textquotedbl{}standard\textquotedbl{} / \textquotedbl{}slow\textquotedbl{})}{}{ro}
\ocverb{run\_validation}{\texttt{run\_validation(name: String,\allowbreak{} inherit\_settings: Bool = false)}}{runs one registered validation in process and returns its code (0 = pass); the session's settings are lifted to its baseline for the run unless inherit\_settings, and the validation's own output is captured unless \texttt{verbose}}{\index{run_validation@\texttt{run\_validation} (verb)}}{rw}
\ocverb{bench}{\texttt{bench(mode: String,\allowbreak{} n: Int,\allowbreak{} ell: Int,\allowbreak{} p: Int = 0,\allowbreak{} r: Int = 1,\allowbreak{} family: List = List,\allowbreak{} repeat: Int = 1,\allowbreak{} lean: Bool = false,\allowbreak{} cross\_checks: Bool = false,\allowbreak{} count\_allocs: Bool = false,\allowbreak{} inherit\_settings: Bool = false)}}{one \texttt{-{}-{}bench} run (cohom / hecke / family / decomp / ext) through bench::RunMode: a record \{mode, params, dims, total\_s, total\_us, median\_s, repeats, rss\_mb, phases, extra, bench\_line\}; times are Strings}{\index{bench@\texttt{bench} (verb)}}{rw}
\ocverb{micro}{\texttt{micro(op: String,\allowbreak{} a: any = nil,\allowbreak{} b: any = nil,\allowbreak{} c: any = nil,\allowbreak{} d: any = nil,\allowbreak{} backend: String = \textquotedbl{}\textquotedbl{},\allowbreak{} repeat: Int = 1,\allowbreak{} inherit\_settings: Bool = false)}}{one \texttt{-{}-{}bench micro} op (gemm / gemv / rref / kernel / transpose / suite) with up to four integer arguments; backend \textquotedbl{}\textquotedbl{} = the ENGINE\_BACKEND default}{\index{micro@\texttt{micro} (verb)}}{rw}
\ocverb{campaign\_spec}{\texttt{campaign\_spec(path: String)}}{parses a campaign spec file (relative to the running script)}{\index{campaign_spec@\texttt{campaign\_spec} (verb)}}{rw}
\ocverb{campaign\_spec\_text}{\texttt{campaign\_spec\_text(text: String)}}{parses campaign spec text}{\index{campaign_spec_text@\texttt{campaign\_spec\_text} (verb)}}{ro}
\ocverb{campaign\_plan}{\texttt{campaign\_plan(spec: CampaignSpec)}}{the plan: one \{signature, jobs\} record per shape group, in execution order}{\index{campaign_plan@\texttt{campaign\_plan} (verb)}}{ro}
\ocverb{run\_campaign}{\texttt{run\_campaign(spec: CampaignSpec,\allowbreak{} output: String = \textquotedbl{}\textquotedbl{},\allowbreak{} inherit\_settings: Bool = false)}}{runs the campaign (writes one JSON per job, index.json and summary.tsv under \texttt{output}, default the spec's output\_dir): \{output\_dir, ran, skipped, failed, jobs, summary\_tsv, index\_json, wall\_s\}}{\index{run_campaign@\texttt{run\_campaign} (verb)}}{rw}
\ocverb{golden\_subset}{\texttt{golden\_subset(cells: List,\allowbreak{} inherit\_settings: Bool = false)}}{recomputes every golden object of the [n, ell] cells and compares both tiers: \{mismatches, digest, records\}}{\index{golden_subset@\texttt{golden\_subset} (verb)}}{ro}
\ocverb{timed}{\texttt{timed(x: any)}}{forces a declared handle and times it: \{value, wall\_s, wall\_us, peak\_rss\_mb, peak\_rss\_kb, allocs, alloc\_bytes\} (peak RSS is the PROCESS high-water mark)}{\index{timed@\texttt{timed} (verb)}}{ro}
\ocverb{bench\_line}{\texttt{bench\_line(record: Record)}}{prints the record's \texttt{BENCH ...} line as a \texttt{\#\textasciitilde{}} annotation}{\index{bench_line@\texttt{bench\_line} (verb)}}{ro}
\ocverb{bench\_line}{\texttt{bench\_line(record: Record,\allowbreak{} mode: String)}}{the same, naming the mode of a timed() record}{}{ro}
\ocverb{rss}{\texttt{rss()}}{the process's resident set: \{peak\_rss\_mb, rss\_mb, peak\_rss\_kb, rss\_kb, scope\} -- the high-water mark of the PROCESS, not of any one verb}{\index{rss@\texttt{rss} (verb)}}{ro}
\end{verbgroup}
\begin{verbgroup}{reference}
\ocverb{help\_markdown}{\texttt{help\_markdown()}}{the console reference manual (docs/console/console.md) as Markdown, generated from this registry -- [P]'s \texttt{help -{}-{}markdown}}{\index{help_markdown@\texttt{help\_markdown} (verb)}}{ro}
\end{verbgroup}
\begin{verbgroup}{extension}
\ocverb{sturm\_bound}{\texttt{sturm\_bound(k: Int,\allowbreak{} ell: Int)}}{the computed Sturm bound B'(k, ell) at the default cap 12*B(k)}{\index{sturm_bound@\texttt{sturm\_bound} (verb)}}{ro}
\ocverb{sturm\_bound}{\texttt{sturm\_bound(k: Int,\allowbreak{} ell: Int,\allowbreak{} cap: Int = 0)}}{the same with the cap pinned}{}{ro}
\ocverb{u\_ell\_decision}{\texttt{u\_ell\_decision(n: Int,\allowbreak{} ell: Int)}}{whether U\_ell fits the \texttt{u\_ell} budget in force, by the frozen RSS model}{\index{u_ell_decision@\texttt{u\_ell\_decision} (verb)}}{ro}
\ocverb{u\_ell\_decision}{\texttt{u\_ell\_decision(n: Int,\allowbreak{} ell: Int,\allowbreak{} budget\_mb: Int)}}{the same against an explicit budget in MB}{}{ro}
\ocverb{bernoulli\_mod}{\texttt{bernoulli\_mod(w: Int,\allowbreak{} ell: Int)}}{B\_w in F\_ell by the F\_ell recursion (von Staudt-Clausen: defined exactly for 0 \textless{}= w \textless{} ell - 1, refused outside it)}{\index{bernoulli_mod@\texttt{bernoulli\_mod} (verb)}}{ro}
\ocverb{glover\_factors}{\texttt{glover\_factors(n: Int,\allowbreak{} ell: Int)}}{the Glover / Ash-Stevens composition factors Sym\textasciicircum{}\{n\_i\} (x) det\textasciicircum{}\{j\_i\} of Sym\textasciicircum{}n(F\_ell\textasciicircum{}2) over F\_ell[GL\_2(F\_ell)], with multiplicities, in the PIPELINE's convention: Sym\textasciicircum{}n of the STANDARD module (not its dual), the one rho\_n acts on -- the dual would shift every det\_twist by -n}{\index{glover_factors@\texttt{glover\_factors} (verb)}}{ro}
\ocverb{canonical\_field\_polynomial}{\texttt{canonical\_field\_polynomial(ell: Int,\allowbreak{} r: Int)}}{C\_\{ell,r\}, the defining polynomial of the degree-r split field, as the \texttt{field\_polynomial} setting in force selects it: \texttt{least} (the default) gives the CanonicalLess-least monic irreducible of [E] def:canon; \texttt{conway} is a legal value this build cannot execute (exit code 5)}{\index{canonical_field_polynomial@\texttt{canonical\_field\_polynomial} (verb)}}{ro}
\ocverb{nilpotent\_part}{\texttt{nilpotent\_part(A: Matrix)}}{the nilpotent part n of A = s + n (Jordan-Chevalley over F\_ell: s = A\textasciicircum{}(ell\textasciicircum{}N))}{\index{nilpotent_part@\texttt{nilpotent\_part} (verb)}}{ro}
\ocverb{nilpotent\_part}{\texttt{nilpotent\_part(ops: List)}}{the nilpotent parts of a COMMUTING family (the generators of its radical ideal)}{}{ro}
\ocverb{radical}{\texttt{radical(alg: any)}}{rad A = the nilpotent elements of A = F\_ell[ops] (ker Frob\textasciicircum{}e), \{dim, basis\}; \texttt{alg} is a list of commuting operators, the Record of algebra(ops) or a HeckeAlgebra}{\index{radical@\texttt{radical} (verb)}}{ro}
\ocverb{tangent\_space}{\texttt{tangent\_space(alg: any)}}{m/m\textasciicircum{}2 = rad/rad\textasciicircum{}2: \{dim (over F\_ell), local, f, embdim, basis\}}{\index{tangent_space@\texttt{tangent\_space} (verb)}}{ro}
\ocverb{socle}{\texttt{socle(alg: any)}}{soc A = Ann\_A(rad A): \{dim, local, f, gorenstein, basis\}}{\index{socle@\texttt{socle} (verb)}}{ro}
\ocverb{nilpotency\_index}{\texttt{nilpotency\_index(alg: any)}}{the least nu with rad\textasciicircum{}nu = 0 (1 iff A is reduced)}{\index{nilpotency_index@\texttt{nilpotency\_index} (verb)}}{ro}
\ocverb{is\_local}{\texttt{is\_local(alg: any)}}{whether A has exactly one maximal ideal (Berlekamp: dim\{x : x\textasciicircum{}ell - x in rad\} - dim rad == 1) -- decided, never searched}{\index{is_local@\texttt{is\_local} (verb)}}{ro}
\ocverb{algebra\_invariants}{\texttt{algebra\_invariants(alg: any)}}{every invariant above in one Record}{\index{algebra_invariants@\texttt{algebra\_invariants} (verb)}}{ro}
\ocverb{loewy}{\texttt{loewy(alg: any,\allowbreak{} V: Matrix = nil)}}{the Loewy (radical) series rad\textasciicircum{}j V of an A-stable subspace V (columns; the whole space when omitted): \{dims, layers, f, layers\_over\_f, length\}}{\index{loewy@\texttt{loewy} (verb)}}{ro}
\ocverb{socle\_series}{\texttt{socle\_series(alg: any,\allowbreak{} V: Matrix = nil)}}{the socle series soc\textasciicircum{}j V, same shape as loewy}{\index{socle_series@\texttt{socle\_series} (verb)}}{ro}
\ocverb{check\_identities}{\texttt{check\_identities(H: Cohomology,\allowbreak{} primes: List = nil,\allowbreak{} items: List = nil)}}{every item HeckeIdentityCertificate defines, at H's cell, as \{items: [\{item, name, status, detail\}], code, k\_max, ...\}; primes = [p, p'] (default [2, 3]); items 6 and 7 are assembled only when listed in \texttt{items}, and an item whose inputs were not assembled reports \texttt{unavailable}}{\index{check_identities@\texttt{check\_identities} (verb)}}{ro}
\ocverb{L\_map}{\texttt{L\_map(H\_lower: Cohomology,\allowbreak{} H: Cohomology)}}{the Dickson map L\_* : H\textasciicircum{}*(M\_m) -\textgreater{} H\textasciicircum{}*(M\_n), m = n - ell - 1: \{L, certificate, induced, ranks\}}{\index{L_map@\texttt{L\_map} (verb)}}{ro}
\ocverb{ext\_record}{\texttt{ext\_record(n: Int,\allowbreak{} ell: Int,\allowbreak{} cone: Bool = false,\allowbreak{} sign\_split: Bool = false,\allowbreak{} presentation: Bool = false,\allowbreak{} sturm: Bool = true,\allowbreak{} commutant: Bool = true)}}{the cell's hecke\_ext/v5 record (Stages 1-5) under the mathematical settings in force (lazy); cone / sign\_split add the E4b-E5 / E3 objects, sturm / commutant select the Stage-4 routes}{\index{ext_record@\texttt{ext\_record} (verb)}}{ro}
\ocverb{local\_algebra}{\texttt{local\_algebra(m: MaximalIdeal)}}{the local Hecke algebra T\_m of a certified-local maximal ideal}{\index{local_algebra@\texttt{local\_algebra} (verb)}}{ro}
\ocverb{module\_structure}{\texttt{module\_structure(m: MaximalIdeal,\allowbreak{} group: String = nil)}}{the LoewyProfile of m in \texttt{group} (default: m's own group)}{\index{module_structure@\texttt{module\_structure} (verb)}}{ro}
\ocverb{presentation}{\texttt{presentation(m: MaximalIdeal)}}{generators and a minimal relation set of H\_m over T\_m}{\index{presentation@\texttt{presentation} (verb)}}{ro}
\ocverb{atlas}{\texttt{atlas(ell: Int)}}{Sigma(ell) and S\_sep(ell), in process at ell in \{5, 7, 11, 13\} (lazy)}{\index{atlas@\texttt{atlas} (verb)}}{ro}
\ocverb{GF}{\texttt{GF(ell: Int,\allowbreak{} r: Int)}}{the canonical field F\_ell[Y]/(C\_\{ell,r\}) of degree r}{\index{GF@\texttt{GF} (verb)}}{ro}
\ocverb{canonical\_field}{\texttt{canonical\_field(ell: Int,\allowbreak{} r: Int)}}{the same as GF(ell, r) ([X] Sec. 7's name)}{\index{canonical_field@\texttt{canonical\_field} (verb)}}{ro}
\ocverb{gf\_add}{\texttt{gf\_add(x: GFElt,\allowbreak{} y: GFElt)}}{x + y in GF}{\index{gf_add@\texttt{gf\_add} (verb)}}{ro}
\ocverb{gf\_sub}{\texttt{gf\_sub(x: GFElt,\allowbreak{} y: GFElt)}}{x - y in GF}{\index{gf_sub@\texttt{gf\_sub} (verb)}}{ro}
\ocverb{gf\_mul}{\texttt{gf\_mul(x: GFElt,\allowbreak{} y: GFElt)}}{x * y in GF}{\index{gf_mul@\texttt{gf\_mul} (verb)}}{ro}
\ocverb{gf\_inv}{\texttt{gf\_inv(x: GFElt)}}{x\textasciicircum{}-1 in GF (x != 0)}{\index{gf_inv@\texttt{gf\_inv} (verb)}}{ro}
\ocverb{gf\_pow}{\texttt{gf\_pow(x: GFElt,\allowbreak{} e: Int)}}{x\textasciicircum{}e in GF}{\index{gf_pow@\texttt{gf\_pow} (verb)}}{ro}
\ocverb{gf\_frobenius}{\texttt{gf\_frobenius(x: GFElt)}}{x\textasciicircum{}ell}{\index{gf_frobenius@\texttt{gf\_frobenius} (verb)}}{ro}
\ocverb{gf\_eval}{\texttt{gf\_eval(f: Poly,\allowbreak{} x: GFElt)}}{f(x) for f over F\_ell and x in GF}{\index{gf_eval@\texttt{gf\_eval} (verb)}}{ro}
\ocverb{settings\_snapshot}{\texttt{settings\_snapshot()}}{the mathematical settings in force, as a Settings value (math\_settings() is the C1a Record form)}{\index{settings_snapshot@\texttt{settings\_snapshot} (verb)}}{ro}
\ocverb{iota}{\texttt{iota(H: Cohomology)}}{the production involution iota\textasciicircum{}* on H and its +/- split (lazy)}{\index{iota@\texttt{iota} (verb)}}{ro}
\ocverb{boundary}{\texttt{boundary(H: Cohomology)}}{H\textasciicircum{}*\_bd on the boundary circle with r and T\_p\textasciicircum{}bd (lazy)}{\index{boundary@\texttt{boundary} (verb)}}{ro}
\ocverb{cone}{\texttt{cone(H: Cohomology)}}{H\textasciicircum{}*\_c = H\textasciicircum{}*(Cone(r)) with its LES and A\_c(T\_p) (lazy)}{\index{cone@\texttt{cone} (verb)}}{ro}
\ocverb{dickson\_les}{\texttt{dickson\_les(n: Int,\allowbreak{} ell: Int)}}{the long exact sequence of 0 -\textgreater{} Sym\textasciicircum{}n (x) det -\textgreater{} Sym\textasciicircum{}\{n+ell+1\} -\textgreater{} Q\_n -\textgreater{} 0 over F\_ell, ell \textgreater{}= 5, as a DicksonLES handle (lazy)}{\index{dickson_les@\texttt{dickson\_les} (verb)}}{ro}
\ocverb{hecke\_ring\_algebra}{\texttt{hecke\_ring\_algebra(n: Int,\allowbreak{} ell: Int,\allowbreak{} kmax: Int = 3)}}{the Z/ell\textasciicircum{}K Hecke-algebra route ([C] cor:stablecert; HeckeRingHeckeAlgebra), ell \textgreater{}= 5, n even, kmax = 3 by default, ell\textasciicircum{}max(kmax, K\_cls) \textless{} 2\textasciicircum{}64, U\_ell within the \texttt{u\_ell} budget: \{schema = hecke\_ring\_algebra/v1, n, ell, kmax, primes, e\_max, K\_cls, K\_run, Lv\_run, level\_mode, blocks, lines\} (K\_run = the levels read, Lv\_run = the levels at which the module route ran, level\_mode = reuse:\textless{}policy\textgreater{} \textbar{} per\_K:\textless{}reason\textgreater{}; [C] obl:precisionreuse); per block \{system, D, free\_rank, torsion, rank\_T, A1\_hilbert, AK\_hilbert, K\_stab, eps, verdict, stable\_cert, stable, cf9, hidden\}; a ring exit 5, ell \textless{} 5 exit 3}{\index{hecke_ring_algebra@\texttt{hecke\_ring\_algebra} (verb)}}{ro}
\ocverb{theta}{\texttt{theta(H: Cohomology,\allowbreak{} H\_up: Cohomology)}}{Dickson theta = L\_* : H\textasciicircum{}q(M\_n (x) det) -\textgreater{} H\textasciicircum{}q(M\_\{n+ell+1\}), source-indexed (the same map as L\_map(H, H\_up)): \{n, n\_up, ell, induced, ranks, ker, image, twist\_check, u\_ell\_zero\}}{\index{theta@\texttt{theta} (verb)}}{ro}
\ocverb{lfunction}{\texttt{lfunction(n: Int,\allowbreak{} ell: Int,\allowbreak{} a: Int,\allowbreak{} M: Int = 2,\allowbreak{} R: Int = 4,\allowbreak{} sign: String = \textquotedbl{}minus\textquotedbl{},\allowbreak{} sep: List = nil)}}{the record lfunction/v1 of branch a: L\_a(T) in Lambda/(ell\textasciicircum{}R, omega\_\{M-1\}(T)) of the ell-stabilised MTT measure ([E] def:lvmtt) of the ordinary eigensymbol picked by sep (List of [p, c]; nil = T\_2+3T\_3+7T\_5; the Eisenstein root if it occurs, else the first root mod ell), sign \textquotedbl{}minus\textquotedbl{} (r\textasciicircum{}-) or \textquotedbl{}plus\textquotedbl{} (r\textasciicircum{}+); ell \textgreater{}= 5 prime, n even, M \textgreater{}= 2, ell\textasciicircum{}M \textless{}= 4*10\textasciicircum{}6: \{schema, n, k, ell, R, M, G, u, sign, geometric, form, normalisation, a, vanishes, mu, lambda, readable, lambda\_mod\_ell, coefficients, precision, values\_modulus, critical\_values, hash, line\} ([E] lem:lvlfunprec: prec(0) = R, prec(t) = min(R, M-1-floor(log\_ell t)))}{\index{lfunction@\texttt{lfunction} (verb)}}{ro}
\ocverb{kummer\_unit}{\texttt{kummer\_unit(ell: Int,\allowbreak{} i: Int)}}{the cyclotomic unit eta\_i of layer 1 ([E] sec:lvnf) and its Kummer l-th-power test, ell prime, 5 \textless{}= ell \textless{}= 65521, i even, 2 \textless{}= i \textless{}= ell - 3: \{schema, ell, i, exponents, coefficients, max\_digits, v\_pi, unramified, lambda\_split, not\_lth\_power, eigen\_ok, unit\_ok, ell\_divides\_Bi, oracle\_agrees, split\_primes\}}{\index{kummer_unit@\texttt{kummer\_unit} (verb)}}{ro}
\ocverb{hida\_tower}{\texttt{hida\_tower(n: Int,\allowbreak{} ell: Int,\allowbreak{} m: Int)}}{the Hida weight-tower cell at class n ([C] def:hidatower), ell \textgreater{}= 5 prime, n even \textless{}= 200, m \textless{}= 3: both readings of the ordinary part at levels e = 1 .. m, \{schema, n, ell, m, e\_max, K, h, exps, levels: [\{e, ord\_rank, image\_free, charpoly, free\_ord\_rank, free\_charpoly, readings\_agree\}], lines\}}{\index{hida_tower@\texttt{hida\_tower} (verb)}}{ro}
\ocverb{level\_n}{\texttt{level\_n(N: Int,\allowbreak{} n: Int,\allowbreak{} ell: Int,\allowbreak{} primes: List = nil)}}{H\textasciicircum{}1(Gamma\_0(N); Sym\textasciicircum{}n F\_ell) by Shapiro's lemma ([C] sec:levelN; the sparse cokernel of the induced module) and the Heilbronn operators T\_p (p not dividing N) / U\_p (p \textbar{} N) on it, primes = List of at most 6 primes \textless{}= 97, none equal to ell (nil = 2, 3, 5 without ell); ell \textgreater{}= 5 prime, n even, N \textless{}= 100, index(N) * (n + 1) \textless{}= 1000: \{schema, N, n, ell, index, D, dim\_S, dim\_ST, h0, h1, cusps, dim\_cusp\_forms, ops: [\{label, p, bad, charpoly\}]\}}{\index{level_n@\texttt{level\_n} (verb)}}{ro}
\ocverb{nebentypus}{\texttt{nebentypus(N: Int,\allowbreak{} m: Int,\allowbreak{} n: Int,\allowbreak{} ell: Int,\allowbreak{} primes: List = nil)}}{H\textasciicircum{}1(Gamma\_0(N); Sym\textasciicircum{}n F\_ell (x) chi) for the Dirichlet character of Conrey label N.m reduced into F\_ell ([C] sec:nebentypus), and the Heilbronn operators on it; ell prime not dividing N with order(chi) \textbar{} ell - 1, n \textgreater{}= 0, N \textless{}= 100, index(N) * (n + 1) \textless{}= 1000; primes = List of at most 4 primes \textless{}= 13 (nil = 2, 3, 5 without ell): \{schema, label, N, m, order, parity\_ok, n, ell, D, h0, h1, ops: [\{label, p, bad, charpoly\}]\} (h0 = h1 = 0 and no ops when chi(-1) != (-1)\textasciicircum{}n)}{\index{nebentypus@\texttt{nebentypus} (verb)}}{ro}
\ocverb{forms\_algebra}{\texttt{forms\_algebra(k: Int,\allowbreak{} ell: Int,\allowbreak{} primes: List = nil,\allowbreak{} terms: Int = 8)}}{the forms side mod ell: S\_k over F\_ell with its Miller q-expansion basis and T\_p on it ([E] sec:forms; the characteristic polynomials), ell \textgreater{}= 5 prime, k even, 12 \textless{}= k \textless{}= 400, primes (at most 6, \textless{}= 97, none equal to ell; nil = 2, 3, 5 without ell), terms \textless{}= 64: \{schema, k, ell, d, k\_prime, precision, basis, ops: [\{p, charpoly\}]\}}{\index{forms_algebra@\texttt{forms\_algebra} (verb)}}{ro}
\ocverb{hida\_lambda}{\texttt{hida\_lambda(ell: Int,\allowbreak{} k: Int,\allowbreak{} W: Int,\allowbreak{} primes: List = nil)}}{the forms side of the Lambda-adic Hecke algebra at one weight ([U] sec:u-lambda): the ordinary parts e M\_k and e S\_k over Z/ell\textasciicircum{}W and the characteristic polynomials of T\_p and U = T\_ell on e S\_k; ell \textgreater{}= 5 prime \textless{}= 101, k even with 4 \textless{}= k \textless{}= 200 and k - 1 \textgreater{}= W, W \textless{}= 6, ell\textasciicircum{}W \textless{} 2\textasciicircum{}62; primes as for forms\_algebra (nil = 2, 3, 5 without ell): \{schema, ell, k, W, dim\_M, dim\_S, rank\_M, rank\_S, comm\_ok, stable\_ok, eisenstein, ops: [\{name, charpoly\}]\}}{\index{hida_lambda@\texttt{hida\_lambda} (verb)}}{ro}
\ocverb{descent}{\texttt{descent(N: Int,\allowbreak{} n: Int,\allowbreak{} primes: List = nil)}}{the integral descent record of H\textasciicircum{}q(Gamma\_0(N); Sym\textasciicircum{}n Z\textasciicircum{}2) over Z ([C] thm:descentcert, sec:descentpair; slot-24 D6), a lazy DescentRecord: groups to q\_max = 3, bad primes, K per prime, char\_Z(T\_p), the torsion Hecke modules, the boundary pair; only under coefficient\_ring = \textquotedbl{}Z\textquotedbl{}; primes = List of 1 .. 6 primes \textless{}= 97 not dividing N, sum \textless{}= 100 (nil = the first three of 2, 3, 5, 7, 11, 13 not dividing N); N \textless{}= 300, n \textless{}= 60, index(N) * (n + 1) \textless{}= 300 (at the bound about 40 s and 1.4 GiB)}{\index{descent@\texttt{descent} (verb)}}{ro}
\ocverb{hecke\_structure}{\texttt{hecke\_structure(N: Int,\allowbreak{} n: Int,\allowbreak{} algebra: String = \textquotedbl{}full\textquotedbl{})}}{the record integral\_hecke\_structure/v1 of the cell (N, n, algebra) ([C] sec:recordschema item 8; slot-25 HZ-REC): status, record\_fnv1a64, cert\_sha256, the status signature and the hashed body (G0 .. G8) as a Record; only under coefficient\_ring = \textquotedbl{}Z\textquotedbl{}; algebra = \textquotedbl{}full\textquotedbl{} \textbar{} \textquotedbl{}anemic\textquotedbl{} (N \textgreater{} 1); N \textless{}= 300, n even \textless{}= 60, index(N) * (n + 1) \textless{}= 300 (the descent verb's bound; beyond it a terminal REPL asks first, elsewhere it refuses unless confirm = \textquotedbl{}yes\textquotedbl{})}{\index{hecke_structure@\texttt{hecke\_structure} (verb)}}{ro}
\end{verbgroup}


\chapter{Type reference}\label{app:types}

\emph{Computed} members force the handle's computation; \emph{structural} members are available while it is
still \texttt{declared}. Members with arguments are methods: \texttt{H.H(1)}, \texttt{T.A(1)}.

% GENERATED by docs/latex/tools/gen_cli_manual_tables.sh, i.e.
% `CLI_MANUAL_EMIT=<dir> hacs CliManualCoverageValidation`, from the
% registries of the build it runs in.  Do not edit by hand: regenerate when a
% verb, type, setting, validation or command-line token changes.
\octype{Fp}{an element of F\_p (plan Sec. 3.4)}{\index{Fp@\texttt{Fp} (type)}}
\begin{membertable}
\texttt{.p} & structural & the modulus \\
\texttt{.value} & structural & the canonical residue in [0, p) \\
\end{membertable}
\octype{Poly}{a polynomial over F\_ell; prints HIGH -\textgreater{} LOW degree, while coeffs / hash use the machine order LOW -\textgreater{} HIGH (plan Sec. 2 C5)}{\index{Poly@\texttt{Poly} (type)}}
\begin{membertable}
\texttt{.ell} & structural & the modulus \\
\texttt{.deg} & structural & the degree (-1 for the zero polynomial) \\
\texttt{.monic} & structural & whether the leading coefficient is 1 \\
\end{membertable}
\octype{Matrix}{a dense matrix over F\_p (FlatMatrixT\textless{}uint32\_t\textgreater{})}{\index{Matrix@\texttt{Matrix} (type)}}
\begin{membertable}
\texttt{.rows} & structural & the row count \\
\texttt{.cols} & structural & the column count \\
\texttt{.p} & structural & the modulus \\
\texttt{.entry(row,\allowbreak{} col)} & structural & the (row, col) entry as a canonical residue \\
\end{membertable}
\octype{Model}{the SL(2,Z) orbifold model: its triangulation and the exact tessellation tables}{\index{Model@\texttt{Model} (type)}}
\begin{membertable}
\texttt{.name} & structural & the descriptor name \\
\texttt{.dims} & structural & the simplex counts by dimension \\
\texttt{.vertices} & structural & the vertex count \\
\end{membertable}
\octype{Complex}{a simplicial complex (the presets of SimplicialCohomology.h and complex(...))}{\index{Complex@\texttt{Complex} (type)}}
\begin{membertable}
\texttt{.name} & structural & the complex's name \\
\texttt{.dim} & structural & the top dimension \\
\texttt{.euler} & structural & the Euler characteristic \\
\texttt{.dims} & structural & the simplex counts by dimension \\
\texttt{.simplices(q)} & structural & the q-simplices, each as a List of vertex indices \\
\end{membertable}
\octype{LocalSystem}{Sym\textasciicircum{}n(F\_ell\textasciicircum{}2) on a model, with its rho\_n(g\_r) relations (lazy)}{\index{LocalSystem@\texttt{LocalSystem} (type)}}
\begin{membertable}
\texttt{.n} & structural & the symmetric power \\
\texttt{.p} & structural & the coefficient prime ell \\
\texttt{.d\_V} & structural & dim Sym\textasciicircum{}n = n + 1 \\
\texttt{.coefficient\_ring} & structural & the ring's canonical text: F\_ell, or Z/ell\textasciicircum{}\textless{}k\textgreater{} ([E] sec:ring (a)) \\
\texttt{.relations} & computed & the number of local-system relations \\
\end{membertable}
\octype{Cohomology}{H\textasciicircum{}q\_eq(K, R; L) with its C / Z / B / H bases and projectors (lazy); over Z/ell\textasciicircum{}k (coefficient\_ring, ell \textgreater{}= 5) its lengths, invariant factors and K-mode}{\index{Cohomology@\texttt{Cohomology} (type)}}
\begin{membertable}
\texttt{.computed} & structural & whether compute\_cohomology() has run \\
\texttt{.dims} & computed & dim H\textasciicircum{}q for every degree q \\
\texttt{.ell} & computed & the coefficient prime \\
\texttt{.d\_V} & computed & the representation dimension \\
\texttt{.q\_max} & computed & the top degree of THIS complex (derived) \\
\texttt{.truncated} & computed & false on a finite model \\
\texttt{.q\_max\_source} & computed & 'derived' unless truncated (E9) \\
\texttt{.degrees} & computed & [0 .. q\_max], derived \\
\texttt{.coefficient\_ring} & structural & the ring's canonical text: F\_ell, or Z/ell\textasciicircum{}\textless{}k\textgreater{} ([E] sec:ring (a)) \\
\texttt{.lengths} & computed & the composition length of H\textasciicircum{}q over Z/ell\textasciicircum{}k, q = 0, 1, 2 (lengths\_H\textasciicircum{}q) \\
\texttt{.ladic\_type\_certified} & computed & whether the K-mode certifies the Z\_(ell)-type at K = k ([E] sec:ring (d)) \\
\texttt{.precision} & computed & on a digit handle (ell\textasciicircum{}k \textgreater{}= 2\textasciicircum{}64): the certified precision bound -- \{precision\_requested, precision\_working, e\_max, word\_level, precision\_certified, digit\_width, digits\_L\} \\
\texttt{.bockstein\_rank} & computed & bockstein\_rank\textasciicircum{}0: the rank of beta\textasciicircum{}0 on H\textasciicircum{}0(F\_ell), from the Z/ell\textasciicircum{}2 and F\_ell objects ([E] sec:ring (d), (g)) \\
\texttt{.higher\_torsion\_count} & computed & dim H\textasciicircum{}0(F\_ell) - bockstein\_rank\textasciicircum{}0 ([E] sec:ring (d)) \\
\texttt{.C(q)} & computed & a basis of C\textasciicircum{}q\_eq, one N\_q x d\_V cochain per element \\
\texttt{.Z(q)} & computed & a basis of Z\textasciicircum{}q\_eq \\
\texttt{.B(q)} & computed & a basis of B\textasciicircum{}q\_eq \\
\texttt{.H(q)} & computed & a basis of H\textasciicircum{}q\_eq \\
\texttt{.P(q)} & computed & the projector C\textasciicircum{}q\_flat -\textgreater{} H\textasciicircum{}q coordinates (legacy shape conventions) \\
\texttt{.dim\_C(q)} & computed & dim C\textasciicircum{}q\_eq (available even in the memory-lean mode) \\
\texttt{.invariant\_factors(q)} & computed & H\textasciicircum{}q = (+)\_i Z/ell\textasciicircum{}\{e\_i\}: the exponents e\_i, descending (invariant\_factors\_H\textasciicircum{}q) \\
\texttt{.r\_K(q)} & computed & r\_K(d\textasciicircum{}q) for q in \{0, 1\}: the K-mode's rank of the coboundary \\
\texttt{.ladic\_type(q)} & computed & \{free\_rank, torsion\} of H\textasciicircum{}q(C)\_(ell); defined only at a certified level (else the object's named std::logic\_error) \\
\texttt{.red(q)} & computed & red\_H\textasciicircum{}q: H\textasciicircum{}q(Z/ell\textasciicircum{}k) -\textgreater{} H\textasciicircum{}q(F\_ell), the (c5) matrix from the cyclic summands to the F\_ell object's D5b basis, over F\_ell \\
\end{membertable}
\octype{HeckeOperator}{T\_p (or U\_ell, the same path) on H\textasciicircum{}*: the coset tables, the certified tau datum and the assembled Theta / A\textasciicircum{}\{(q)\}}{\index{HeckeOperator@\texttt{HeckeOperator} (type)}}
\begin{membertable}
\texttt{.p} & structural & the Hecke prime \\
\texttt{.r} & structural & the exponent of N = p\textasciicircum{}r \\
\texttt{.cosets} & computed & the number of cosets m \\
\texttt{.certificates} & computed & the in-process certificate bundle compute() already ran (O12(a) and the counters) \\
\texttt{.json} & computed & the operator's own report JSON (levels (a) + (b)) \\
\texttt{.dims} & computed & dim H\textasciicircum{}q of the underlying complex \\
\texttt{.coefficient\_ring} & structural & the ring's canonical text: F\_ell, or Z/ell\textasciicircum{}\textless{}k\textgreater{} ([E] sec:ring (a)) \\
\texttt{.charpoly\_free} & computed & over Z/ell\textasciicircum{}k: \{precision, coeffs\} -- the charpoly of A\textasciicircum{}\{(1)\}(T\_p) on the free quotient, over Z/ell\textasciicircum{}precision, coefficients low -\textgreater{} high; defined only at a certified level \\
\texttt{.A(q)} & computed & A\textasciicircum{}\{(q)\}(T\_p), the cohomology-level matrix in degree q (over Z/ell\textasciicircum{}k: the (c5) matrix on the cyclic summands, a List of rows, entry (i, j) in [0, ell\textasciicircum{}\{e\_i\})) \\
\texttt{.Theta(q)} & computed & Theta\textasciicircum{}\{(q)\}\_flat on the flat cochains C\textasciicircum{}q(K; V) \\
\texttt{.Theta\_eq(q)} & computed & Theta\textasciicircum{}\{(q)\} restricted to C\textasciicircum{}q\_eq in the stored C\_eq basis \\
\texttt{.apply(q,\allowbreak{} v)} & computed & Theta\textasciicircum{}\{(q)\} applied to a flat cochain (the fused route) \\
\texttt{.apply\_route2(q,\allowbreak{} v)} & computed & the INDEPENDENT second route of [C, eq:heckechain] on the same cochain (matrix-free; it does not read the fused tables) \\
\end{membertable}
\octype{Family}{the Hecke operator family on one cohomology: T\_p, T\_\{p\textasciicircum{}r\} (recursion and the independent direct route), T\_n and the diamond T(p,p)}{\index{Family@\texttt{Family} (type)}}
\begin{membertable}
\texttt{.T(p)} & computed & T\_p per degree (this IS U\_ell at p = ell) \\
\texttt{.T\_pow(p,\allowbreak{} r)} & computed & the FULL T\_\{p\textasciicircum{}r\} by the algebra recursion of [C, thm:relations] \\
\texttt{.T\_pow\_direct(p,\allowbreak{} r)} & computed & the same operator by the INDEPENDENT double-coset route -- what T\_pow is cross-checked against \\
\texttt{.T\_pow\_primitive(p,\allowbreak{} r)} & computed & the PRIMITIVE double-coset operator T'\_\{p\textasciicircum{}r\} (equal to T\_p at r = 1; it differs from the standard T\_\{p\textasciicircum{}r\} for r \textgreater{}= 2) \\
\texttt{.T\_n(n)} & computed & T\_n for a composite index n \\
\texttt{.T\_pp(p)} & computed & the diamond T(p,p) = p\textasciicircum{}n I ([C, prop:diamond]) \\
\texttt{.theta(p,\allowbreak{} q)} & computed & Theta\textasciicircum{}\{(q)\}\_flat of the family's r = 1 pipeline for p \\
\end{membertable}
\octype{ContragredientLocalSystem}{D\_n = Sym\textasciicircum{}n(F\_ell\textasciicircum{}2)\textasciicircum{}vee on the 7-vertex model (rho\_n\textasciicircum{}vee(h) = rho\_n(adj h)\textasciicircum{}T), its cohomology, Phi\_n : M\_n -\textgreater{} D\_n and the Hecke action A(D\_n)(T\_p) (E7; lazy)}{\index{ContragredientLocalSystem@\texttt{ContragredientLocalSystem} (type)}}
\begin{membertable}
\texttt{.n} & structural & the symmetric power \\
\texttt{.p} & structural & the coefficient prime ell \\
\texttt{.d\_V} & structural & dim D\_n = n + 1 \\
\texttt{.coefficient\_ring} & structural & F\_ell (E7 is a field object) \\
\texttt{.phi\_is\_isomorphism} & structural & whether Phi\_n (x) F\_ell is an isomorphism, i.e. n \textless{} ell ([C] lem:symdual) \\
\texttt{.dims} & computed & dim H\textasciicircum{}q(Gamma; D\_n) for every degree q (golden dim\_H\_Dn\textasciicircum{}q) \\
\texttt{.dims\_M} & computed & dim H\textasciicircum{}q(Gamma; M\_n), the standard module, for comparison \\
\texttt{.Phi} & computed & Phi\_n : M\_n -\textgreater{} D\_n, (n+1) x (n+1) over F\_ell \\
\texttt{.induced\_Phi(q)} & computed & P\textasciicircum{}q = H\textasciicircum{}q(Phi\_n) : H\textasciicircum{}q(M\_n) -\textgreater{} H\textasciicircum{}q(D\_n) \\
\texttt{.A(p,\allowbreak{} q)} & computed & A(D\_n)(T\_p)\textasciicircum{}q, the Hecke action on H\textasciicircum{}q(D\_n) (U\_ell at p = ell); charpoly(D.A(p, q)) is the golden charpoly\_A\_Dn(T\_p)\textasciicircum{}q \\
\texttt{.A\_M(p,\allowbreak{} q)} & computed & A(M\_n)(T\_p)\textasciicircum{}q on the standard module, built from the SAME coset / tau data \\
\texttt{.intertwines(p)} & computed & \{q0, q1, q2\}: whether A(D\_n)(T\_p) P\textasciicircum{}q == P\textasciicircum{}q A(M\_n)(T\_p) exactly (UNTWISTED, factor 1; [C] prop:pairing) \\
\end{membertable}
\octype{Pairing}{the cone-level pairing H\textasciicircum{}t\_c(M\_n) x H\textasciicircum{}\{2-t\}(D\_n) -\textgreater{} F\_ell and the circle pairing: Grams, perfectness, the PD adjointness gate of [C] cl:pd (E7; lazy)}{\index{Pairing@\texttt{Pairing} (type)}}
\begin{membertable}
\texttt{.coefficient\_ring} & structural & F\_ell (E7 is a field object) \\
\texttt{.certificate} & computed & the pairing certificate re-derived: \{failing, failing\_items\} \\
\texttt{.perfect(t)} & computed & whether the cone pairing in H\_c degree t is perfect (golden pairing\_perfect\textasciicircum{}t) \\
\texttt{.gram(t)} & computed & G\textasciicircum{}t: dim H\textasciicircum{}t\_c x dim H\textasciicircum{}\{2-t\}(D\_n) (golden pairing\_matrix\textasciicircum{}t) \\
\texttt{.boundary\_perfect(s)} & computed & whether the circle pairing in degree s is perfect \\
\texttt{.boundary\_gram(s)} & computed & the circle Gram in degree s \\
\texttt{.pd\_gate(primes)} & computed & \{t0, t1, t2\}: A\_c\textasciicircum{}\{(t)T\} G\textasciicircum{}t == G\textasciicircum{}t A(D\_n)\textasciicircum{}\{(2-t)\} for EVERY p in \texttt{primes} (factor 1; t = 0 vacuous and true); at primes [2, 3, 5] it is the golden pd\_adjoint\textasciicircum{}t. A TEST of [C] cl:pd, which is open \\
\end{membertable}
\octype{Block}{one primary block of a decomposition: its dimension, residue degree, the per-prime minimal / characteristic polynomials and generalized-Jordan partition, and the E0b locality certificate}{\index{Block@\texttt{Block} (type)}}
\begin{membertable}
\texttt{.dim} & computed & dim\_F of the block \\
\texttt{.residue\_degree} & computed & f = lcm\_j deg(pi\_j) \\
\texttt{.certified} & computed & whether f is CERTIFIED (a witness proves it), not merely a lower bound \\
\texttt{.local} & computed & the E0b locality certificate: this block is ONE maximal ideal \\
\texttt{.phase} & computed & which phase reached the verdict (linear / berlekamp / residue\_search) \\
\texttt{.w\_hat\_exponent} & computed & N: the Teichmueller exponent of the witness \\
\texttt{.witness\_poly} & computed & pi\_w, the certifying witness's minimal polynomial \\
\texttt{.basis} & computed & the d x dim matrix whose columns span the block \\
\texttt{.primes} & computed & the family's primes, ascending \\
\texttt{.witness} & computed & the linear residue witness as [[p, c], ...] (empty in strict mode, where the witness is a matrix) \\
\texttt{.min\_poly(p)} & computed & the minimal polynomial of T\_p on this block \\
\texttt{.char\_poly(p)} & computed & the characteristic polynomial of T\_p on this block \\
\texttt{.character(p)} & computed & P\_p, the character polynomial of the locality certificate \\
\texttt{.jordan(p)} & computed & the generalized-Jordan partition of T\_p on this block, parts descending \\
\end{membertable}
\octype{Decomposition}{the certified primary decomposition of H\textasciicircum{}q as a module over the commutative F\_ell-algebra generated by a commuting family of Hecke matrices}{\index{Decomposition@\texttt{Decomposition} (type)}}
\begin{membertable}
\texttt{.q} & structural & the cohomological degree \\
\texttt{.ell} & structural & the coefficient prime \\
\texttt{.n} & structural & n\_sym of the coefficient system \\
\texttt{.primes} & structural & the family's primes, in the order given \\
\texttt{.blocks} & computed & the primary blocks, in canonical order \\
\texttt{.num\_blocks} & computed & how many primary blocks there are \\
\texttt{.U} & computed & the change of basis: the block bases, in block order \\
\texttt{.witness\_hash} & computed & the golden \texttt{decomp\_*\_witness} hash (the residue-witness data), as 0x... text \\
\texttt{.locality} & computed & the E0b locality loop's state: mode, budget, splits, completeness \\
\texttt{.json} & computed & the decomposition's own JSON record \\
\texttt{.certify()} & computed & the UNCONDITIONAL certificate: 0 iff every claim re-derives (never trusting the randomized factor search) \\
\texttt{.conjugated(p)} & computed & U\textasciicircum{}\{-1\} A\textasciicircum{}\{(q)\}(T\_p) U, verified exactly block-diagonal \\
\end{membertable}
\octype{Report}{the hecke\_report/v2 record of one degree: level (b) A\textasciicircum{}\{(q)\}(T\_p), level (c) the decomposition, and optionally level (a) Theta\textasciicircum{}\{(q)\}\_flat}{\index{Report@\texttt{Report} (type)}}
\begin{membertable}
\texttt{.schema} & structural & the frozen schema name \\
\texttt{.level\_a} & structural & whether level (a) was included \\
\texttt{.path} & structural & the file it was written to, or \textquotedbl{}\textquotedbl{} \\
\texttt{.degree} & computed & the degree q this record is about \\
\texttt{.q\_max} & computed & the top degree of THIS complex \\
\texttt{.truncated} & computed & false on a finite model \\
\texttt{.q\_max\_source} & computed & \textquotedbl{}derived\textquotedbl{} unless a setting pinned q\_max \\
\texttt{.degrees} & computed & [0 .. q\_max] \\
\texttt{.json} & computed & the report text \\
\texttt{.write(path)} & computed & write the report to a file (an empty path writes to stdout) \\
\texttt{.has\_key(key)} & computed & whether the record carries a given JSON key -- the schema assertion a transcript can make without freezing the whole text \\
\end{membertable}
\octype{CampaignSpec}{a parsed campaign spec (docs/matrix-engine/campaign.md) -- campaign\_spec(path), campaign\_spec\_text(text)}{\index{CampaignSpec@\texttt{CampaignSpec} (type)}}
\begin{membertable}
\texttt{.name} & structural & the spec's \texttt{name} \\
\texttt{.output\_dir} & structural & where run\_campaign writes by default \\
\texttt{.jobs} & structural & the number of jobs the plan has \\
\end{membertable}
\octype{Sturm}{the computed Sturm bound B'(k, ell) of [E] def:bprime: the least prime bound at which the F\_ell-ALGEBRA generated by T\_q, q \textless{}= B, reaches dim M\_k}{\index{Sturm@\texttt{Sturm} (type)}}
\begin{membertable}
\texttt{.k} & structural & the weight k = n + 2 \\
\texttt{.ell} & structural & the coefficient prime \\
\texttt{.outcome} & computed & found / none\_at\_cap / in\_progress \\
\texttt{.found} & computed & whether B' exists at or below the cap \\
\texttt{.B\_prime} & computed & B', meaningful iff \texttt{found} \\
\texttt{.cap} & computed & the cap in force (12*B(k) by default) \\
\texttt{.dim\_M\_k} & computed & dim M\_k(F\_ell) \\
\texttt{.closure\_dim} & computed & the dimension the algebra closure reached \\
\texttt{.primes\_done} & computed & the primes processed, ascending \\
\texttt{.certified} & computed & SturmBoundCertificate::verify() == 0, re-derived from the reported data \\
\texttt{.verify} & computed & the certificate's code (0 = every clause re-derives) \\
\texttt{.json} & computed & the certificate's JSON, \texttt{resumable} object included \\
\end{membertable}
\octype{Atlas}{Sigma(ell): the anemic level-1 eigensystems over weights \textless{}= ell + 1, their theta-twists, and the separating set S\_sep(ell) (HeckeAtlas, in process at ell in \{5, 7, 11, 13\}); lazy}{\index{Atlas@\texttt{Atlas} (type)}}
\begin{membertable}
\texttt{.ell} & structural & the prime \\
\texttt{.n\_max} & computed & the largest even weight cell n \textless{}= ell - 1 \\
\texttt{.primes} & computed & the anemic pool T \\
\texttt{.num\_systems} & computed & \textbar{}Sigma(ell)\textbar{} \\
\texttt{.chi\_count} & computed & distinct chi\_\{a,b\} as T-systems \\
\texttt{.chi\_in\_sigma} & computed & of which in Sigma \\
\texttt{.separating\_primes} & computed & the greedy S\_sep(ell) \\
\texttt{.separates} & computed & S\_sep separates Sigma u chi \\
\texttt{.s0\_separates} & computed & \{2,3,5,7\} \textbackslash{} \{ell\} already separates \\
\texttt{.cells} & computed & the weight cells \{n, dims, ideals, local\_ideals\} \\
\texttt{.json} & computed & the hecke\_atlas/v1 store text \\
\texttt{.system(i)} & computed & the i-th system \{id, r, type, base, twist, chi, eigenvalues, collected, in\_H1\} \\
\texttt{.lookup(eigenvalues)} & computed & the id of the system with these eigenvalues on \texttt{primes} (Ints at r = 1, coefficient lists at r \textgreater{} 1; any Frobenius conjugate), or -1 \\
\end{membertable}
\octype{GF}{the project-canonical field F\_ell[Y]/(C\_\{ell,r\}) ([E] def:canon; HeckeCanonicalField), C chosen by the \texttt{field\_polynomial} setting}{\index{GF@\texttt{GF} (type)}}
\begin{membertable}
\texttt{.ell} & structural & the prime \\
\texttt{.degree} & structural & r = [F : F\_ell] \\
\texttt{.polynomial} & structural & C\_\{ell,r\} \\
\texttt{.canonical} & structural & C is the CanonicalLess-least monic irreducible \\
\texttt{.zero} & structural & 0 \\
\texttt{.one} & structural & 1 \\
\texttt{.gen} & structural & Y, the class of the variable \\
\texttt{.elt(coeffs)} & structural & the element sum c\_i Y\textasciicircum{}i (coefficients low -\textgreater{} high, reduced mod C) \\
\texttt{.element\_at(i)} & structural & the i-th element of the field's fixed enumeration \\
\end{membertable}
\octype{GFElt}{an element of a GF field (coefficient vector in the basis 1, Y, ..., Y\textasciicircum{}\{r-1\})}{\index{GFElt@\texttt{GFElt} (type)}}
\begin{membertable}
\texttt{.coeffs} & structural & the coefficients, low -\textgreater{} high, width r \\
\texttt{.is\_zero} & structural & x == 0 \\
\texttt{.field} & structural & the GF handle \\
\texttt{.poly} & structural & the representative polynomial in Y \\
\end{membertable}
\octype{Settings}{a snapshot of the MATHEMATICAL settings (hecke::MathSettings, [X] Sec. 1.5): every registered key with its value and source, and the fingerprint}{\index{Settings@\texttt{Settings} (type)}}
\begin{membertable}
\texttt{.fingerprint} & structural & FNV-1a-64 of the non-default entries (0x0 at the defaults) \\
\texttt{.is\_default} & structural & every key at its default \\
\texttt{.keys} & structural & the registered keys, in registry order \\
\texttt{.values} & structural & \{key: value text\} \\
\texttt{.text} & structural & the canonical text \\
\texttt{.json} & structural & the settings JSON (values and sources) \\
\texttt{.get(key)} & structural & the value text of one key \\
\end{membertable}
\octype{ExtRecord}{one cell's hecke\_ext/v5 record (HeckeExtDriver, Stages 1-5): its maximal ideals, groups, u\_ell verdict and settings; lazy}{\index{ExtRecord@\texttt{ExtRecord} (type)}}
\begin{membertable}
\texttt{.n} & structural & the symmetric power n \\
\texttt{.ell} & structural & the coefficient prime \\
\texttt{.schema} & computed & the record schema \\
\texttt{.weight} & computed & k = n + 2 \\
\texttt{.coefficient\_ring} & computed & the record's ring (F\_ell) \\
\texttt{.dims} & computed & dim H\textasciicircum{}q of the bulk, every degree \\
\texttt{.degrees\_run} & computed & the bulk degrees Stages 2-5 ran on \\
\texttt{.operators} & computed & the primes whose operators Stage 1 computed \\
\texttt{.anemic\_pool} & computed & the anemic pool \{q \textless{}= 31\} \textbackslash{} \{ell\} \\
\texttt{.s\_sep} & computed & S\_sep(ell) as Stage 2(3) ran it (empty if not computed) \\
\texttt{.dim\_M\_k} & computed & dim M\_k \\
\texttt{.dim\_S\_k} & computed & dim S\_k \\
\texttt{.dim\_dickson} & computed & the Dickson count \\
\texttt{.H1\_shriek} & computed & dim H\textasciicircum{}1\_! (cone = true), else nil \\
\texttt{.les\_exact} & computed & the cone's LES exact (cone = true), else nil \\
\texttt{.groups} & computed & the record's groups, \{name, available, dim\} \\
\texttt{.u\_ell} & computed & the cell's U\_ell verdict and its resumable state \\
\texttt{.ideals} & computed & the maximal ideals, as MaximalIdeal handles \\
\texttt{.num\_ideals} & computed & the number of maximal ideals \\
\texttt{.has\_sign\_records} & computed & whether sign\_split = true built H\textasciicircum{}+ / H\textasciicircum{}- \\
\texttt{.ideals\_plus} & computed & the ideals of the H\textasciicircum{}+ record (sign\_split = true) \\
\texttt{.ideals\_minus} & computed & the ideals of the H\textasciicircum{}- record (sign\_split = true) \\
\texttt{.settings} & computed & the MathSettings the record was computed under \\
\texttt{.verify} & computed & HeckeExtRecord::verify() (0 = consistent) \\
\texttt{.json} & computed & the hecke\_ext/v5 record text (carries the build's commit) \\
\texttt{.digest} & computed & FNV-1a-64 of the record text with \texttt{commit} redacted (NOT the golden ext\_json\_* hash) \\
\texttt{.ideal(i)} & computed & the i-th maximal ideal \\
\end{membertable}
\octype{MaximalIdeal}{one maximal ideal m of a cell's Hecke algebra (an ideal of the ExtRecord)}{\index{MaximalIdeal@\texttt{MaximalIdeal} (type)}}
\begin{membertable}
\texttt{.id} & computed & the ideal's id in its record \\
\texttt{.sign} & computed & all / plus / minus \\
\texttt{.dim} & computed & dim\_F\_ell of the block H\_m \\
\texttt{.group} & computed & the ambient group of the block \\
\texttt{.local} & computed & certified local \\
\texttt{.residue\_degree} & computed & f = [T/m : F\_ell] \\
\texttt{.residue\_degree\_certified} & computed & whether f is certified \\
\texttt{.witness\_poly} & computed & pi\_w, the witness's minimal polynomial \\
\texttt{.w\_hat\_exponent} & computed & N of the Teichmueller lift (local only) \\
\texttt{.final} & computed & true / false / heuristic \\
\texttt{.final\_reason} & computed & why not final ('' when final) \\
\texttt{.certificate} & computed & cyclic / commutant / sturm(B') / dickson / boundary / unavailable \\
\texttt{.certificate\_routes} & computed & the routes' outcomes \{a, b, c, d0, d1\} (d0 / d1: the finality certificates, hecke\_ext/v4) \\
\texttt{.generators\_certified} & computed & the certified generating primes \\
\texttt{.boundary\_reps} & computed & R(ell), the class representatives (d1) tested (empty unless it reached them) \\
\texttt{.operator\_primes} & computed & the primes of the restricted operators \\
\texttt{.tangent\_active} & computed & the primes whose n\_p is active in m/m\textasciicircum{}2 \\
\texttt{.embdim} & computed & dim\_K m/m\textasciicircum{}2 (local only) \\
\texttt{.nilpotency} & computed & the nilpotency index nu (local only) \\
\texttt{.gorenstein} & computed & dim soc = f (local only) \\
\texttt{.label} & computed & the atlas label \{type, system, theta\_twist, ...\} or nil \\
\texttt{.predicted} & computed & the predictions attached to the ideal \\
\texttt{.unpredicted} & computed & what was found and not predicted \\
\texttt{.modules} & computed & the LoewyProfile of every group the ideal lives in \\
\texttt{.algebra} & computed & the local Hecke algebra T\_m (HeckeAlgebra) \\
\texttt{.t(p)} & computed & the restricted operator t\_p = T\_p\textbar{}\_\{H\_m\} \\
\texttt{.character(p)} & computed & P\_p, the eigenvalue of T\_p as a polynomial in the witness (local only) \\
\end{membertable}
\octype{HeckeAlgebra}{the local Hecke algebra T\_m = A\_S of a certified-local ideal (HeckeLocalAlgebra; [X] Sec. 7's \texttt{LocalAlgebra})}{\index{HeckeAlgebra@\texttt{HeckeAlgebra} (type)}}
\begin{membertable}
\texttt{.ell} & computed & the prime \\
\texttt{.D} & computed & dim\_F\_ell of the block \\
\texttt{.f} & computed & the residue degree \\
\texttt{.dim} & computed & dim\_F\_ell A\_S \\
\texttt{.dim\_m} & computed & dim\_F\_ell m\_S \\
\texttt{.nilpotency} & computed & nu: m\textasciicircum{}nu = 0 != m\textasciicircum{}\{nu-1\} \\
\texttt{.hilbert} & computed & dim(m\textasciicircum{}i/m\textasciicircum{}\{i+1\})/f \\
\texttt{.socle\_dim} & computed & dim\_F\_ell soc A\_S \\
\texttt{.gorenstein} & computed & dim soc A\_S == f \\
\texttt{.embdim} & computed & dim\_K m/m\textasciicircum{}2 \\
\texttt{.primes} & computed & S, ascending \\
\texttt{.active\_set} & computed & the primes whose n\_p-bar is nonzero \\
\texttt{.minimal\_active} & computed & a K-spanning subset, greedy in ascending p \\
\texttt{.min\_generators} & computed & dim(V/m V)/f \\
\texttt{.cyclic} & computed & certificate (a): V cyclic over A\_S \\
\texttt{.verify} & computed & certificate (a)'s verify() (0 = re-derives) \\
\texttt{.basis} & computed & an F\_ell-basis of A\_S (canonical RREF) \\
\texttt{.radical} & computed & an F\_ell-basis of m\_S \\
\texttt{.w\_hat} & computed & the Teichmueller lift of the witness \\
\texttt{.invariants} & computed & the SAME invariants by the certificate-free route of algebra\_invariants(t\_p) -- a cross-check \\
\texttt{.json} & computed & the LocalAlgebra JSON \\
\texttt{.t(p)} & computed & t\_p = T\_p restricted to the block \\
\texttt{.n(p)} & computed & n\_p = t\_p - P\_p(w\_hat), the nilpotent generator \\
\end{membertable}
\octype{LoewyProfile}{the Stage-5 module structure of a block at m (HeckeModuleStructure; [X] Sec. 7's \texttt{ModuleStructure}): Loewy / socle series, Jordan types, structure maps}{\index{LoewyProfile@\texttt{LoewyProfile} (type)}}
\begin{membertable}
\texttt{.group} & structural & the ambient group \\
\texttt{.D} & computed & dim\_F\_ell V \\
\texttt{.f} & computed & the residue degree \\
\texttt{.loewy} & computed & Loewy layer dimensions in units of f \\
\texttt{.socle} & computed & socle layer dimensions in units of f \\
\texttt{.m\_dims} & computed & dim\_F\_ell m\textasciicircum{}j V, j = 0 .. nu\_V \\
\texttt{.socle\_dims} & computed & dim\_F\_ell soc\textasciicircum{}j V \\
\texttt{.primes} & computed & the generators' primes \\
\texttt{.cyclic} & computed & V cyclic over T\_m \\
\texttt{.free} & computed & V free of rank one \\
\texttt{.min\_generators} & computed & dim(V/mV)/f \\
\texttt{.semisimple} & computed & mV = 0 (no extension at m) \\
\texttt{.verify} & computed & ModuleStructure::verify() (0 = re-derives) \\
\texttt{.json} & computed & the module structure JSON \\
\texttt{.structure\_maps} & computed & the ExtensionClass of every consecutive layer pair \\
\texttt{.jordan(p)} & computed & the Jordan type of n\_p on V, in units of f \\
\texttt{.layer(j)} & computed & the canonical basis of the Loewy layer m\textasciicircum{}j V / m\textasciicircum{}\{j+1\} V (columns) \\
\texttt{.structure\_map(j)} & computed & the ExtensionClass of the layer pair (j, j+1) \\
\end{membertable}
\octype{ExtensionClass}{the structure map theta\_j between two consecutive Loewy layers ([E] D4): the extension class the generators n\_p glue the layers by}{\index{ExtensionClass@\texttt{ExtensionClass} (type)}}
\begin{membertable}
\texttt{.top\_layer} & computed & j \\
\texttt{.bottom\_layer} & computed & j + 1 \\
\texttt{.top\_dim} & computed & the top layer's K-dimension \\
\texttt{.bottom\_dim} & computed & the bottom layer's K-dimension \\
\texttt{.rank} & computed & F\_ell rank of the stacked map \\
\texttt{.rank\_over\_f} & computed & its K-rank \\
\texttt{.active} & computed & the primes whose own map is nonzero \\
\texttt{.rank\_p} & computed & the F\_ell rank per generator \\
\texttt{.split} & computed & rank == 0: the two layers are not glued \\
\texttt{.matrix(p)} & computed & the matrix of n\_p between the layers (rows bottom, columns top) \\
\end{membertable}
\octype{Presentation}{generators and a minimal relation set of a certified-local block over T\_m (HeckePresentation, [X] Sec. 2.2 Stage 5, OQ11)}{\index{Presentation@\texttt{Presentation} (type)}}
\begin{membertable}
\texttt{.available} & structural & whether the block was local \\
\texttt{.min\_generators} & structural & t, over A\_S \\
\texttt{.min\_relations} & structural & over A\_S \\
\texttt{.dim\_relation\_module} & structural & t dim A\_S - D \\
\texttt{.free} & structural & no relations \\
\texttt{.generator\_indices} & structural & the generators' standard indices \\
\texttt{.generators} & structural & the generators as columns \\
\texttt{.relations} & structural & the relations as t-tuples over the A\_S basis \\
\texttt{.verify} & structural & Presentation::verify() (0 = re-derives) \\
\texttt{.json} & structural & the presentation JSON \\
\end{membertable}
\octype{Involution}{the production involution iota\textasciicircum{}* = -iota\textasciicircum{}*\_geo on H ([C] def:iota(iv); HeckeInvolution) with the +/- split and, per degree, the sign split of the Hecke action and the Eisenstein line (HeckeSignSplit); lazy}{\index{Involution@\texttt{Involution} (type)}}
\begin{membertable}
\texttt{.n} & computed & the symmetric power \\
\texttt{.ell} & computed & the prime \\
\texttt{.status} & computed & the split's status \\
\texttt{.coefficient\_ring} & computed & F\_ell, or Z/ell\textasciicircum{}k (slot 20: iota honours the ring) \\
\texttt{.lengths\_plus} & computed & length H\textasciicircum{}q\_+ per degree (= dims\_plus over F\_ell) \\
\texttt{.lengths\_minus} & computed & length H\textasciicircum{}q\_- per degree (= dims\_minus over F\_ell) \\
\texttt{.dims} & computed & dim H\textasciicircum{}q \\
\texttt{.dims\_plus} & computed & dim H\textasciicircum{}q\_+ \\
\texttt{.dims\_minus} & computed & dim H\textasciicircum{}q\_- \\
\texttt{.certificate} & computed & InvolutionCertificate::verify() (0 = obl:iota discharged; over Z/ell\textasciicircum{}k: iota\textasciicircum{}2 = id and H = H\textasciicircum{}+ (+) H\textasciicircum{}- were ENGINE\_CHECKed) \\
\texttt{.json} & computed & hecke\_involution/v1 (over Z/ell\textasciicircum{}k: hecke\_involution\_ring/v1) \\
\texttt{.invariant\_factors\_plus(q)} & computed & H\textasciicircum{}q\_+ = (+)\_i Z/ell\textasciicircum{}\{e\_i\}: the e\_i, descending (F\_ell: all 1) \\
\texttt{.invariant\_factors\_minus(q)} & computed & H\textasciicircum{}q\_- = (+)\_i Z/ell\textasciicircum{}\{e\_i\}: the e\_i, descending (F\_ell: all 1) \\
\texttt{.matrix(q)} & computed & iota\textasciicircum{}* on H\textasciicircum{}q \\
\texttt{.plus(q)} & computed & a basis of H\textasciicircum{}q\_+ (columns, in H\textasciicircum{}q coordinates) \\
\texttt{.minus(q)} & computed & a basis of H\textasciicircum{}q\_- (columns, in H\textasciicircum{}q coordinates) \\
\texttt{.sign\_split(q)} & computed & the Hecke action's sign split at \{2,3,5,7\} \textbackslash{} \{ell\} and the Eisenstein line \\
\texttt{.eisenstein\_line(q)} & computed & the Eisenstein line in the stored H\textasciicircum{}q basis (defined only) \\
\end{membertable}
\octype{BoundaryCohomology}{H\textasciicircum{}*\_bd on the boundary circle (SL2ZBoundaryPair), with the restriction r and T\_p\textasciicircum{}bd (HeckeBoundaryOperator); lazy}{\index{BoundaryCohomology@\texttt{BoundaryCohomology} (type)}}
\begin{membertable}
\texttt{.n} & computed & the symmetric power \\
\texttt{.ell} & computed & the prime \\
\texttt{.model} & computed & the boundary model \\
\texttt{.is\_group\_cohomology} & computed & 2 in R\textasciicircum{}x (the circle computes group cohomology) \\
\texttt{.dims} & computed & dim H\textasciicircum{}q\_bd \\
\texttt{.bulk\_dims} & computed & dim H\textasciicircum{}q of the bulk \\
\texttt{.q\_max} & computed & the top degree of THIS complex (derived) \\
\texttt{.truncated} & computed & false on a finite model \\
\texttt{.q\_max\_source} & computed & 'derived' unless truncated (E9) \\
\texttt{.degrees} & computed & [0 .. q\_max], derived \\
\texttt{.coefficient\_ring} & computed & the complex's ring (the extension layer is F\_ell) \\
\texttt{.r(q)} & computed & the restriction H\textasciicircum{}q -\textgreater{} H\textasciicircum{}q\_bd \\
\texttt{.A(p,\allowbreak{} q)} & computed & T\_p\textasciicircum{}bd on H\textasciicircum{}q\_bd \\
\end{membertable}
\octype{ConeCohomology}{H\textasciicircum{}*\_c = H\textasciicircum{}*(Cone(r)) with its long exact sequence (EquivariantMappingCone) and, per prime, the cone action A\_c(T\_p) with h DERIVED from the certified pair homotopy eta over Z (HeckeConeAction); lazy}{\index{ConeCohomology@\texttt{ConeCohomology} (type)}}
\begin{membertable}
\texttt{.n} & computed & the symmetric power \\
\texttt{.ell} & computed & the prime \\
\texttt{.name} & computed & the cone's name \\
\texttt{.dims} & computed & dim H\textasciicircum{}q\_c \\
\texttt{.bulk\_dims} & computed & dim H\textasciicircum{}q of the bulk \\
\texttt{.boundary\_dims} & computed & dim H\textasciicircum{}q\_bd \\
\texttt{.q\_max} & computed & the top degree of THIS complex (derived) \\
\texttt{.truncated} & computed & false on a finite model \\
\texttt{.q\_max\_source} & computed & 'derived' unless truncated (E9) \\
\texttt{.degrees} & computed & [0 .. q\_max], derived \\
\texttt{.coefficient\_ring} & computed & the complex's ring (the extension layer is F\_ell) \\
\texttt{.certificate} & computed & ConeCertificate::verify() (0 = re-derives) \\
\texttt{.H1\_shriek} & computed & dim H\textasciicircum{}1\_! = dim ker r\textasciicircum{}1 (nil if not derivable) \\
\texttt{.les} & computed & the long exact sequence spots \{term, q, dim, rank\_in, rank\_out, asserted, exact\} \\
\texttt{.les\_exact} & computed & every asserted spot of the LES is exact \\
\texttt{.pi(q)} & computed & H\textasciicircum{}q\_c -\textgreater{} H\textasciicircum{}q \\
\texttt{.r(q)} & computed & H\textasciicircum{}q -\textgreater{} H\textasciicircum{}q\_bd \\
\texttt{.iota(q)} & computed & H\textasciicircum{}q\_bd -\textgreater{} H\textasciicircum{}\{q+1\}\_c (the connecting map) \\
\texttt{.A(p,\allowbreak{} q)} & computed & A\_c(T\_p) on H\textasciicircum{}q\_c \\
\texttt{.h(p,\allowbreak{} q)} & computed & the cochain homotopy h at p in degree q, DERIVED from eta (not stored) \\
\texttt{.pair\_homotopy(p)} & computed & the certified pair homotopy eta at p: \{eta\_zero, slots, certify, solve\_regime, ...\} \\
\texttt{.action\_certificate(p)} & computed & ConeActionCertificate::verify() of the action at p (0 = re-derives) \\
\end{membertable}
\octype{DicksonLES}{the long exact sequence of 0 -\textgreater{} Sym\textasciicircum{}n (x) det -\textgreater{} Sym\textasciicircum{}\{n+ell+1\} -\textgreater{} Q\_n -\textgreater{} 0 on the model over F\_ell, ell \textgreater{}= 5 ([E] sec:thetaobj; HeckeDicksonTheta): the six terms, the five maps in the stored bases (theta = L\_* source-indexed, delta by the snake construction), exactness at all six spots as subspaces, H\textasciicircum{}0 the Dickson ring; lazy}{\index{DicksonLES@\texttt{DicksonLES} (type)}}
\begin{membertable}
\texttt{.n} & structural & the symmetric power of the source term \\
\texttt{.ell} & structural & the prime \\
\texttt{.n\_up} & structural & n + ell + 1, the middle term's symmetric power \\
\texttt{.dims} & computed & \{H0\_n, H0\_up, H0\_Q, H1\_n, H1\_up, H1\_Q\} \\
\texttt{.ranks} & computed & \{theta0, pi0, delta, theta1, pi1\} \\
\texttt{.maps} & computed & \{theta0, pi0, delta, theta1, pi1\} as Matrix values in the stored bases (column convention) \\
\texttt{.exact} & computed & exactness certified at all six spots as subspaces \\
\texttt{.ker\_theta} & computed & a basis of ker theta\textasciicircum{}1 (columns, in H1\_n's stored coordinates) \\
\texttt{.im\_delta} & computed & a basis of im delta (columns, in H1\_n's stored coordinates) \\
\texttt{.ker\_equals\_im} & computed & ker theta\textasciicircum{}1 = im delta as subspaces of H1\_n \\
\texttt{.dickson\_count} & structural & \#\{(a, b) : a(ell + 1) + b(ell\textasciicircum{}2 - ell) = n\} \\
\texttt{.certificate} & computed & DicksonThetaCertificate code (0 = every check of [E] (6a)-(6d) held) \\
\texttt{.coefficient\_ring} & structural & the complex's ring (theta is F\_ell only) \\
\end{membertable}
\octype{DescentRecord}{the integral descent record integral\_descent\_record/v1 of H\textasciicircum{}q(Gamma\_0(N); Sym\textasciicircum{}n Z\textasciicircum{}2) over Z ([C] thm:descentcert; slot-24 D6): the groups to q\_max = 3 (object = group, truncated, periodic from q = 2), the bad primes, the certified precision K per prime, char\_Z(T\_p \textbar{} H\textasciicircum{}1/tors), the torsion Hecke modules by their invariants, the boundary pair (sec:descentpair); only under coefficient\_ring = Z; lazy}{\index{DescentRecord@\texttt{DescentRecord} (type)}}
\begin{membertable}
\texttt{.N} & structural & the level \\
\texttt{.n} & structural & the symmetric power \\
\texttt{.primes} & structural & the Hecke primes P (none dividing N) \\
\texttt{.coefficient\_ring} & structural & Z (the record exists only under coefficient\_ring = Z) \\
\texttt{.object} & structural & group (group cohomology, never Bredon; OD-10 (a)) \\
\texttt{.q\_max} & structural & 3, derived from thm:descentcert(6) \\
\texttt{.truncated} & structural & true: degrees above q\_max are not reported (they are periodic, thm:descentcert(6)) \\
\texttt{.q\_max\_source} & structural & derived \\
\texttt{.degrees} & structural & [0, 1, 2, 3] \\
\texttt{.groups} & computed & H\textasciicircum{}0 .. H\textasciicircum{}3 as abelian groups (\textquotedbl{}Z\textasciicircum{}3 + Z/2 + Z/4\textquotedbl{}); empty when refused \\
\texttt{.bad\_primes} & computed & B(M\_N) of thm:descentcert(1): List of [ell, count] (prop:badprimesN(1)), ell \textgreater{}= 5 \\
\texttt{.torsion\_primes} & computed & the primes whose torsion is read: 2, 3 (n even) and B(M\_N) \\
\texttt{.K} & computed & per prime of the groups' run: \{ell, route, K, W, certified\} (certification provenance, not hashed) \\
\texttt{.charpoly} & computed & per p: \{p, r, coefficients (high -\textgreater{} low, String: an integer, \textquotedbl{}[sign,D digits]\textquotedbl{} beyond 2\textasciicircum{}63, \textquotedbl{}?\textquotedbl{} undetermined), exact (high -\textgreater{} low, String: every coefficient as an exact integer, \textquotedbl{}?\textquotedbl{} undetermined), determined, digest, refused\} -- char\_Z(T\_p \textbar{} H\textasciicircum{}1/tors), prop:crtcharpolyN's output rule \\
\texttt{.charpoly\_factored} & computed & per p: \{p, status (ok \textbar{} unavailable), expanded, factored (\textquotedbl{}(X + 24)\textasciicircum{}2 (X - 2049)\textquotedbl{}), factors: [\{coefficients (high -\textgreater{} low, String), degree, multiplicity, irreducible\}], complete, verified, note\} -- char\_Z(T\_p \textbar{} H\textasciicircum{}1/tors) factored over Z (Yun + Zassenhaus, exact; the product of the factors is checked against char\_Z); presentation only, not in the record's JSON or digest \\
\texttt{.torsion} & computed & per (q, ell): \{q, ell, exponents, module\_certified, ops: [\{p, invariants\}], refused\} -- the torsion Hecke module by its choice-free invariants (prop:torsionhecke(1)) \\
\texttt{.pair} & computed & the boundary pair (sec:descentpair): \{bulk, boundary, compact\} each \{q\_max, truncated, q\_max\_source, degrees, groups\}, q = 0 .. 4 (compact, the cone, also derived\_from = [bulk, boundary]); refused \\
\texttt{.groups\_inv} & computed & H\textasciicircum{}0 .. H\textasciicircum{}3 as structures: per q \{free\_rank, torsion: [[p, k, mult], ...]\} = Z\textasciicircum{}free\_rank + (+) (Z/p\textasciicircum{}k)\textasciicircum{}mult (the same groups as \texttt{groups}); empty when refused \\
\texttt{.pair\_inv} & computed & the boundary pair's groups as structures: \{bulk, boundary, compact\}, each a List over q = 0 .. 4 of \{free\_rank, torsion: [[p, k, mult], ...]\} (the same groups as \texttt{pair}) \\
\texttt{.status} & computed & ok \textbar{} refused \\
\texttt{.refusal} & computed & the first refused part's reason (\textquotedbl{}\textquotedbl{} when ok) \\
\texttt{.digest} & computed & record\_fnv1a64: FNV-1a-64 of the record's invariants and the settings' canonical text \\
\texttt{.json} & computed & the record as one JSON line (schema integral\_descent\_record/v1) \\
\texttt{.module(ell,\allowbreak{} K)} & structural & H\textasciicircum{}1(Gamma\_0(N); Sym\textasciicircum{}n Z\textasciicircum{}2) (x) Z/ell\textasciicircum{}K with T\_p (p in primes) by cochain-level lifts (slot-25 EXPORT; [C] prop:torsionhecke, prop:strictliftN): \{N, n, ell, K, K\_max, route, carrier, status (ok \textbar{} blocked(export) \textbar{} refused), refusal, free\_rank, torsion (exponents at K), invariant\_factors, generators: [\{kind, exponent, cocycle (a cocycle mod ell\textasciicircum{}K)\}], ops: [\{p, matrix (column convention, row i mod ell\textasciicircum{}exponent\_i), charpoly\_free (high -\textgreater{} low, mod ell\textasciicircum{}K)\}]\}; one run at the route's width serves every K \textless{}= K\_max (computed on the first call per ell; it does not force the record); ell in \{2, 3\}: the tree model, ell \textgreater{}= 5: the cokernel presentation (K\_max \textless{}= 1 for ell \textgreater{}= 46341) \\
\texttt{.module\_Z(mode)} & computed & H\textasciicircum{}1(Gamma\_0(N); Sym\textasciicircum{}n Z\textasciicircum{}2) over Z as a presented Z[T\_p]-module ([C] prop:torsionhecke): T\_p = [[tau, M], [0, F]] in generators torsion first, then free; module\_Z() = EXPORT-Z step 2 (F exact over Z), module\_Z(\textquotedbl{}local\textquotedbl{}) = step 1 (slot-25 EXPORT-Z, F pending). \{N, n, primes, status (ok \textbar{} blocked(export)), refusal, splitting (export\_retraction/v1), retractions: [\{ell, route, K\_run, K, fnv, [step 2: pi\_B (t\_ell x r, row a mod d\_a: the retraction on B, the torsion coordinates of the classes [b\_j]), pi\_B\_fnv]\}], torsion\_primes, largest\_torsion\_prime, invariant\_factors (the elementary divisors d\_i = ell\textasciicircum{}sigma\_i, ell ascending; key name kept), torsion\_generators: [\{ell, exponent, d\}], free\_rank (a blocked value prints these four as \textquotedbl{}blocked\textquotedbl{}), free\_block (step 2: \textquotedbl{}exact\textquotedbl{}, or \textquotedbl{}refused\textquotedbl{} with free\_refusal; local: \textquotedbl{}pending(step 2)\textquotedbl{}), M\_basis (step 2: global, M in the basis B; local: M's ell-block is against module(ell, K)'s free generators), entry\_encoding (step 2: exact, F as decimal integer strings; local: word), split (unknown), presentation, [step 2: basis: \{rule (iso\_dixon/v3), q0, saturation (det\_pm1), basis\_fnv, phi\_fnv, bits\_B, bits\_Phi, bound (phi0 \textbar{} phi\_cert by the Auto size rule; - at rank 0)\}, free\_refusal], ops: [\{p, tau (row i mod d\_i, block-diagonal over ell), M: [\{ell, K, matrix\}] (step 2: against the classes [b\_j] of B, T\_p b\_j = sum\_i F\_ij b\_i + sum\_a M\_aj g\_a, i.e. M = M\_local C - pi\_B F + tau pi\_B; local: against module(ell, K)'s free generators), F (step 2: r x r decimal strings; local: \textquotedbl{}pending(step 2)\textquotedbl{}), F\_local: [\{ell, modulus, matrix\}] (the export's free block mod ell\textasciicircum{}K, not F), charpoly\_Z (the record's char\_Z(T\_p \textbar{} H\textasciicircum{}1/tors), high -\textgreater{} low), [step 2: F\_fnv, F\_proof (apriori \textbar{} apost \textbar{} rank0), F\_primes, M\_fnv (FNV-1a-64 of M over every ell, span convention), C: [\{ell, modulus, matrix\}] (B in the export's local free generators: C F = F\_local C mod ell\textasciicircum{}K)]\}], torsion\_canon: [\{ell, ops: [\{p, invariants\}]\}] (the record's)\}; unique up to [[Aut T, Hom(L, T)], [0, GL\_r(Z)]] \\
\end{membertable}


\chapter{Settings reference}\label{app:settings}

\texttt{math.}\ marks a mathematical setting (may change a value, moves the fingerprint); \texttt{engine}
settings never change a value, and their default is named by the environment variable that supplies it.
\emph{env-only} keys are shown by \vb{settings()} but cannot be \texttt{set}. The eleven mathematical rows
share one registry boilerplate sentence; their values and meanings are in Section~\ref{sec:mathsettings}.

% GENERATED by docs/latex/tools/gen_cli_manual_tables.sh, i.e.
% `CLI_MANUAL_EMIT=<dir> hacs CliManualCoverageValidation`, from the
% registries of the build it runs in.  Do not edit by hand: regenerate when a
% verb, type, setting, validation or command-line token changes.
\begin{settingtable}
\ocsetting{print\_max\_entries}\index{print_max_entries@\texttt{print\_max\_entries} (setting)} & engine & Int & session & entries printed before a value prints its summary form \\
\ocsetting{verbose}\index{verbose@\texttt{verbose} (setting)} & engine & Bool & session & let the pipeline's own stdout/stderr chatter through (default: true in the REPL, false in scripts) \\
\ocsetting{locality}\index{locality@\texttt{locality} (setting)} & \textbf{math.} & String & \texttt{\textquotedbl{}budget=64\textquotedbl{}} & mathematical setting \texttt{locality} (hecke::MathSettings -- it changes what a result MEANS, so a non-default value moves the fingerprint and \texttt{golden(...)} refuses) \\
\ocsetting{u\_ell}\index{u_ell@\texttt{u\_ell} (setting)} & \textbf{math.} & String & \texttt{\textquotedbl{}auto\textquotedbl{}} & mathematical setting \texttt{u\_ell} (hecke::MathSettings -- it changes what a result MEANS, so a non-default value moves the fingerprint and \texttt{golden(...)} refuses) \\
\ocsetting{sturm}\index{sturm@\texttt{sturm} (setting)} & \textbf{math.} & String & \texttt{\textquotedbl{}auto\textquotedbl{}} & mathematical setting \texttt{sturm} (hecke::MathSettings -- it changes what a result MEANS, so a non-default value moves the fingerprint and \texttt{golden(...)} refuses) \\
\ocsetting{sturm\_unavailable}\index{sturm_unavailable@\texttt{sturm\_unavailable} (setting)} & \textbf{math.} & String & \texttt{\textquotedbl{}nonfinal\textquotedbl{}} & mathematical setting \texttt{sturm\_unavailable} (hecke::MathSettings -- it changes what a result MEANS, so a non-default value moves the fingerprint and \texttt{golden(...)} refuses) \\
\ocsetting{commutant\_budget}\index{commutant_budget@\texttt{commutant\_budget} (setting)} & \textbf{math.} & String & \texttt{\textquotedbl{}20000000000\textquotedbl{}} & mathematical setting \texttt{commutant\_budget} (hecke::MathSettings -- it changes what a result MEANS, so a non-default value moves the fingerprint and \texttt{golden(...)} refuses) \\
\ocsetting{final\_mode}\index{final_mode@\texttt{final\_mode} (setting)} & \textbf{math.} & String & \texttt{\textquotedbl{}closure\textquotedbl{}} & mathematical setting \texttt{final\_mode} (hecke::MathSettings -- it changes what a result MEANS, so a non-default value moves the fingerprint and \texttt{golden(...)} refuses) \\
\ocsetting{field\_polynomial}\index{field_polynomial@\texttt{field\_polynomial} (setting)} & \textbf{math.} & String & \texttt{\textquotedbl{}least\textquotedbl{}} & mathematical setting \texttt{field\_polynomial} (hecke::MathSettings -- it changes what a result MEANS, so a non-default value moves the fingerprint and \texttt{golden(...)} refuses) \\
\ocsetting{split\_field\_report}\index{split_field_report@\texttt{split\_field\_report} (setting)} & \textbf{math.} & String & \texttt{\textquotedbl{}canonical\textquotedbl{}} & mathematical setting \texttt{split\_field\_report} (hecke::MathSettings -- it changes what a result MEANS, so a non-default value moves the fingerprint and \texttt{golden(...)} refuses) \\
\ocsetting{presentation}\index{presentation@\texttt{presentation} (setting)} & \textbf{math.} & String & \texttt{\textquotedbl{}off\textquotedbl{}} & mathematical setting \texttt{presentation} (hecke::MathSettings -- it changes what a result MEANS, so a non-default value moves the fingerprint and \texttt{golden(...)} refuses) \\
\ocsetting{predictions}\index{predictions@\texttt{predictions} (setting)} & \textbf{math.} & String & \texttt{\textquotedbl{}annotate\textquotedbl{}} & mathematical setting \texttt{predictions} (hecke::MathSettings -- it changes what a result MEANS, so a non-default value moves the fingerprint and \texttt{golden(...)} refuses) \\
\ocsetting{coefficient\_ring}\index{coefficient_ring@\texttt{coefficient\_ring} (setting)} & \textbf{math.} & String & \texttt{\textquotedbl{}F\_ell\textquotedbl{}} & mathematical setting \texttt{coefficient\_ring} (hecke::MathSettings -- it changes what a result MEANS, so a non-default value moves the fingerprint and \texttt{golden(...)} refuses) \\
\ocsetting{threads}\index{threads@\texttt{threads} (setting)} & engine & Int & env \texttt{ENGINE\_THREADS} & the engine consumer pool's worker count (engine::tape::consumer\_pool().size()); thread-local, propagated onto run\_concurrently tasks \\
\ocsetting{backend}\index{backend@\texttt{backend} (setting)} & engine & String & env \texttt{ENGINE\_BACKEND} & engine \textbar{} legacy -- the Phase 2/3 rollback; \texttt{legacy} holds engine::elim::ScopedLegacy and forces the legacy cohomology and Hecke-operator arms (thread-local; OC\_ENABLE\_LEGACY=ON builds only -- an OFF build refuses \texttt{legacy}, exit 5) (one of: engine, legacy) \\
\ocsetting{device}\index{device@\texttt{device} (setting)} & engine & String & env \texttt{ENGINE\_DEVICE} & cpu \textbar{} gpu \textbar{} hybrid -- pass 6's device policy (PassOptions); process-global.  \texttt{cpu} is the golden path; the GPU is never the default (one of: cpu, gpu, hybrid) \\
\ocsetting{batch}\index{batch@\texttt{batch} (setting)} & engine & Bool & env \texttt{ENGINE\_BATCH} & pass 5, batching (PassOptions; process-global) \\
\ocsetting{ksplit}\index{ksplit@\texttt{ksplit} (setting)} & engine & Bool & env \texttt{ENGINE\_KSPLIT} & pass 4, k-splitting of GEMMs with k \textgreater{} k\_max (PassOptions; process-global) \\
\ocsetting{arena}\index{arena@\texttt{arena} (setting)} & engine & Bool & env \texttt{ENGINE\_ARENA} & the temp arena (false: every Temp owns its storage) (PassOptions; process-global) \\
\ocsetting{cse}\index{cse@\texttt{cse} (setting)} & engine & Bool & env \texttt{ENGINE\_CSE} & pass 2, common subexpressions (PassOptions; process-global) \\
\ocsetting{fuse}\index{fuse@\texttt{fuse} (setting)} & engine & Bool & env \texttt{ENGINE\_FUSE} & pass 2b, product-into-AXPY fusion (PassOptions; process-global) \\
\ocsetting{layout}\index{layout@\texttt{layout} (setting)} & engine & Bool & env \texttt{ENGINE\_LAYOUT} & pass 3, layout (PassOptions; process-global) \\
\ocsetting{lean}\index{lean@\texttt{lean} (setting)} & engine & Bool & session & the memory-lean configuration of \texttt{-{}-{}bench -{}-{}lean}: drop the cohomology objects no production consumer reads and skip the Theta\_eq solve (A\textasciicircum{}(q) unchanged); thread-local, captured at a handle's CREATION like every setting \\
\ocsetting{program\_cache}\index{program_cache@\texttt{program\_cache} (setting)} & engine & Bool & session & whether compiled tape programs are served from the program cache (false: a fresh compile every time); process-global \\
\ocsetting{confirm}\index{confirm@\texttt{confirm} (setting)} & engine & String & session & a call outside a verb's silent cost envelope (NBCONFIRM): \texttt{ask} = on a terminal REPL show the estimate and ask, in a served session answer confirm\_required with a per-cell token, elsewhere refuse as before (exit 3); \texttt{yes} = run it (hard memory / disk refusals still refuse); \texttt{no} = always refuse as before; not settable in a served session (one of: ask, yes, no) \\
\ocsetting{confirm\_wall\_s}\index{confirm_wall_s@\texttt{confirm\_wall\_s} (setting)} & engine & Int & session & seconds: on a terminal REPL or in a served session, a call outside its silent envelope whose estimated wall time (upper edge) is above this is confirmed before it runs (NBCONFIRM; default 60; serve-settable) \\
\ocsetting{host\_threads}\index{host_threads@\texttt{host\_threads} (setting)} & engine & String & env \texttt{ENGINE\_HOST\_THREADS} (env-only) & environment-only: read once by the engine from ENGINE\_HOST\_THREADS (empty = unset); there is no in-process override \\
\ocsetting{tile\_mult}\index{tile_mult@\texttt{tile\_mult} (setting)} & engine & String & env \texttt{ENGINE\_TILE\_MULT} (env-only) & environment-only: read once by the engine from ENGINE\_TILE\_MULT (empty = unset); there is no in-process override \\
\ocsetting{tile\_cols}\index{tile_cols@\texttt{tile\_cols} (setting)} & engine & String & env \texttt{ENGINE\_TILE\_COLS} (env-only) & environment-only: read once by the engine from ENGINE\_TILE\_COLS (empty = unset); there is no in-process override \\
\ocsetting{spin}\index{spin@\texttt{spin} (setting)} & engine & String & env \texttt{ENGINE\_SPIN} (env-only) & environment-only: read once by the engine from ENGINE\_SPIN (empty = unset); there is no in-process override \\
\ocsetting{elim\_base}\index{elim_base@\texttt{elim\_base} (setting)} & engine & String & env \texttt{ENGINE\_ELIM\_BASE} (env-only) & environment-only: read once by the engine from ENGINE\_ELIM\_BASE (empty = unset); there is no in-process override \\
\ocsetting{cuda\_disable}\index{cuda_disable@\texttt{cuda\_disable} (setting)} & engine & String & env \texttt{ENGINE\_CUDA\_DISABLE} (env-only) & environment-only: read once by the engine from ENGINE\_CUDA\_DISABLE (empty = unset); there is no in-process override \\
\ocsetting{cuda\_device}\index{cuda_device@\texttt{cuda\_device} (setting)} & engine & String & env \texttt{ENGINE\_CUDA\_DEVICE} (env-only) & environment-only: read once by the engine from ENGINE\_CUDA\_DEVICE (empty = unset); there is no in-process override \\
\ocsetting{cuda\_program\_cache}\index{cuda_program_cache@\texttt{cuda\_program\_cache} (setting)} & engine & String & env \texttt{ENGINE\_CUDA\_PROGRAM\_CACHE} (env-only) & environment-only: read once by the engine from ENGINE\_CUDA\_PROGRAM\_CACHE (empty = unset); there is no in-process override \\
\ocsetting{cuda\_graph}\index{cuda_graph@\texttt{cuda\_graph} (setting)} & engine & String & env \texttt{ENGINE\_CUDA\_GRAPH} (env-only) & environment-only: read once by the engine from ENGINE\_CUDA\_GRAPH (empty = unset); there is no in-process override \\
\ocsetting{cuda\_merge\_levels}\index{cuda_merge_levels@\texttt{cuda\_merge\_levels} (setting)} & engine & String & env \texttt{ENGINE\_CUDA\_MERGE\_LEVELS} (env-only) & environment-only: read once by the engine from ENGINE\_CUDA\_MERGE\_LEVELS (empty = unset); there is no in-process override \\
\ocsetting{cuda\_pinned\_mb}\index{cuda_pinned_mb@\texttt{cuda\_pinned\_mb} (setting)} & engine & String & env \texttt{ENGINE\_CUDA\_PINNED\_MB} (env-only) & environment-only: read once by the engine from ENGINE\_CUDA\_PINNED\_MB (empty = unset); there is no in-process override \\
\ocsetting{cuda\_copy\_threads}\index{cuda_copy_threads@\texttt{cuda\_copy\_threads} (setting)} & engine & String & env \texttt{ENGINE\_CUDA\_COPY\_THREADS} (env-only) & environment-only: read once by the engine from ENGINE\_CUDA\_COPY\_THREADS (empty = unset); there is no in-process override \\
\ocsetting{trace\_nodes}\index{trace_nodes@\texttt{trace\_nodes} (setting)} & engine & String & env \texttt{ENGINE\_TRACE\_NODES} (env-only) & environment-only: read once by the engine from ENGINE\_TRACE\_NODES (empty = unset); there is no in-process override \\
\ocsetting{trace\_rss}\index{trace_rss@\texttt{trace\_rss} (setting)} & engine & String & env \texttt{ENGINE\_TRACE\_RSS} (env-only) & environment-only: read once by the engine from ENGINE\_TRACE\_RSS (empty = unset); there is no in-process override \\
\ocsetting{trace\_elim}\index{trace_elim@\texttt{trace\_elim} (setting)} & engine & String & env \texttt{ENGINE\_TRACE\_ELIM} (env-only) & environment-only: read once by the engine from ENGINE\_TRACE\_ELIM (empty = unset); there is no in-process override \\
\ocsetting{trace\_levels}\index{trace_levels@\texttt{trace\_levels} (setting)} & engine & String & env \texttt{ENGINE\_TRACE\_LEVELS} (env-only) & environment-only: read once by the engine from ENGINE\_TRACE\_LEVELS (empty = unset); there is no in-process override \\
\ocsetting{trace\_passes}\index{trace_passes@\texttt{trace\_passes} (setting)} & engine & String & env \texttt{ENGINE\_TRACE\_PASSES} (env-only) & environment-only: read once by the engine from ENGINE\_TRACE\_PASSES (empty = unset); there is no in-process override \\
\end{settingtable}


\chapter{Validation roster}\label{app:validations}

The \CMNumValidations{} registered validations of this build, sorted by name, with their tiers
(\CMNumFast{} \texttt{fast}, \CMNumStandard{} \texttt{standard}, \CMNumSlow{} \texttt{slow}).

% GENERATED by docs/latex/tools/gen_cli_manual_tables.sh, i.e.
% `CLI_MANUAL_EMIT=<dir> hacs CliManualCoverageValidation`, from the
% registries of the build it runs in.  Do not edit by hand: regenerate when a
% verb, type, setting, validation or command-line token changes.
\begin{validationlist}
\ocvalidation{BarCocycleSolverValidation}{fast}
\ocvalidation{BigIntNttValidation}{standard}
\ocvalidation{BinaryF2Validation}{standard}
\ocvalidation{BlockEliminationSlowValidation}{slow}
\ocvalidation{BlockEliminationValidation}{standard}
\ocvalidation{BlockFlatMatrixValidation}{standard}
\ocvalidation{BlockMatValidation}{standard}
\ocvalidation{BlockModelPipelineValidation}{standard}
\ocvalidation{CampaignCertificateValidation}{fast}
\ocvalidation{CampaignRingValidation}{standard}
\ocvalidation{CampaignSmokeValidation}{standard}
\ocvalidation{CliManualCoverageValidation}{fast}
\ocvalidation{CoefficientRingKthRootValidation}{fast}
\ocvalidation{CohomologyBlockModelValidation}{standard}
\ocvalidation{ConsoleC0Validation}{fast}
\ocvalidation{ConsoleCancelValidation}{standard}
\ocvalidation{ConsoleCliPreservationValidation}{fast}
\ocvalidation{ConsoleCostModelValidation}{fast}
\ocvalidation{ConsoleDeterminismValidation}{standard}
\ocvalidation{ConsoleErrorValidation}{fast}
\ocvalidation{ConsoleReadOnlyTagValidation}{fast}
\ocvalidation{ConsoleScriptEmit}{fast}
\ocvalidation{ConsoleScriptValidation}{fast}
\ocvalidation{ConsoleServeAtlasValidation}{standard}
\ocvalidation{ConsoleServeValidation}{standard}
\ocvalidation{ConsoleSettingsValidation}{fast}
\ocvalidation{ConsumerMigrationValidation}{fast}
\ocvalidation{ContractCensusValidation}{fast}
\ocvalidation{CudaBitIdentityValidation}{standard}
\ocvalidation{CudaImmaBitIdentityValidation}{standard}
\ocvalidation{CudaImmaElimValidation}{standard}
\ocvalidation{CudaTransferValidation}{standard}
\ocvalidation{DenseElimPipelineValidation}{standard}
\ocvalidation{DenseElimSlowValidation}{slow}
\ocvalidation{DenseElimValidation}{standard}
\ocvalidation{DescentSmithGpuValidation}{standard}
\ocvalidation{DescentSmithMtValidation}{standard}
\ocvalidation{EngineArenaFirstFitValidation}{fast}
\ocvalidation{EngineArenaValidation}{standard}
\ocvalidation{EngineCoreScalingValidation}{standard}
\ocvalidation{EngineCrossDeviceValidation}{standard}
\ocvalidation{EngineCudaMemoryValidation}{standard}
\ocvalidation{EngineDeterminismCoreValidation}{standard}
\ocvalidation{EngineDeterminismValidation}{standard}
\ocvalidation{EngineDeviceTrapValidation}{slow}
\ocvalidation{EngineElimPassOptionsValidation}{fast}
\ocvalidation{EngineElimSessionValidation}{standard}
\ocvalidation{EngineFingerprintDeviceBitValidation}{fast}
\ocvalidation{EngineLadicInductionValidation}{standard}
\ocvalidation{EngineTapeConsumerValidation}{standard}
\ocvalidation{EngineTapeEliminationValidation}{standard}
\ocvalidation{EngineTapeValidation}{standard}
\ocvalidation{EngineU128CompositeValidation}{standard}
\ocvalidation{FlatMatrixValidation}{fast}
\ocvalidation{GoldenDump}{fast}
\ocvalidation{GoldenEmit}{fast}
\ocvalidation{GoldenRegressionValidation}{standard}
\ocvalidation{GpuLeanSessionValidation}{standard}
\ocvalidation{HeckeAtlasSkeletonValidation}{fast}
\ocvalidation{HeckeAtlasValidation}{slow}
\ocvalidation{HeckeBadPrimeValidation}{fast}
\ocvalidation{HeckeBenchmarkSmoke}{slow}
\ocvalidation{HeckeBlockFormValidation}{standard}
\ocvalidation{HeckeBocksteinCuspValidation}{fast}
\ocvalidation{HeckeBoundaryFinalityValidation}{standard}
\ocvalidation{HeckeBoundaryOperatorValidation}{standard}
\ocvalidation{HeckeBoundaryParabolicValidation}{standard}
\ocvalidation{HeckeBoundaryTauValidation}{standard}
\ocvalidation{HeckeBredonTreeRingValidation}{standard}
\ocvalidation{HeckeBredonTreeValidation}{fast}
\ocvalidation{HeckeBredonU64FlatRouteValidation}{fast}
\ocvalidation{HeckeCanonicalFieldValidation}{fast}
\ocvalidation{HeckeCensusCellListValidation}{fast}
\ocvalidation{HeckeCensusLatticeValidation}{standard}
\ocvalidation{HeckeCensusSplitValidation}{standard}
\ocvalidation{HeckeConeActionValidation}{standard}
\ocvalidation{HeckeConeCrossRouteRingValidation}{standard}
\ocvalidation{HeckeConeValidation}{standard}
\ocvalidation{HeckeCosetTablesSelfCheckValidation}{fast}
\ocvalidation{HeckeCosetTablesValidation}{fast}
\ocvalidation{HeckeCrossCheckValidation}{standard}
\ocvalidation{HeckeDecompositionSmokeValidation}{fast}
\ocvalidation{HeckeDicksonThetaCoverageSlowValidation}{slow}
\ocvalidation{HeckeDicksonThetaHeckeCoverageSlowValidation}{slow}
\ocvalidation{HeckeDicksonThetaHeckeValidation}{standard}
\ocvalidation{HeckeDicksonThetaValidation}{standard}
\ocvalidation{HeckeDualityCoverageSlowValidation}{slow}
\ocvalidation{HeckeDualityValidation}{standard}
\ocvalidation{HeckeEndToEndValidation}{standard}
\ocvalidation{HeckeExactIntSmokeValidation}{fast}
\ocvalidation{HeckeExtAlgebraInvariantsValidation}{standard}
\ocvalidation{HeckeExtAlgebraValidation}{standard}
\ocvalidation{HeckeExtArchitectureValidation}{fast}
\ocvalidation{HeckeExtCampaignValidation}{standard}
\ocvalidation{HeckeExtDeterminismValidation}{standard}
\ocvalidation{HeckeExtEndToEndValidation}{slow}
\ocvalidation{HeckeExtKummerValidation}{slow}
\ocvalidation{HeckeExtRamanujanValidation}{standard}
\ocvalidation{HeckeExtRecordValidation}{fast}
\ocvalidation{HeckeExtSeparationValidation}{standard}
\ocvalidation{HeckeExtSmallEllValidation}{standard}
\ocvalidation{HeckeExtSturmSweepValidation}{standard}
\ocvalidation{HeckeExtendedRegimeValidation}{fast}
\ocvalidation{HeckeFamilyConcurrentBuildValidation}{fast}
\ocvalidation{HeckeFormsFastMulValidation}{standard}
\ocvalidation{HeckeFormsHeckeAlgebraValidation}{standard}
\ocvalidation{HeckeFormsNtt3Validation}{standard}
\ocvalidation{HeckeFpFilterValidation}{standard}
\ocvalidation{HeckeFreeModelValidation}{standard}
\ocvalidation{HeckeHidaCMValidation}{standard}
\ocvalidation{HeckeHidaComponentsValidation}{standard}
\ocvalidation{HeckeHidaControlValidation}{standard}
\ocvalidation{HeckeHidaDiscLambdaValidation}{standard}
\ocvalidation{HeckeHidaEisLambdaValidation}{standard}
\ocvalidation{HeckeHidaHeavyCellProbe}{slow}
\ocvalidation{HeckeHidaLevelEllValidation}{standard}
\ocvalidation{HeckeHidaTowerSlowValidation}{slow}
\ocvalidation{HeckeHidaTowerValidation}{standard}
\ocvalidation{HeckeIdentityValidation}{fast}
\ocvalidation{HeckeIndependenceValidation}{standard}
\ocvalidation{HeckeInducedLocalSystemSmokeValidation}{fast}
\ocvalidation{HeckeIntegralActionValidation}{standard}
\ocvalidation{HeckeIntegralCharpolyValidation}{standard}
\ocvalidation{HeckeInvolutionValidation}{standard}
\ocvalidation{HeckeIotaSignValidation}{standard}
\ocvalidation{HeckeIotaStreamValidation}{standard}
\ocvalidation{HeckeIotaTauSymmetryValidation}{fast}
\ocvalidation{HeckeIsotropyModelEllValidation}{standard}
\ocvalidation{HeckeIsotropyModelTreeSlowValidation}{slow}
\ocvalidation{HeckeIsotropyModelTreeValidation}{standard}
\ocvalidation{HeckeIsotropyModelValidation}{standard}
\ocvalidation{HeckeLValueCoefficientRingValidation}{standard}
\ocvalidation{HeckeLValueEisensteinMeasureValidation}{standard}
\ocvalidation{HeckeLValueIwasawaInvariantsValidation}{standard}
\ocvalidation{HeckeLValueMeasureValidation}{standard}
\ocvalidation{HeckeLValueShapiroValidation}{standard}
\ocvalidation{HeckeLValueValidation}{standard}
\ocvalidation{HeckeLValueWildCharacterValidation}{standard}
\ocvalidation{HeckeLadicCellValidation}{standard}
\ocvalidation{HeckeLadicSplitValidation}{slow}
\ocvalidation{HeckeLambdaModSymValidation}{standard}
\ocvalidation{HeckeLambdaQExpValidation}{standard}
\ocvalidation{HeckeLambdaTowerValidation}{standard}
\ocvalidation{HeckeLambdaTwoVarValidation}{standard}
\ocvalidation{HeckeLambdaValidation}{standard}
\ocvalidation{HeckeLatticeModuleValidation}{slow}
\ocvalidation{HeckeLevelConeValidation}{standard}
\ocvalidation{HeckeLevelInfResValidation}{standard}
\ocvalidation{HeckeLevelNCurveValidation}{standard}
\ocvalidation{HeckeLevelNHeilbronnValidation}{standard}
\ocvalidation{HeckeLevelNSparseValidation}{standard}
\ocvalidation{HeckeLevelNValidation}{standard}
\ocvalidation{HeckeLevelNormalityValidation}{standard}
\ocvalidation{HeckeLocalityProbe}{slow}
\ocvalidation{HeckeLocalitySlowValidation}{slow}
\ocvalidation{HeckeLocalityValidation}{standard}
\ocvalidation{HeckeMazurEisValidation}{standard}
\ocvalidation{HeckeMixedFinalityValidation}{slow}
\ocvalidation{HeckeModuleMtValidation}{slow}
\ocvalidation{HeckeModuleStreamSlowValidation}{slow}
\ocvalidation{HeckeModuleStreamValidation}{standard}
\ocvalidation{HeckeMorphismGeneralizationValidation}{fast}
\ocvalidation{HeckeMorphismRegressionValidation}{standard}
\ocvalidation{HeckeNebentypusValidation}{standard}
\ocvalidation{HeckeNebentypusWildValidation}{standard}
\ocvalidation{HeckeOperatorFamilySmokeValidation}{fast}
\ocvalidation{HeckeOperatorSmokeValidation}{fast}
\ocvalidation{HeckeOracleValidation}{standard}
\ocvalidation{HeckeOrdinaryPrecisionValidation}{standard}
\ocvalidation{HeckePairHomotopyValidation}{standard}
\ocvalidation{HeckePeriodicityClassValidation}{standard}
\ocvalidation{HeckePrecisionReuseSlowValidation}{slow}
\ocvalidation{HeckePrecisionReuseValidation}{standard}
\ocvalidation{HeckePredictionsValidation}{fast}
\ocvalidation{HeckePresentationValidation}{fast}
\ocvalidation{HeckeReductionCoverageSlowValidation}{slow}
\ocvalidation{HeckeReductionValidation}{standard}
\ocvalidation{HeckeRelationsValidation}{standard}
\ocvalidation{HeckeRho2StreamValidation}{standard}
\ocvalidation{HeckeRingLoewyDigitSlowValidation}{slow}
\ocvalidation{HeckeRingLoewyValidation}{standard}
\ocvalidation{HeckeSatCertDumpValidation}{fast}
\ocvalidation{HeckeSatPerfValidation}{slow}
\ocvalidation{HeckeSaturationDefectValidation}{standard}
\ocvalidation{HeckeShapiroValidation}{fast}
\ocvalidation{HeckeSimultaneousSmokeValidation}{fast}
\ocvalidation{HeckeSimultaneousValidation}{standard}
\ocvalidation{HeckeSmallCarriesValidation}{fast}
\ocvalidation{HeckeStrictLiftRelationsValidation}{standard}
\ocvalidation{HeckeStrictLiftValidation}{standard}
\ocvalidation{HeckeSturmBoundValidation}{fast}
\ocvalidation{HeckeTameFamilyValidation}{standard}
\ocvalidation{HeckeTauDumpValidation}{fast}
\ocvalidation{HeckeTauTableSelfCheckValidation}{fast}
\ocvalidation{HeckeTauTableValidation}{fast}
\ocvalidation{HeckeTessellationSmokeValidation}{fast}
\ocvalidation{HeckeTessellationValidation}{fast}
\ocvalidation{HeckeThetaBlockValidation}{standard}
\ocvalidation{HeckeThetaCensusRingRefusalValidation}{fast}
\ocvalidation{HeckeThetaRssDietValidation}{slow}
\ocvalidation{HeckeUEllBudgetValidation}{fast}
\ocvalidation{HeckeUEllLowRankValidation}{slow}
\ocvalidation{HeckeUEllReachSlow}{slow}
\ocvalidation{HeckeUEllReachValidation}{standard}
\ocvalidation{HeckeUEllStreamValidation}{slow}
\ocvalidation{HeckeWideTauSlowValidation}{slow}
\ocvalidation{HeckeWideTauValidation}{standard}
\ocvalidation{HeckeZCertDumpValidation}{slow}
\ocvalidation{HeckeZCrossedHomOracleValidation}{standard}
\ocvalidation{HeckeZEisUpCharValidation}{standard}
\ocvalidation{HeckeZEisUpValidation}{standard}
\ocvalidation{HeckeZExportValidation}{standard}
\ocvalidation{HeckeZExportWordValidation}{standard}
\ocvalidation{HeckeZExportZ2SlowValidation}{slow}
\ocvalidation{HeckeZExportZ2Validation}{standard}
\ocvalidation{HeckeZExportZValidation}{standard}
\ocvalidation{HeckeZFreeBasisPairsValidation}{standard}
\ocvalidation{HeckeZFreeBasisSatValidation}{standard}
\ocvalidation{HeckeZFreeBlockCrtValidation}{standard}
\ocvalidation{HeckeZLocalCensusXcheck}{slow}
\ocvalidation{HeckeZLocalParValidation}{standard}
\ocvalidation{HeckeZLocalValidation}{standard}
\ocvalidation{HeckeZRationalValidation}{fast}
\ocvalidation{HeckeZRecordValidation}{standard}
\ocvalidation{HeckeZThetaSearchValidation}{standard}
\ocvalidation{HeckeZTorValidation}{standard}
\ocvalidation{HeckeZUTagValidation}{fast}
\ocvalidation{HeckeZUpValidation}{standard}
\ocvalidation{HostElimThreadDeterminismValidation}{standard}
\ocvalidation{IntegralDescentGroupsSlowValidation}{slow}
\ocvalidation{IntegralDescentGroupsValidation}{standard}
\ocvalidation{IntegralDescentHeckeSlowValidation}{slow}
\ocvalidation{IntegralDescentHeckeValidation}{standard}
\ocvalidation{IntegralDescentLevelNSlowValidation}{slow}
\ocvalidation{IntegralDescentLevelNValidation}{standard}
\ocvalidation{IntegralDescentMapsSlowValidation}{slow}
\ocvalidation{IntegralDescentMapsValidation}{standard}
\ocvalidation{IntegralDescentPairSlowValidation}{slow}
\ocvalidation{IntegralDescentPairValidation}{standard}
\ocvalidation{IntegralDescentParAsmValidation}{slow}
\ocvalidation{IntegralDescentParLnValidation}{slow}
\ocvalidation{IntegralDescentRecordValidation}{standard}
\ocvalidation{IntegralDescentSmithBlockedValidation}{standard}
\ocvalidation{IntegralDescentTreePairSlowValidation}{slow}
\ocvalidation{IntegralDescentTreePairValidation}{standard}
\ocvalidation{IsotropyModelDeltaSquaredValidation}{fast}
\ocvalidation{KernelAvx2BitIdentityValidation}{fast}
\ocvalidation{KernelBitIdentityValidation}{standard}
\ocvalidation{L23AmalgamReaderValidation}{standard}
\ocvalidation{L23Bin3o1Validation}{slow}
\ocvalidation{LeanCensusFileValidation}{fast}
\ocvalidation{LicenceGuardValidation}{fast}
\ocvalidation{LocalSystemCohomologyValidation}{fast}
\ocvalidation{MatrixRetirementValidation}{fast}
\ocvalidation{MatrixValidation}{fast}
\ocvalidation{ModContextValidation}{standard}
\ocvalidation{ModCtxMontgomeryPairValidation}{fast}
\ocvalidation{NumberFieldCyclicCubicValidation}{standard}
\ocvalidation{NumberFieldLllValidation}{fast}
\ocvalidation{NumberFieldPRankValidation}{standard}
\ocvalidation{NumberFieldRibetFieldValidation}{slow}
\ocvalidation{NumberFieldValidation}{fast}
\ocvalidation{OrbifoldCampaignL23Validation}{standard}
\ocvalidation{OrbifoldCensusL23SlowValidation}{slow}
\ocvalidation{OrbifoldCensusL23Validation}{standard}
\ocvalidation{OrbifoldF3lValidation}{slow}
\ocvalidation{OrbifoldMPAllDegreesValidation}{standard}
\ocvalidation{OrbifoldMPBoundaryValidation}{standard}
\ocvalidation{OrbifoldMPComparisonValidation}{standard}
\ocvalidation{OrbifoldMPHeckeValidation}{standard}
\ocvalidation{OrbifoldMPLevelConeValidation}{standard}
\ocvalidation{OrbifoldMPNerveH3Slow}{slow}
\ocvalidation{OrbifoldMPNerveValidation}{standard}
\ocvalidation{OrbifoldMPRefinementValidation}{standard}
\ocvalidation{OrbifoldMPRulesValidation}{fast}
\ocvalidation{OrbifoldMPTeardropValidation}{fast}
\ocvalidation{OrbifoldMod3kValidation}{slow}
\ocvalidation{OrbifoldRing2kValidation}{slow}
\ocvalidation{PerformanceSmokeValidation}{standard}
\ocvalidation{ProgramCacheReleaseValidation}{fast}
\ocvalidation{Ring2kLadderValidation}{standard}
\ocvalidation{Ring2kR5SimdValidation}{standard}
\ocvalidation{SharedHelpersValidation}{fast}
\ocvalidation{SimplicialCohomologyValidation}{fast}
\ocvalidation{SmallEllPolicyScopeValidation}{fast}
\ocvalidation{SubstrateDifferentialValidation}{standard}
\ocvalidation{SymmetricPowerValidation}{standard}
\ocvalidation{SymmetricPowerValidationSlow}{slow}
\ocvalidation{TernaryF3Validation}{standard}
\ocvalidation{ThetaEqSolveValidation}{standard}
\ocvalidation{VacuousRelationValidation}{standard}
\ocvalidation{ValidationHarnessCatchValidation}{fast}
\ocvalidation{VnniGemmBitIdentityValidation}{standard}
\ocvalidation{mod\_arithmetic\_validation}{fast}
\end{validationlist}


\chapter{Regenerating this manual}\label{app:regen}

The files under \file{docs/latex/cli_manual/} are generated from a built binary:
\begin{shell}
docs/latex/tools/gen_cli_manual_tables.sh <path/to/hacs>           # rewrite the tables
docs/latex/tools/gen_cli_manual_tables.sh <path/to/hacs> --check   # diff only, exit 1 if stale
docs/latex/build_cli_manual.sh <outdir>       # pdflatex + makeindex; outdir in scratch
\end{shell}
The generator runs \texttt{CLI\_MANUAL\_EMIT=<scratch>} \texttt{<binary>} \texttt{CliManualCoverageValidation}, which writes
the tables into a scratch directory, and copies them into the tree. The fast validation
\vb{CliManualCoverageValidation}, part of every \cli{--all}, checks that (1) every committed table equals
what the binary generates; (2) every registered verb, type, setting, validation and command-line token
appears in them; (3) this file inputs each of them; (4) the token table equals the set of strings the
parsers of \file{main.cpp}, \file{Bench/}, \file{Campaign/}, \file{Console/Core/ConsoleCli.cpp},
\file{HeckeExt/HeckeAtlasCli.cpp} and the \cli{--backend} parser compare against; (5) the top-level usage
equals the frozen \file{cli_frozen/help.stdout}; (6) every excerpt line of the manual
occurs in its script and every quoted transcript exists; (7) every registered verb is called in a numbered script that the
manual names (by an excerpt, a transcript or a \verb|\file{}| reference), or is on the allow-list \file{docs/latex/cli_manual_gate.tsv} with a
reason (an allow-list entry for a verb that has an example, or is not registered, is stale and fails); (8) every
\texttt{ENGINE\_*}, \texttt{HECKE\_DESCENT\_*} and \texttt{L23\_*} name read in the sources, and every \texttt{lever} row of that file, appears
as \verb|\env{NAME}| in the environment chapter (Sections~\ref{sec:env}--\ref{sec:env-ln23}); (9) every \verb|\runoutput{<name>}| has its
\file{docs/latex/cli_manual_runs/<name>.txt}. The run outputs are regenerated by
\verb|docs/latex/tools/gen_cli_manual_runs.sh <binary>| (\verb|--check| compares them with timings masked). Prose outside the generated
files is not checked beyond (6)--(9): when a front end changes, revise the matching section too.

\chapter{Versions}\label{app:versions}

\textbf{1.5.15} (2026-11-13, S27 TS; not binding). The campaign key \texttt{hz\_theta\_order} (\texttt{lex} /
\texttt{lex\_then\_greedy}, \texttt{structure = hecke} only): the order of the record's $\vartheta$ search, a row of the generated spec-token
table; the validation roster gains \texttt{HeckeZThetaSearchValidation} (S27 TS-FIX). No other table or count changed.

\textbf{1.5.14} (2026-11-12, S26 LEVERS-26; not binding). Wording only: the rollback arms \texttt{=0} of the eleven
\texttt{HECKE\_DESCENT\_*} levers, \texttt{HECKE\_FIELD\_DIET=0} and \texttt{ENGINE\_VNNI=0}/\texttt{off} are retired (refused in a default
build, read in a legacy build), Sections~\ref{sec:env-cpu}, \ref{sec:env-descent}, \ref{sec:env-module};
\texttt{HECKE\_DESCENT\_SMITH\_PANEL} is a tuning knob; the \texttt{ENGINE\_GPU\_SMITH2K*} switches belong to the opt-in device route;
the default of \texttt{ENGINE\_VNNI2\_FLOOR} is now $2^{20}$ (a speed choice, measured).
No table, count, verb, type, setting or token changes.

\textbf{1.5.13} (2026-11-09, release audit fixes; not binding). No change of mathematical content and no table, count,
verb, type, setting or token changes. The program is called \texttt{hacs} throughout (title, quick start, examples, the
quoted run outputs; \texttt{cohomology\_app} remains a link to it, Section~\ref{sec:quickstart}). The two sections on
running a census with the project's own batch scripts on its own machines (formerly ``Running a census on one machine''
and ``A second machine'') are replaced by one neutral section, Section~\ref{sec:localcensus} (``Long census runs''): the
scripts are not part of the release, and the general rules they implemented are kept. The presets paragraph of
Section~\ref{sec:build} names the \texttt{termux} / \texttt{aarch64-cpu} presets and points to \file{BUILDING.md}.

\textbf{1.5.12} (2026-11-09, slot 26 LAYOUT; not binding). Paths only, no change of content (repository layout, owner
ruling 2026-11-09): the sources moved under \file{src/} (\file{src/hacs/}, \file{src/ModArithmetic/},
\file{src/Engine/}, \file{src/Cuda/}); the script folder is \file{src/hacs/Console/scripts/}; builds are out-of-source
only (Section~\ref{sec:build}, and one troubleshooting item). No table, count, verb, type, setting or token changes.

\textbf{1.5.11} (2026-11-09, slot 26; not binding). No change of mathematical content: machine names and private paths removed for the public release (owner ruling 2026-11-09); the section on the second machine is now written without its host name (label \texttt{sec:second-machine}).

\textbf{1.5.10} (2026-11-08, slot-25 PERF-VF; not binding). The environment row \texttt{ENGINE\_VNNI2\_FLOOR} (the two-limb VNNI route's
work-per-thread floor, Section~\ref{sec:env-cpu}): its default, its parsing (digits only) and when its EFFECTIVE line prints. No other
table, count, verb, type, setting or token changes.

\textbf{1.5.9} (2026-11-07, slot-25 LEVERS-25; not binding). Wording only: the env arms \texttt{=0} of \texttt{ENGINE\_GPU\_LEAN},
\texttt{HECKE\_RHO2\_STREAM}, \texttt{HECKE\_MODULE\_STREAM}, \texttt{HECKE\_MODULE\_EVAL}, \texttt{HECKE\_MODULE\_PARALLEL},
\texttt{HECKE\_ROUTE2\_SHARED\_MEMO}, \texttt{HECKE\_RHO\_SIMD} (also \texttt{scalar}), \texttt{HECKE\_PRECISION\_REUSE},
\texttt{HECKE\_FORMS\_FASTMUL}, \texttt{L23\_TERNARY\_F3}, \texttt{L23\_BIN3O1} (and \texttt{HECKE\_IOTA\_STREAM}, not listed in the tables) are
described as retired: refused in a default build, the earlier route in a legacy build (Sections~\ref{sec:env-device}, \ref{sec:env-module},
\ref{sec:env-l23}). No table, count, verb, type, setting or token changes.

\textbf{1.5.8} (2026-11-07, slot-25 NBCONFIRM S5; not binding). Interactive mode: \texttt{Ctrl-C} cancels the running statement at its next
cancellation point (the old sentence ``a running verb cannot be interrupted'' is withdrawn); the served \texttt{interrupt} likewise. One
validation added to the generated tables (\texttt{ConsoleCancelValidation}); no verb, type, setting or token changes.

\textbf{1.5.7} (2026-11-06, LaTeX polish stage 1; not binding). Uniform author block (the former author line was ``hecke-extensions project'') and the
dedication ``In memoriam Prof.\ Dr.\ G\"unter Harder (1938--2025).'' at the head of the abstract. No table, count, verb, type, setting or token changes.

\textbf{1.5.6} (2026-11-06, slot~25 HZ-REC; not binding). The verb \texttt{hecke\_structure} and the campaign keys
\texttt{structure}, \texttt{algebra}, \texttt{hz\_cert} (paragraph in the chapter on integral cohomology; script \file{060_s25_hecke_structure.oc}).

\textbf{1.5.5} (2026-11-05, slot~24 close; CL-R2 audit finding M6; not binding). Wording only: the rollback arms
\texttt{ENGINE\_ELIM\_SESSION=0}, \texttt{ENGINE\_GPU\_IMMA=0} and \texttt{ENGINE\_GPU\_ELIM=0} are described as retired (refused in a
default build, comparison arms in a legacy build), and the four $U_\ell$ memory levers as retired by CL-R1
(Sections~\ref{sec:env-device}, \ref{sec:env-slot22}, \ref{sec:env-module}). No table, count, verb, type, setting or token changes.

\textbf{1.5.4} (2026-11-03, slot~24 FMTMP; not binding). The descent record's \texttt{.charpoly} entries gain \texttt{exact}, every
coefficient as an exact integer; \texttt{.charpoly\_factored} reads the coefficients beyond $2^{63}$ from the record's exact values
({\ttfamily[C,~sec:recordschema]} item~7~(l)); the generated type table is regenerated.

\textbf{1.5.3} (2026-10-31, slot~24 CL-7; owner ruling~4 of 2026-10-30; not binding). With no argument the
program prints the top-level usage and exits~1; the legacy benchmark is gone (Section~\ref{sec:cli}). The frozen
usage text now lists \cli{--help}, \cli{--atlas}, the console forms and \cli{--serve}; the dispatch lives in
\file{main.cpp}.

\textbf{1.5.2} (2026-10-29, slot~24 NB-R34; not binding). The verb \vb{display} (Section~\ref{sec:printing}, with an
example; [P] 3.14.0-pre.2) and what the notebook page renders (Section~\ref{sec:serve}); the generated tables count
one more verb and overload.

\textbf{1.5.1} (2026-10-29, slot~24 DOCS-FIX; not binding). Referee fixes C1--C12
(\file{docs/implementation/slot24/DOCS\_referee\_1029.md}): $5\equiv1+2^{11}\pmod 7$ in the level-1 example (C1); one
Eisenstein class at level~11, weight~2 (C2); the golden reference pins nine descent cells, 156 tier-A entries (C3); the
invariants string has \texttt{ker=}/\texttt{cok=} only when certified with at most 400 summands, and \texttt{B=} blocks above 300
summands read \texttt{-} (C4); the \texttt{shortcut} route at~3 and the rollback routes listed (C5); \texttt{D digits} are
mixed-radix digits (C6); every piece is certified or refused (C7); the boundary pair is certified by
{\ttfamily[C,~prop:descentpair]} (C8); the orbifold census introduced as the second route at level~1, with its cross-check
(C9); \S\ref{sec:orbifold} opens with the mathematics, the \texttt{orbifold\_cell} v1$\to$v2 history in a footnote (C10);
stabiliser orders $2,4,6$ (C11); \texttt{Z} for rank one (C12). No generated table or run output changed.

\textbf{1.5.0} (2026-10-29, slot~24 DOCS-CLI; not binding). \emph{Restructured by subject} (owner direction 2026-10-29):
getting started; the mathematical objects (new: the model, $\Sym^n$, the coefficient rings, levels and nebentypus);
cohomology over a field (with $\Z/\ell^k$ and the Moerdijk--Pronk model); Hecke operators; \emph{integral cohomology by
descent} (new chapter: what is computed, the algorithm step by step, reading the \vb{DescentRecord} including the
torsion-invariants grammar, worked examples, limits); extensions, Hida theory and the forms algebra (new prose for
\vb{pow}, \vb{kummer\_unit}, \vb{hida\_tower}, \vb{level\_n}, \vb{nebentypus}, \vb{forms\_algebra}, \vb{hida\_lambda},
\vb{lfunction}); the console and its language (new: a tour of the small verbs, the scripts by subject); campaigns (new:
campaigns over $\Z$, a sample orbifold run); performance and engine switches (new: a complete environment chapter, among
them \texttt{ENGINE\_KERNEL}, \texttt{ENGINE\_AVX2\_PAIR}, \texttt{ENGINE\_GPU\_SMITH2K*}, \texttt{ENGINE\_GPU\_LEAN},
\texttt{ENGINE\_TERNARY\_F3}, the \texttt{HECKE\_DESCENT\_*} levers with the panel rule of \texttt{dca97bb8}, the module and
forms levers, the \texttt{L23\_*} levers). New console scripts \file{054_s24_descent.oc} and
\file{055_s24_language_tour.oc} (frozen transcripts); real command-line outputs generated by the new
\file{docs/latex/tools/gen_cli_manual_runs.sh} into \file{docs/latex/cli_manual_runs/}. Corrected: the orbifold record is
\texttt{orbifold\_cell/v2}; the \texttt{ENGINE\_DEVICE} row (the GPU is the fast route of the descent's 2-adic Smith form);
the counts (\CMNumValidations{} validations). \vb{CliManualCoverageValidation} now also requires a worked example per verb
(allow-list \file{docs/latex/cli_manual_gate.tsv}), the environment chapter to name every engine, descent and
\texttt{L23} lever and every listed rollback lever, and every quoted run output to exist. Version history and all
section labels of 1.4.1 are kept.\par\smallskip

\textbf{1.4.1} (2026-10-26, slot~24 D6; wording only; not binding). \texttt{coefficient\_ring = Z} is no longer described as refused: it is stored (fingerprint non-zero) and honoured by the integral descent only (the verb \vb{descent}, the campaign's descent cells, the record \texttt{integral\_descent\_record/v1}; {\ttfamily[C,~sec:recordschema]} item~7, {\ttfamily[C,~def:ringsemantics]}(3)); every other consumer refuses it by name with exit~5; \texttt{Z/<m>} stays reserved. Corrected: the exit-code table (exit~5 row), the \vb{coefficient\_ring} row of the settings table, the troubleshooting item on exit code~5. No generated table moves.\par\smallskip
\textbf{1.4.0} (2026-10-12, slot~22 CLOSE-T). The release of 1.4.0-pre.2: only the version changes (no table, count, verb, type, setting or token; not binding).

\textbf{1.4.0-pre.2} (2026-10-11, slot~22; pre-release). The generated tables checked against the slot-22 census binary (\texttt{gen\_cli\_manual\_tables.sh --check}: current, 22 files; \CMNumValidations{} validations, no new verb, type, setting or token). Prose only. New text: Section~\ref{sec:orbifold} (the orbifold census at $\ell=2,3$: the keys \texttt{orbifold}, \texttt{orbifold\_primes}, \texttt{orbifold\_m\_target}, \texttt{orbifold\_k\_cap}; the \texttt{orbifold\_cell/v1} record (\texttt{v2} since slot~24); reading the \texttt{EFFECTIVE campaign arm=orbifold} line; resume, which stops on a failed record; measured memory and time); Section~\ref{sec:second-machine} (the second machine, a second workstation: the CPU-only build, the anchor gate, the bulk command, PAUSE and resume, the gate); Section~\ref{sec:env-ln23} (the LN23 hooks of \texttt{HeckeLevelNCurveValidation}, the $\Lambda$-tower cells and caps of \texttt{HeckeLambdaTowerValidation}, the saturation-defect hooks); one row in the engine-switch table for the GPU tuning knobs \texttt{ENGINE\_CUDA\_GRAPH}, \texttt{ENGINE\_CUDA\_MERGE\_LEVELS}, \texttt{ENGINE\_CUDA\_PROGRAM\_CACHE}, \texttt{ENGINE\_ELIM\_SESSION} and the trace \texttt{ENGINE\_TRACE\_DEVICE}. The paragraph on the orbifold section in Section~\ref{sec:campaign} is replaced by Section~\ref{sec:orbifold}. Corrected the same day after the audit \file{docs/implementation/slot22/CLI\_pre2\_audit.md} (the second machine's build commit and roster count, the $\Lambda$-tower caps, seven clarifications).

\textbf{1.4.0-pre.1} (2026-10-10, slot~22; pre-release). The prose brought up to the slot-22 tables (regenerated in the tree): \CMNumValidations{} validations (\CMNumFast{} \texttt{fast}, \CMNumStandard{} \texttt{standard}, \CMNumSlow{} \texttt{slow}), \CMNumAll{} run by \cli{--all}; verbs, types and settings unchanged. New text: Section~\ref{sec:env-slot22} (\texttt{ENGINE\_BINARY\_F2}, \texttt{ENGINE\_GPU\_IMMA}, \texttt{ENGINE\_GPU\_ELIM*}, the \texttt{EFFECTIVE DEVICE\_ELIM} line); Section~\ref{sec:mp} (the ten Moerdijk--Pronk validations, \texttt{ORBIFOLD\_MP\_*}); Section~\ref{sec:localcensus} (\texttt{run\_local\_census.sh}, the GPU lane); the \cli{--atlas --query} paragraph; new rows in the engine-switch and read-only-hook tables.

\textbf{1.3.1-pre.1} (2026-10-08, slot~22 DOC22; pre-release). One paragraph added to Section~\ref{sec:serve} (the notebook's page themes, client-side only, [P] Section~18.17). No table, count, verb or token changes.

\textbf{1.3.0} (2026-10-06, the slot-21 close). The tables regenerated from the binary of the release commit
(\texttt{gen\_cli\_manual\_tables.sh --check}: current); \CMNumValidations{} validations (\CMNumFast{} \texttt{fast},
\CMNumStandard{} \texttt{standard}), \CMNumAll{} run by \cli{--all}; \CMNumVerbs{} verbs with \CMNumOverloads{} overloads,
\CMNumReadOnly{} of them served. New since 1.2.0, all served (\texttt{read\_only}): the verbs \vb{pow}, \vb{kummer\_unit} and
\vb{hida\_tower} (slot~21) and \vb{level\_n}, \vb{nebentypus}, \vb{forms\_algebra}, \vb{hida\_lambda} (early slot~22, WP-VERB-2,
landed before the slot-21 close); eleven new validations --- \vb{LicenceGuardValidation}, \vb{NumberFieldLllValidation}
(\texttt{fast}), \vb{HeckeLValueCoefficientRingValidation}, \vb{HeckeLValueEisensteinMeasureValidation},
\vb{HeckeLValueIwasawaInvariantsValidation}, \vb{HeckeLValueWildCharacterValidation}, \vb{HeckeLambdaModSymValidation},
\vb{HeckeLambdaQExpValidation}, \vb{HeckeLambdaTwoVarValidation}, \vb{HeckeLevelNCurveValidation},
\vb{NumberFieldRibetFieldValidation} (\texttt{standard}). Appendix~\ref{app:versions}'s entries for 1.2.0 and earlier now give their
counts literally. No front-end section changes.

\textbf{1.2.0} (2026-10-05, the slot-20 close). The tables regenerated from the close binary
(\texttt{gen\_cli\_manual\_tables.sh --check}: current); $182$ validations, $165$ run by
\cli{--all}; $155$ verbs with $168$ overloads, $159$ of them served. New since 1.1.0:
the verb \vb{lfunction} (slot-20 S1; the record \texttt{lfunction/v1}, the one-variable $\ell$-adic series of
a linear route, \Ecite{sec:lvlfunction}; served); the involution \vb{iota} honours
\vb{coefficient\_ring} \vb{Z/ell\^{}<k>} (slot-20 WP-LV3, \Ecite{sec:ring}; no longer exit~5); two new
validations, \vb{HeckeLValueMeasureValidation} (WP-MEAS, \texttt{standard}; the $691$ arm
\texttt{HECKE\_LVALUE\_MEASURE\_691=1}) and \vb{ConsoleServeAtlasValidation} (NB-4, \texttt{standard});
three token rows made precise (slot-20 WP-8: the spec value \texttt{off} of \texttt{level\_a}, the range
$5\le\ell\le100000$ of \cli{--atlas}, \cli{-} reads the script from stdin); and, from 1.2.0-pre.1, the atlas server
\cli{--serve-atlas} (Section~\ref{sec:serve-atlas}; protocol \texttt{oc-serve-atlas/1}, [P] v3.11
Section~18.16; its tokens \cli{--store}, \cli{--index}, \cli{--expect-fnv}, \cli{--max-reply-bytes},
\cli{--rss-cap-mib}) and the sidecar writer \cli{--atlas --index} / \cli{--out}; one validation,
\vb{ConsoleServeAtlasValidation} (\texttt{standard}).

\textbf{1.1.0} (2026-10-04, the slot-19 close). The tables regenerated from the close binary
(\texttt{gen\_cli\_manual\_tables.sh --check}: current). New since 1.0.0: the served console \cli{--serve}
(Section~\ref{sec:serve}; protocol \texttt{oc-serve/1}, [P] v3.10 Section~18; its tokens \cli{--rss-cap-mib},
\cli{--max-reply-bytes}; built only with \texttt{OC\_BUILD\_NOTEBOOK=ON}); the per-overload \texttt{read\_only}
tag, shown in Appendix~\ref{app:verbs} as \textsf{[served]} / \textsf{[not served]}, with its twin
\file{cli_manual/read_only_verbs.json} and the counts $158$ of $167$ (at 1.1.0); the verb
\vb{report} split into a served pathless overload and the file-writing overload ($167$
overloads at 1.1.0, one more than 1.0.0); the build switches and presets of Section~\ref{sec:build}; six new
validations --- \vb{NumberFieldValidation}, \vb{ConsoleReadOnlyTagValidation} (\texttt{fast}),
\vb{HeckeLValueValidation}, \vb{ConsoleServeValidation}, \vb{HeckeLambdaValidation} (\texttt{standard}),
\vb{HeckeLatticeModuleValidation} (\texttt{slow}) --- so $180$ in all at 1.1.0, $163$ run by \cli{--all}.

\textbf{1.0.0} (2026-10-03, the slot-18 close). The first release; refereed as 1.0.0-pre.1 (Fable F-M, ACCEPT WITH
REPAIRS), repairs at 1.0.0-pre.2.

\clearpage
\addcontentsline{toc}{chapter}{Index}
\printindex

\end{document}
